% Security model and analysis of the composed lattice-ZK video system.
%
% Intended use in main.tex: replace the informal "Security analysis" section
% by % Security model and analysis of the composed lattice-ZK video system.
%
% Intended use in main.tex: replace the informal "Security analysis" section
% by % Security model and analysis of the composed lattice-ZK video system.
%
% Intended use in main.tex: replace the informal "Security analysis" section
% by % Security model and analysis of the composed lattice-ZK video system.
%
% Intended use in main.tex: replace the informal "Security analysis" section
% by \input{sections/security}, and after \appendices add
% \input{sections/security_appendix}. Requires amsmath, amssymb, amsthm and the
% theorem environments definition, lemma, theorem defined in main.tex. Extra
% bibliography entries are in sections/security_refs.bib
% (\bibliography{references,sections/security_refs}).
%
% Every code reference is to commit ea96d60 of feature/lattice-pq-stack.

\providecommand{\Rq}{\mathcal{R}_q}
\providecommand{\norm}[1]{\lVert #1 \rVert}
\providecommand{\Adv}{\mathsf{Adv}}
\providecommand{\pid}{\mathit{pid}}
\providecommand{\Cset}{\mathcal{C}}
\providecommand{\Cdiff}{\bar{\mathcal{C}}}
\providecommand{\Com}{\mathsf{Com}}
\providecommand{\vid}{\mathsf{vid}}
\providecommand{\Hfs}{H_{\mathsf{fs}}}
\providecommand{\Stmt}{\mathsf{Stmt}}
\providecommand{\Canon}{\mathsf{Canon}}
\providecommand{\MSIS}{\textsf{M-SIS}}
\providecommand{\MLWE}{\textsf{M-LWE}}
\providecommand{\epszk}{\varepsilon_{\mathsf{zk}}}
\providecommand{\Amax}{A_{\max}}

\section{Mô hình và phân tích an toàn}
\label{sec:security-model}

Mục này phát biểu hệ thống chứng minh được điều gì và dưới những giả thiết nào.
Mọi kết quả đều nằm trong mô hình oracle ngẫu nhiên cổ điển (ROM). Không có định
lý nào trong mô hình oracle ngẫu nhiên lượng tử (QROM) được khẳng định
(Mục~\ref{sec:sec-qrom}). Các chứng minh đầy đủ nằm trong
Phụ lục~\ref{app:security-proofs}; việc kiểm tra chéo từng định nghĩa đối chiếu
với phần hiện thực nằm trong Phụ lục~\ref{app:code-check}.

\subsection{Ký hiệu, tham số và giả thiết}
\label{sec:sec-notation}

Đặt $\mathcal R=\mathbb Z[X]/(X^N+1)$ và $\Rq=\mathcal R/q\mathcal R$, với
$N=256$ và số nguyên tố $q=2^{30}-35\equiv5\pmod 8$. Khóa cam kết là
$\mathsf{ck}=(B_1,b_2)$ với $B_1=[I_n\,|\,B_1']\in\Rq^{n\times k}$,
$b_2=[0^n\,|\,1\,|\,b_2']\in\Rq^{1\times k}$, $k=n+\lambda+1$, $n=6$, $\lambda=7$ (nên $k=14$).
$S_1$ là tập các đa thức có hệ số trong $\{-1,0,1\}$. Với một thông điệp
$m\in\Rq$ và $r\in S_1^k$,
$\Com(m;r)=(t_1,t_2)=(B_1r,\ \langle b_2,r\rangle+m)$. Tập thách thức là
$\Cset=\{c\in\mathcal R:\norm{c}_\infty=1,\ \norm{c}_1=\kappa\}$ với
$\kappa=60$, nên $|\Cset|=\binom{256}{60}2^{60}=2^{257.01}$; tập hiệu của nó
là $\Cdiff=\{c-c':c\neq c'\in\Cset\}$. Các mặt nạ đều (uniform) trên
$Y=[-\gamma+1,\gamma]^{kN}$ với $\gamma=2^{17}$; các đáp ứng được chấp nhận khi
và chỉ khi chúng nằm trong $Z=[-B,B]^{kN}$, $B=\gamma-\kappa-1$. Chúng tôi dùng
$\norm{fv}_\infty\le\norm{f}_1\norm{v}_\infty$ xuyên suốt. Các đại lượng cụ thể
dùng trong các cận được tập hợp trong Bảng~\ref{tab:sec-numbers}.

\begin{table}[t]
\centering
\caption{Các đại lượng cụ thể cho \texttt{bdlop-open-128-v2}}
\label{tab:sec-numbers}
\footnotesize
\begin{tabular}{@{}lr@{}}
\toprule
Quantity & Value \\
\midrule
$|\Cset|$ & $2^{257.01}$\\
Fiat--Shamir seed length & 256 bits\\
$p_{\mathrm{acc}}=((2B+1)/2\gamma)^{kN}$ (opening) & $0.1912$\\
$p_{\mathrm{acc}}^2$ (equality proof) & $0.0365$\\
Attempt cap $\Amax$ (opening; equality) & $512$; $4\,096$\\
$(1-p_{\mathrm{acc}})^{512}$; $(1-p_{\mathrm{acc}}^2)^{4096}$ & $2^{-156.7}$; $2^{-220.0}$\\
$\eta=nN\log_2(2B+1)$ (min-entropy of $w$) & $27\,646.98$ bits\\
$\beta=8\kappa B$ (M-SIS bound) & $62\,885\,280\approx2^{25.91}$\\
$2\kappa+2B$ (one honest, one relaxed opening) & $262\,142\approx2^{18.0}$\\
$\sqrt{q/2}$ (invertibility threshold) & $23\,170.5$\\
$3^{kN}q^{-(n+1)N}$ & $\approx2^{-48\,079}$\\
Statement $x$ & 274 bytes\\
Envelope, in-band / registered & 14\,885 / 8\,165 bytes\\
\bottomrule
\end{tabular}
\end{table}

\begin{lemma}[Hiệu của các thách thức là khả nghịch]
\label{lem:invertible}
Mọi $f\in\Cdiff$ đều khả nghịch trong $\Rq$.
\end{lemma}
\begin{proof}
Vì $q\equiv5\pmod8$, $X^N+1$ tách theo modulo $q$ thành đúng hai nhân tử bất
khả quy, và mọi $y\in\Rq$ khác không với $\norm{y}_\infty<q^{1/2}/\sqrt2$ đều
khả nghịch (Lyubashevsky--Seiler~\cite{ref:ls18}, Hệ quả~1.2 với $k=2$). Mọi
$f\in\Cdiff$ đều khác không với $\norm{f}_\infty\le2<23\,170$. Hàm khởi tạo của
bộ tham số bắt buộc cả hai điều kiện (\texttt{params.py:49--51},
\texttt{ring.py:52--58}).
\end{proof}

\begin{definition}[Các phần mở]
\label{def:relaxed}
Một \emph{phần mở nới lỏng} của $t=(t_1,t_2)$ là một cặp $(\bar r,f)$ với
$f\in\Cdiff$, $\norm{\bar r}_\infty\le2B$ và $ft_1=B_1\bar r$. Nó xác định thông
điệp $\bar m=t_2-f^{-1}\langle b_2,\bar r\rangle$, là $\bar m$ duy nhất thỏa
$ft_2=\langle b_2,\bar r\rangle+f\bar m$ (Bổ đề~\ref{lem:invertible}). Một
\emph{phần mở trung thực} là $(r,1)$ với $r\in S_1^k$ và $t_1=B_1r$, có thông
điệp $t_2-\langle b_2,r\rangle$; đó là thứ mà một lần tiết lộ bộc lộ ra.
\end{definition}

\paragraph{Các giả thiết}
$\Adv^{\MSIS}$ là lợi thế của việc tìm $x\in\mathcal R^k$ với
$B_1x\equiv0\pmod q$, $0<\norm{x}_\infty\le\beta=8\kappa B$, cho $B_1'$ đều.
$\Adv^{\MLWE}$ là lợi thế của việc phân biệt $(A,As+e)$ với đều, cho
$A\in\Rq^{(n+1)\times\lambda}$ đều và $s\in S_1^\lambda$, $e\in S_1^{n+1}$.
Đây là hai thể hiện được ước lượng ở 140/127 và 164/150 bit (core-SVP cổ
điển/lượng tử, Bảng~\ref{tab:params}); ước lượng M-SIS lượng tử chỉ vừa dưới
128 bit. Mọi nghiệm M-SIS trong mục này đều có chuẩn tối đa $\beta$, nên không
cần thêm giả thiết lưới nào nữa.
$\Adv^{\mathsf{cr}}_{\mathsf{SHA2}}$ và $\Adv^{\mathsf{cr}}_{\mathsf{SHA3}}$ là
các lợi thế tìm va chạm đối với SHA-256 và SHA3-256, và
$\Adv^{\mathsf{euf\text{-}cma}}_{\mathsf{ML\text{-}DSA}}$ là lợi thế EUF-CMA đối
với ML-DSA-65~\cite{ref:fips204}.

\paragraph{Các oracle ngẫu nhiên}
Chúng tôi mô hình hóa như các oracle ngẫu nhiên độc lập các thể hiện SHAKE-256
được dùng cho hạt giống thách thức Fiat--Shamir $\Hfs$ (và $H_{\mathsf{eq}}$ cho
chứng minh đẳng thức), phép mở rộng thách thức $G:\{0,1\}^{256}\to\Cset$
(\textsf{SampleInBall}), luồng mặt nạ, phép mở rộng khóa, và hàm băm phòng hộ
(hedging) SHA3-256. Các thẻ miền của chúng là không-tiền-tố đôi một
(Phụ lục~\ref{app:code-check}), nên một oracle SHAKE-256 tách thành các oracle
độc lập. $G$ đều chính xác trên $\Cset$ với một đầu vào mới: các dấu và các chỉ
số Fisher--Yates được rút bằng lấy mẫu loại bỏ chính xác. Phép mở rộng khóa
dưới dạng một oracle ngẫu nhiên làm cho $\mathsf{ck}$ đều; đó là một khóa toàn
cục duy nhất được dẫn xuất từ một hạt giống công khai (không có khóa riêng cho
từng người dùng). $Q_H$ ký hiệu tổng số truy vấn oracle ngẫu nhiên của đối thủ,
bao gồm cả những truy vấn được thực hiện thay mặt nó bởi các truy vấn kiểm
chứng, cộng thêm một.

\subsection{Cú pháp của lược đồ hợp thành}
\label{sec:sec-syntax}

Tầng vận chuyển là một tập các hàm tất định của một hồ sơ vật mang đã ghim
$\mathsf{prof}$ và một khóa đồng bộ $k_s$: $\mathsf{Place}(\vid,\ell)$ trả về
các vị trí vật mang có thứ tự $P$ cho một envelope $\ell$ byte;
$\Canon(\vid,P)$ trả về digest đã chuẩn hóa theo vật mang $D$;
$\mathsf{Write}(\vid,P,E)$ ghi envelope $E$ vào các vật mang;
$\mathsf{Read}(\vid)$ trả về $(E,P,D)$ hoặc $\bot$. \emph{Tính đúng đắn vận
chuyển} nghĩa là $\mathsf{Read}(\mathsf{Write}(\vid,P,E))=(E,P,\Canon(\vid,P))$
với $P=\mathsf{Place}(\vid,|E|)$; đó là một tính chất hàm, và việc nó thất bại
chỉ gây ra từ chối.

\begin{itemize}
\item $\mathsf{Setup}$: cố định bộ tham số, $\pid$ (một digest SHA3-256 của tất
  cả các tham số và hạt giống thiết lập) và $\mathsf{ck}$.
\item $\mathsf{RegKeyGen}$, $\mathsf{RegSign}$: bên phát hành nắm một cặp khóa
  ML-DSA-65 $(\mathit{pk}_I,\mathit{sk}_I)$ và ký một sổ đăng ký
  $\mathsf{Reg}=(e,\Delta)$: một epoch và một tập các mô tả quan hệ. Một mô tả
  băm $\pid$, khóa, module quan hệ và chính sách (chính sách codec và hồ sơ vật
  mang); trong chế độ đã đăng ký nó còn băm một cam kết chứng chỉ $t_c$.
\item $\mathsf{Pins}$: cho $\mathsf{Reg}$, $\mathit{pk}_I$, $\mathsf{prof}$,
  một ngưỡng chống quay lui $e_{\min}$ và (chế độ đã đăng ký) $t_c$, kiểm tra
  chữ ký, tính lại mô tả kỳ vọng, đòi hỏi đúng một kết quả khớp, và xuất ra
  $\mathsf{pins}=(\mathit{rid},\mathit{root},e,\mathit{cpid},
  \mathit{prof\_id})$, hoặc thất bại.
\item $\mathsf{Issue}$: rút một thách thức đều 32 byte $N_s$ và lưu nó cùng ràng
  buộc $\mathsf{bind}(\mathsf{pins})$ và một thời điểm hết hạn.
\item $\mathsf{Embed}$: trên vật phủ $\vid$, tải trọng $p$ (in-band) hoặc chứng
  chỉ $(t_c,r_c)$ (đã đăng ký), $\mathsf{pins}$ và $N_s$: tính
  $P=\mathsf{Place}(\vid,\ell)$ từ một placeholder bằng không có độ dài cuối
  cùng $\ell$, $D=\Canon(\vid,P)$ và $x=\Stmt(\mathsf{pins},N_s,P,D)$; với
  in-band, lấy mẫu $r\leftarrow S_1^k$, đặt $t=\Com(\mathsf{enc}(p);r)$; chứng
  minh $\pi\leftarrow\mathsf{P}(x,t,r)$; lập $E=\texttt{"ZKLB"}\|1\|\pid\|N_s\|t\|\pi$
  (hoặc $\texttt{"ZKLR"}\|1\|\pid\|N_s\|\pi$) và xuất ra
  $\mathsf{Write}(\vid,P,E)$ và, với in-band, phần mở $(p,r)$.
\item $\mathsf{Verify}$: trên $\vid^*$: $(E^*,P^*,D^*)=\mathsf{Read}(\vid^*)$;
  phân tích $E^*$ một cách nghiêm ngặt; lấy $N_s^*$ từ $E^*$ (chế độ có lưu
  trữ) hoặc dùng một $N_s$ cố định (chế độ không trạng thái);
  $x^*=\Stmt(\mathsf{pins},N_s^*,P^*,D^*)$; chấp nhận khi và chỉ khi
  $\mathsf{V}(x^*,t^*,\pi^*)=1$ với $t^*$ từ $E^*$ (in-band) hoặc $t_c$ đã ghim
  (đã đăng ký) và, trong chế độ có lưu trữ, $N_s^*$ còn sống đối với
  $\mathsf{bind}(\mathsf{pins})$ và được tiêu thụ một cách nguyên tử.
\item $\mathsf{VerDisc}$: chấp nhận $(E,p,r)$ khi và chỉ khi $(r,1)$ là một phần
  mở trung thực của cam kết trong $E$ tới $\mathsf{enc}(p)$.
\item $\mathsf{Link}$, $\mathsf{VerLink}$: chứng minh đẳng thức
  $\Pi_{\mathsf{eq}}$ trên hai cam kết, được ràng buộc vào một digest ngữ cảnh
  $\mathit{lx}=\mathrm{SHA3}(\text{tag}\|\pid\|\text{ngữ cảnh có tiền tố độ dài})$.
\end{itemize}

Phát biểu là chuỗi 274 byte
\begin{multline*}
\Stmt=\mathtt{u16}(1)\|\mathit{rid}\|\pid\|\mathit{root}\|\mathtt{u64}(e)\|N_s\|
\mathit{cpid}\|\mathit{prof\_id}\\
\|\mathtt{u64}(|P|)\|\mathrm{SHA256}(P)\|V,\quad
V=\mathrm{SHA256}(\text{tag}\|D),
\end{multline*}
với $\mathrm{SHA256}(P)$ được lấy trên một mã hóa JSON chuẩn tắc của các vị trí
có thứ tự. Chứng minh mở $\Pi_{\mathsf{open}}$ xuất ra $\pi=(h,z)$ với
$h=\Hfs(\text{tag}\|\pid\|\mathtt{u64}(|x|)\|x\|\mathsf{enc}(t)\|\mathsf{enc}(w))$,
$w=B_1y$, $c=G(h)$, $z=y+cr\in Z$; bên kiểm chứng tính lại $w'=B_1z-ct_1$ và
kiểm tra $h$. Mọi trường sau $x$ đều có độ dài được cố định bởi $\pid$, nên đầu
vào của hàm băm giải mã được một cách duy nhất.

\subsection{Mô hình đối thủ}
\label{sec:sec-adversary}

Đối thủ $\mathcal A$ là một thuật toán xác suất cổ điển. Nó kiểm soát mọi kênh,
chọn các vật phủ và các tải trọng, biết $k_s$ (không tính chất nào dưới đây phụ
thuộc vào tính bí mật của nó), và có các oracle sau. $\mathcal O_{\mathsf{Iss}}$
trả về một $N_s$ mới. $\mathcal O_{\mathsf{Emb}}(\vid,p,N_s)$ chạy một
$\mathsf{Embed}$ in-band trung thực và trả về video stego;
$\mathcal O_{\mathsf{EmbR}}(j,\vid,N_s)$ làm điều tương tự cho chứng chỉ đã đăng
ký của bên nắm giữ trung thực $j$. Chúng tôi cho phép các oracle này chứng minh
các phát biểu $x$ tùy ý do $\mathcal A$ chọn, điều này chỉ làm mạnh thêm các kết
quả. Các đầu ra của chúng tạo thành danh sách $L$ các bộ ba $(x,t,\pi)$.
$\mathcal O_{\mathsf{Ver}}$ chạy $\mathsf{Verify}$ với kho dùng chung.
$\mathcal O_{\mathsf{Link}}$ trả về các chứng minh liên kết trung thực giữa các
cam kết đã ghi. $\mathcal O_{\mathsf{Disc}}(i)$ trả về phần mở của cam kết trung
thực thứ $i$ và đánh dấu nó là \emph{đã tiết lộ}.
$\mathcal O_{\mathsf{Reg}}(e,\Delta,T)$ khiến bên phát hành ký một sổ đăng ký
với epoch $e$ mà các mô tả của nó là các mô tả in-band $\Delta$ và các mô tả đã
đăng ký mà bên phát hành tính từ các cam kết chứng chỉ trong $T$, tất cả do
$\mathcal A$ chọn; $T_{\mathsf{reg}}$ là hợp của mọi $T$ đã truy vấn. (Một bên
phát hành chỉ ký một mô tả đã đăng ký cho một chứng chỉ mà nó đã thẩm định, nên
các mô tả đã đăng ký chỉ đi vào qua $T$.) $Q_P$ chặn số chứng minh được tạo bởi
các bên trung thực.

Bên phát hành và bên kiểm chứng là trung thực; kho phiên là tuyến tính hóa được
(linearizable) và bền vững. Việc làm hỏng bên nắm giữ là tĩnh (đối thủ đăng ký
các chứng chỉ của riêng nó qua $\mathcal O_{\mathsf{Reg}}$); tiết lộ là hình
thức làm hỏng thích nghi duy nhất, và các cam kết đã tiết lộ bị loại khỏi các
điều kiện thắng. Các kênh phụ, lỗi codec, và tính đúng đắn của tầng vận chuyển
nằm ngoài phạm vi.

\subsection{Tính không tiết lộ tri thức của chứng minh Fiat--Shamir-with-aborts}
\label{sec:sec-zk}

\begin{lemma}[Chấp nhận độc lập với nhân chứng]
\label{lem:hvzk}
Với mọi $r\in S_1^k$ và $c\in\Cset$, một lần thử với $y\leftarrow Y$ được chấp
nhận với xác suất đúng bằng $p_{\mathrm{acc}}$, và với điều kiện được chấp nhận
thì $z=y+cr$ đều trên $Z$.
\end{lemma}
\noindent
Mỗi hệ số $v$ của $cr$ thỏa $|v|\le\kappa$, nên $2B+1$ giá trị
$z_i\in[-B,B]$ tương ứng với các mặt nạ $z_i-v\in[-\gamma+1,\gamma-1]\subset
[-\gamma+1,\gamma]$, mỗi cái có xác suất $1/2\gamma$.

\begin{lemma}[Mô phỏng]
\label{lem:sim}
Tồn tại một bộ mô phỏng mà, chỉ với $(x,t)$ và quyền lập trình $\Hfs$, xuất ra
các chứng minh có phân phối hợp với các oracle ngẫu nhiên nằm trong khoảng
\begin{multline*}
\epszk(Q_P,Q_H)=2Q_P(Q_H+Q_P)2^{-256}\\
+(\Amax+1)Q_P\bigl(Q_H+(\Amax+1)Q_P\bigr)2^{-\eta}
\end{multline*}
so với $Q_P$ chứng minh thực, với mọi đối thủ thực hiện $Q_H$ truy vấn. Điều
tương tự đúng cho $\Pi_{\mathsf{eq}}$ với $p_{\mathrm{acc}}^2$ thay cho $p_{\mathrm{acc}}$.
\end{lemma}
\noindent
Bộ mô phỏng lấy mẫu số lần thử từ phân phối hình học bị cắt cụt với tham số
$p_{\mathrm{acc}}$, xuất ra $\bot$ với xác suất $(1-p_{\mathrm{acc}})^{\Amax}$,
và ngược lại lấy mẫu $h\leftarrow\{0,1\}^{256}$, $z\leftarrow Z$, đặt
$w=B_1z-G(h)t_1$ và lập trình $\Hfs$ tại $(\pid,x,t,w)$. Bước tế nhị là các truy
vấn oracle ngẫu nhiên mà bên chứng minh thực thực hiện trong các lần thử bị loại
và không bao giờ tiết lộ: các phân tích trước đây về Fiat--Shamir with aborts đã
xử lý sai việc tái lập trình tương ứng trong bước CMA-to-NMA, như được Barbosa
và cộng sự~\cite{ref:barbosa} chỉ ra và sửa chữa, và được Devevey và cộng
sự~\cite{ref:devevey} phân tích chi tiết. Ở đây bước này được giải quyết trong
ROM cổ điển nhờ hai sự kiện đặc thù của hiện thực này: sự chấp nhận là
\emph{hoàn toàn} độc lập với nhân chứng và với $c$ (Bổ đề~\ref{lem:hvzk}), và
mọi điểm truy vấn ẩn đều chứa $w=B_1y$ có min-entropy ít nhất $\eta$ (bởi vì
$B_1=[I_n|B_1']$). Tính độc lập chính xác cũng có nghĩa là số lần thử không rò
rỉ gì, dù rằng ngoài ra phần hiện thực không chạy trong thời gian hằng.

\subsection{Tính đúng đắn tri thức và ràng buộc vào video}
\label{sec:sec-soundness}

\begin{definition}[Ràng buộc video trích xuất được qua mô phỏng]
\label{def:se}
Trong thí nghiệm $\mathsf{SE}$, $\mathcal A$ tương tác với mọi oracle và xuất
ra $\vid^*$. Gọi $x^*$ là phát biểu mà bên kiểm chứng tính lại từ $\vid^*$ và
các mục ghim của chính nó, và $(t^*,\pi^*)$ là cam kết và chứng minh mà nó kiểm
chứng. $\mathcal A$ \emph{thắng} nếu $\mathsf{Verify}$ chấp nhận và
$(x^*,t^*,\pi^*)\notin L$. Một bộ trích xuất $\mathcal E^{\mathcal A}$
\emph{thành công} nếu nó xuất ra $x^*$ và một phần mở nới lỏng của $t^*$, trong
đó $x^*$ là phát biểu của bên kiểm chứng trong một lần chạy thắng của $\mathcal A$.
\end{definition}

\begin{theorem}[Tính đúng đắn tri thức ràng buộc vào phát biểu]
\label{thm:se}
Giả sử $\mathcal A$ thắng $\mathsf{SE}$ với xác suất $\varepsilon$. Tồn tại một
bộ trích xuất $\mathcal E$ chạy $\mathcal A$ hai lần với các oracle được mô
phỏng và thành công với xác suất $\varepsilon_{\mathsf{ext}}$ sao cho
\begin{multline*}
\varepsilon\le\epszk(Q_P,Q_H)+\Adv^{\MSIS}+Q_H2^{-256}\\
+\sqrt{Q_H\bigl(\varepsilon_{\mathsf{ext}}+Q_H2^{-256}+2^{-257}\bigr)}.
\end{multline*}
\end{theorem}
\noindent
\emph{Ý tưởng chứng minh.} Thay các chứng minh trung thực bằng các chứng minh
được mô phỏng (Bổ đề~\ref{lem:sim}). Nếu $\mathcal A$ tái dùng một điểm đã lập
trình, tức $(x^*,t^*)=(x_i,t_i)$ và $w^*=w_i$, thì $h^*=h_i$, do đó $c^*=c_i$ và
$B_1(z^*-z_i)=0$ với $0<\norm{z^*-z_i}_\infty\le2B$: một nghiệm M-SIS. Ngược
lại, truy vấn $(\pid,x^*,t^*,w^*)$ của bên kiểm chứng được trả lời bằng một giá
trị ngẫu nhiên mới, và bổ đề rẽ nhánh tổng quát của Bellare và
Neven~\cite{ref:bn06} (Bổ đề~1) với không gian đầu ra $2^{256}$ cho hai chứng
minh được chấp nhận có cùng $(x^*,t^*,w^*)$ và các hạt giống $h\neq h'$. Khi đó
$c=G(h)\neq G(h')=c'$ ngoại trừ với xác suất $2^{-257}+Q_H2^{-256}$, và
$B_1(z-z')=(c-c')t_1^*$ cho phần mở nới lỏng $(z-z',c-c')$. Hiệu thách thức là
khả nghịch theo Bổ đề~\ref{lem:invertible}, chính điều này khiến phần mở xác
định một thông điệp duy nhất. Bởi vì truy vấn bị rẽ nhánh chứa $x^*$, phần mở
được ràng buộc vào đúng phát biểu do bên kiểm chứng tính lại: cam kết video, các
vị trí vật mang, thách thức phiên, gốc sổ đăng ký và epoch, quan hệ và chính
sách của nó. \qed

Tổn thất căn bậc hai là cái giá thường thấy của việc tua lại (rewinding). Với
$Q_H=2^{64}$, cận biến một xác suất trích xuất $2^{-128}$ thành $\varepsilon
\lesssim2^{-32}$. Đối với tính xác thực (Định lý~\ref{thm:auth}), có một cận
chặt nhờ hướng SelfTargetMSIS của Kiltz, Lyubashevsky và
Schaffner~\cite{ref:kls18}, với cái giá là giả thiết phi chuẩn đó.

\begin{theorem}[Tính ràng buộc đối với các phần mở nới lỏng]
\label{thm:binding}
Không đối thủ nào có thể xuất ra $t$ và hai phần mở của $t$, mỗi phần mở trung
thực hoặc nới lỏng, với các thông điệp khác nhau, ngoại trừ với xác suất
$\Adv^{\MSIS}$.
\end{theorem}
\noindent
Cho $(\bar r,f,m)$ và $(\bar r',f',m')$, vector
$x=f'\bar r-f\bar r'$ thỏa $B_1x=0$; nếu $x=0$ thì
$\langle b_2,x\rangle=ff'(m-m')=0$, mâu thuẫn với tính khả nghịch. Chuẩn của nó
tối đa là $2\cdot2\kappa\cdot2B=\beta$ cho hai phần mở nới lỏng, $2\kappa+2B$
cho một trung thực và một nới lỏng, và $2$ cho hai trung thực, tất cả đều dưới
$q/2$ (\texttt{params.py:53}), nên $x\not\equiv0$ là một nghiệm M-SIS.

\noindent
Gộp lại, các Định lý~\ref{thm:se} và~\ref{thm:binding} nói rằng: ai tạo ra một
video in-band được chấp nhận đều bị ràng buộc vào thông điệp $\bar m$ mà bộ
trích xuất tính ra, và bất kỳ lần tiết lộ về sau nào mà $\mathsf{VerDisc}$ chấp
nhận đều bộc lộ một tải trọng $p$ với $\mathsf{enc}(p)=\bar m$, trừ khi M-SIS
bị phá.

\emph{Việc tiết lộ giữ cho tính ngẫu nhiên riêng tư.} Một lần tiết lộ là tải
trọng $p$ cùng với một chứng minh đẳng thức giữa $t$ và cam kết tất định
$\mathsf{Com}(\mathsf{enc}(p);0)$, được Fiat--Shamir ràng buộc vào một ngữ cảnh
do bên nhận chọn (định danh phát biểu và nonce). Bên nhận biết $p$ nhưng không
biết phần mở nào của $t$, nên nó không thể tạo các chứng minh mở mới cho $t$;
theo tính đúng đắn của chứng minh đẳng thức, một lần tiết lộ được chấp nhận vẫn
thỏa $\mathsf{enc}(p)=\bar m$ trừ khi M-SIS bị phá. Một phiên bản trước đây bộc
lộ chính $r$; đây là mục~D3 trong Phụ lục~\ref{app:code-check} và đã được sửa.
Các định lý dưới đây vẫn coi các cam kết mà \emph{phần mở} của chúng đã bị rò rỉ
(ví dụ do lộ khóa) là đã tiết lộ.

\begin{lemma}[Tính một chiều của các cam kết trung thực]
\label{lem:row}
Cho $\mathsf{ck}$ và $t^{(j)}=\Com(m_j;r_j)$, $j\le Q$, với
$r_j\leftarrow S_1^k$ và các thông điệp $m_j$ được chọn thích nghi trước khi mỗi
$t^{(j)}$ được rút, không đối thủ nào xuất ra một phần mở nới lỏng của bất kỳ
$t^{(j)}$ nào ngoại trừ với xác suất
$\Adv^{\mathsf{row}}_Q=\Adv^{\MSIS}+Q\bigl(\Adv^{\MLWE}+3^{kN}q^{-(n+1)N}\bigr)$.
\end{lemma}
\noindent
Đặt $x=fr_j-\bar r$, nên $B_1x=0$ và $\norm{x}_\infty\le2\kappa+2B$. Nếu
$x\neq0$ thì nó là một nghiệm M-SIS (phép rút gọn đã sinh mọi $t^{(j)}$ và biết
$r_j$). Nếu $x=0$ thì đối thủ đã tìm được $r_j=f^{-1}\bar r$, điều mà một bộ
phân biệt cấy mẫu M-LWE của nó tại một $j$ được đoán sẽ kiểm tra đối chiếu với
$t^{(j)}$; với $t^{(j)}$ đều, một phần mở bậc ba tồn tại với xác suất tối đa
$3^{kN}q^{-(n+1)N}$. Chỉ riêng tính ràng buộc là không đủ ở đây: một đối thủ
khôi phục được $r_j$ sẽ mở $t^{(j)}$ tới đúng thông điệp.

\subsection{Che giấu tải trọng}
\label{sec:sec-hiding}

\begin{definition}[Tính không phân biệt được tải trọng]
\label{def:hiding}
Trong $\mathsf{IND}^b$, $\mathcal A$ có mọi oracle và nộp một lần
$(\vid,p_0,p_1,N_s)$ với $|p_0|,|p_1|\le766$ byte; bên thách thức trả về
$\mathcal O_{\mathsf{Emb}}(\vid,p_b,N_s)$. $\mathcal A$ không được truy vấn
$\mathcal O_{\mathsf{Disc}}$ hoặc $\mathcal O_{\mathsf{Link}}$ trên cam kết
thách thức. $\Adv^{\mathsf{ind}}=|\Pr[\mathsf{IND}^1=1]-\Pr[\mathsf{IND}^0=1]|$.
\end{definition}

\begin{theorem}[Che giấu]
\label{thm:hiding}
$\Adv^{\mathsf{ind}}\le2\,\epszk(Q_P+1,Q_H)+2\Adv^{\MLWE}$.
\end{theorem}
\noindent
Video stego chỉ phụ thuộc vào $b$ thông qua $E$: $P$ và $D$ được cố định từ một
placeholder trước khi chứng minh, và envelope có kích thước cố định, nên độ dài
tải trọng cũng được che giấu. Mô phỏng chứng minh thách thức
(Bổ đề~\ref{lem:sim}), rồi thay $t$ bằng một cặp đều: với
$B_1'=[B_{1,2}'\,|\,B_{1,3}']$ và $r=(r_1,r_2,r_3)$,
$(t_1-B_{1,2}'(t_2-m),\,t_2-m)=A r_3+(r_1,r_2)$ với
$A=\bigl[\begin{smallmatrix}B_{1,3}'-B_{1,2}'b_2'\\ b_2'\end{smallmatrix}\bigr]$
đều, một mẫu M-LWE hạng $\lambda$ với $n+1$ mẫu vành; ánh xạ là một song ánh với
mỗi $m$ cố định. Khi đó thí nghiệm độc lập với $b$. Tính che giấu mang tính tính
toán và là tính chất duy nhất dựa trên M-LWE; nó đúng cho tải trọng, chứ không
cho sự kiện rằng một tải trọng tồn tại, vốn là một tính chất giấu tin được đánh
giá bằng thực nghiệm.

\subsection{Khả năng chống phát lại}
\label{sec:sec-replay}

\begin{definition}[Phát lại]
\label{def:replay}
Trong $\mathsf{RPL}$, $\mathcal A$ có mọi oracle đối với một bên kiểm chứng có
kho phiên và thắng nếu (i) hai lời gọi $\mathcal O_{\mathsf{Ver}}$ chấp nhận với
cùng $N_s$; hoặc (ii) một lời gọi chấp nhận với một $N_s$ không được phát hành
cho $\mathsf{bind}(\mathsf{pins})$, hoặc đã hết hạn; hoặc (iii) một lời gọi chấp
nhận $\vid^*$ có bộ ba $(x^*,t^*,\pi^*)\notin L$ và $\mathcal A$ đã truy vấn
$\Hfs$ tại điểm $(\pid,x^*,t^*,w^*)$ của bên kiểm chứng trước khi $N_s^*$ được
phát hành.
\end{definition}

\begin{theorem}[Khả năng chống phát lại]
\label{thm:replay}
$\Pr[\text{(i) hoặc (ii)}]\le\Adv^{\mathsf{cr}}_{\mathsf{SHA2}}+Q_I^22^{-257}$ và
$\Pr[\text{(iii)}]\le Q_HQ_I2^{-256}$, với $Q_I$ phiên được phát hành.
\end{theorem}
\noindent
(i) và (ii) suy ra từ ngữ nghĩa của kho: các hàng chỉ tồn tại cho các thách
thức đã phát hành (khóa theo một digest SHA-256 của $N_s$), và việc tiêu thụ là
một \texttt{UPDATE \dots\ WHERE consumed\_at IS NULL AND expires\_at>now} có
điều kiện đơn nhất bên trong một giao dịch ghi, nên nhiều nhất một bên gọi quan
sát được một hàng đã cập nhật (\texttt{zkp\_sessions.py:181--241}). Với (iii),
$N_s^*$ là một chuỗi con của $x^*$, nên một truy vấn thực hiện trước đó dự đoán
một giá trị 256 bit đều. Cùng với Định lý~\ref{thm:se}, một phiên được chấp nhận
tương ứng với một bên đã biết một phần mở và đã tính chứng minh sau khi phiên
được phát hành.

\emph{Điều mà một bên kiểm chứng không trạng thái không thể đem lại.} Với một
$N_s$ cố định (\texttt{expected\_session}), (i) thất bại: bất kỳ bản sao nào của
một video được chấp nhận đều lại được chấp nhận, bởi bất kỳ ai và bao nhiêu lần
tùy ý, vì một chứng minh không tương tác có thể chuyển nhượng. Nếu $N_s$ bị tái
dùng qua nhiều lần kiểm chứng, (iii) cũng thất bại và chỉ còn lại ràng buộc
video (Định lý~\ref{thm:se}). Không biến thể nào ngăn được việc tiếp chuyển
(relay): chế độ in-band không ràng buộc $N_s$ vào một bên chứng minh, nên ai
trình ra một video hợp lệ trước thì tiêu thụ phiên.

\subsection{Tính xác thực chế độ đã đăng ký}
\label{sec:sec-auth}

\begin{definition}[Tính xác thực]
\label{def:auth}
Trong $\mathsf{AUTH}$, bên thách thức tạo các chứng chỉ trung thực
$t_c^{(j)}=\Com(m_j;r_j)$, $j\le Q_C$, với $m_j$ do $\mathcal A$ chọn, và
$\mathcal A$ có mọi oracle. $\mathcal A$ xuất ra $(\mathsf{Reg}^*,\vid^*,t^*)$
và thắng nếu một bên kiểm chứng ghim $\mathit{pk}_I$ và $t^*$ chấp nhận $\vid^*$
với các mục ghim từ $\mathsf{Reg}^*$, bộ ba $(x^*,t^*,\pi^*)\notin L$, và hoặc
(a) $t^*\notin T_{\mathsf{reg}}$ (chưa bao giờ được bên phát hành đăng ký), hoặc
(b) $t^*=t_c^{(j)}$ cho một $j$ trung thực chưa tiết lộ.
\end{definition}

\begin{theorem}[Tính xác thực]
\label{thm:auth}
\begin{multline*}
\Adv^{\mathsf{auth}}\le\Adv^{\mathsf{euf\text{-}cma}}_{\mathsf{ML\text{-}DSA}}
+\Adv^{\mathsf{cr}}_{\mathsf{SHA2}}+\epszk(Q_P,Q_H)\\
+\Adv^{\MSIS}+Q_H2^{-256}\\
+\sqrt{Q_H\bigl(\Adv^{\mathsf{row}}_{Q_C}+Q_H2^{-256}+2^{-257}\bigr)}.
\end{multline*}
\end{theorem}
\noindent
Trường hợp (a) cho một chữ ký giả mạo ML-DSA trên các byte ký của sổ đăng ký
hoặc, nếu $\mathsf{Reg}^*$ đã được ký, một va chạm SHA-256 giữa các đầu vào có
độ dài cố định $B_1'\|b_2'\|t^*$ và $B_1'\|b_2'\|t'$ với một
$t'\in T_{\mathsf{reg}}$ nào đó, vì mô tả khớp chứa digest đó. Trong trường hợp
(b), Định lý~\ref{thm:se} áp dụng cho trò chơi được mô phỏng trích xuất một phần
mở nới lỏng của một chứng chỉ trung thực, điều mà Bổ đề~\ref{lem:row} loại trừ.
Định lý này nói về các chứng chỉ; sự liên kết tới một con người chỉ mạnh ngang
với việc bên phát hành kiểm tra bên nắm giữ trước khi ký. Khi bên phát hành ràng
buộc một định danh của bên nắm giữ, mô tả đã ký chứa một digest có tách miền của
nó bên cạnh cam kết, nên quan hệ đã ghim của bên kiểm chứng nêu tên cả hai
(Phụ lục~\ref{app:code-check}, mục~D2, đã sửa).

\subsection{Liên kết}
\label{sec:sec-linkage}

\begin{definition}[Tính đúng đắn và tính che giấu của liên kết]
\label{def:link}
Trong $\mathsf{LNK\text{-}S}$, $\mathcal A$ xuất ra một ngữ cảnh, $t_a,t_b$, một
$\pi_{\mathsf{eq}}$ được chấp nhận không do $\mathcal O_{\mathsf{Link}}$ xuất ra
cho bộ ba đó, và các phần mở (trung thực hoặc nới lỏng) của $t_a$ và $t_b$ với
các thông điệp khác nhau. Trong $\mathsf{LNK\text{-}H}^b$, $\mathcal A$ nộp
$(p_0,p_1,\mathit{ctx})$ và nhận $t_a=\Com(p_b;r_a)$,
$t_b=\Com(p_b;r_b)$ và $\pi_{\mathsf{eq}}$ cho $\mathit{ctx}$.
\end{definition}

\begin{theorem}[Liên kết]
\label{thm:link}
$\Adv^{\mathsf{lnk\text{-}s}}$ tuân theo cận của Định lý~\ref{thm:se} với
$\varepsilon_{\mathsf{ext}}\le\Adv^{\MSIS}$, và
$\Adv^{\mathsf{lnk\text{-}h}}\le2\,\epszk(Q_P+1,Q_H)+4\Adv^{\MLWE}$.
\end{theorem}
\noindent
Hai bản ghi bị rẽ nhánh với một thông báo $(w_a,w_b,u)$ cho các phần mở nới lỏng
$(\bar r_a,f)$, $(\bar r_b,f)$ với một $f$ \emph{chung} và
$f(t_{2,a}-t_{2,b})=\langle b_2,\bar r_a-\bar r_b\rangle$, do đó
$\bar m_a=\bar m_b$; hai phần mở của đối thủ khi đó mâu thuẫn với
Định lý~\ref{thm:binding} cho $t_a$ hoặc cho $t_b$. Việc liên kết tới một chứng
chỉ đã đăng ký còn xác thực thêm bên nắm giữ của nó, như trong
Định lý~\ref{thm:auth}, bởi vì việc trích xuất cũng cho ra một phần mở nới lỏng
của $t_c$. Tính che giấu dùng bộ mô phỏng cho $\Pi_{\mathsf{eq}}$ và hai bước
lai M-LWE; sự bằng nhau của hai tải trọng được bộc lộ theo thiết kế.

\subsection{Tính toàn vẹn của nội dung video}
\label{sec:sec-integrity}

$\Canon$ chia luồng bằng bộ phân đoạn đã ghim (\texttt{ffmpeg -c:v copy
-f segment}, $F$ khung mỗi đoạn) và kết hợp các digest của từng đoạn với
$F$ và các số đếm khung. Một đoạn mang các chunk envelope được băm qua một phép
tuần tự hóa có gắn thẻ, phân định bằng độ dài: với mỗi đơn vị NAL là chỉ số,
kiểu, \texttt{forbidden\_zero\_bit}, \texttt{nal\_ref\_idc} và kích thước mã
khởi đầu (start-code); RBSP nguyên văn cho các đơn vị NAL không có vật mang; và,
bên trong một đơn vị NAL IDR có vật mang, các bit RBSP giữa các khối vật mang
cho tới bit dừng (stop bit), với mỗi khối vật mang được thay bằng vector hệ số
của nó trong đó mỗi vật mang được chuẩn hóa (độ lớn làm tròn xuống số chẵn, giữ
dấu). Các đoạn khác được băm từng byte.

\begin{lemma}[Digest và envelope xác định những gì]
\label{lem:canon}
Phép tuần tự hóa giải mã được một cách duy nhất, nên $V$ bằng nhau kéo theo
$\Canon$ bằng nhau ngoại trừ với xác suất $3\Adv^{\mathsf{cr}}_{\mathsf{SHA2}}$.
Dưới quy tắc vật mang ổn định ($|\ell|\ge4$, một bit chẵn-lẻ), một mức vật mang
là $\mathrm{sgn}(\ell)(2\lfloor|\ell|/2\rfloor+b)$, nên $(\Canon,P)$ và các bit
vật mang xác định toàn bộ vector hệ số của mọi khối vật mang. Nếu chuỗi bit vật
mang là mã hóa duy nhất của $E$ dưới hồ sơ đã ghim, thì $(\Canon,P,E)$ xác định:
sự phân đoạn và các số đếm khung; mọi byte của các đoạn không có vật mang; mọi
trường header NAL và bit RBSP bên ngoài các khối vật mang; và các vector hệ số
của các khối vật mang.
\end{lemma}

\begin{theorem}[Tính toàn vẹn]
\label{thm:integrity}
Giả sử một lần nhúng trung thực tạo ra $\vid_h$ cho phiên $N_s$ với cam kết $t$
(hoặc chứng chỉ $t_c$), và $\mathcal A$ xuất ra $\vid^*$ được chấp nhận cho
$N_s$ với cùng cam kết (hoặc chứng chỉ), không tiết lộ. Khi đó
$\Canon(\vid^*)=\Canon(\vid_h)$, $P^*=P_h$ và $E^*=E_h$, ngoại trừ với xác suất
tối đa
\begin{multline*}
4\Adv^{\mathsf{cr}}_{\mathsf{SHA2}}+\epszk(Q_P,Q_H)+\Adv^{\MSIS}+Q_H2^{-256}\\
+\sqrt{Q_H\bigl(\Adv^{\mathsf{row}}_{Q_P}+Q_H2^{-256}+2^{-257}\bigr)}.
\end{multline*}
\end{theorem}
\noindent
Nếu $x^*\neq x_h$, bộ ba là mới và việc trích xuất cho ra một phần mở nới lỏng
của một cam kết trung thực. Nếu $x^*=x_h$, giá trị băm vị trí và $V$ khớp nhau,
và Bổ đề~\ref{lem:canon} được áp dụng; khi đó chứng minh chính là chứng minh
trung thực, ngược lại trường hợp điểm-đã-lập-trình của Định lý~\ref{thm:se} cho
một nghiệm M-SIS.

\paragraph{Các giới hạn của khẳng định toàn vẹn}
\begin{itemize}
\item Chế độ in-band ràng buộc nội dung vào một \emph{cam kết}, chứ không vào
  một con người: bất kỳ ai nắm một $N_s$ còn sống đều có thể nhúng lại bất kỳ
  video đã chỉnh sửa nào với một cam kết mới của riêng nó, và nó sẽ được chấp
  nhận.
\item Các vật mang chỉ được bảo vệ như các bit envelope. Định lý~\ref{thm:integrity}
  cần các bit vật mang là mã hóa duy nhất của $E$. Bên kiểm chứng chia lại chunk
  của $E$ theo kích thước chunk đã ghim và từ chối trừ khi các mảnh trích xuất
  được đúng là mã hóa đó; nếu không có phép kiểm này, một envelope một-chunk sẽ
  để trống trường \texttt{chunk\_size} 16 bit của nó (mục~D1 trong
  Phụ lục~\ref{app:code-check}, đã sửa).
\item Không được ràng buộc: các byte không sống sót qua bộ phân đoạn đã ghim (dữ
  liệu bộ chứa, các luồng khác, dữ liệu ngoài các đơn vị NAL); các bit căn lề
  RBSP sau bit dừng trong các đơn vị NAL IDR có mang vật mang; các từ mã CAVLC
  thay thế cho cùng một vector hệ số của một khối vật mang; tầng chống giả lập
  (digest bao phủ RBSP); và phiên bản của bộ phân đoạn, vốn được chính sách
  codec nêu tên nhưng không ghim.
\item Không có khẳng định về tính đúng đắn codec: digest bao phủ cú pháp, không
  bao phủ các điểm ảnh đã giải mã. Không có khẳng định về nguồn gốc: không gì
  nói rằng một máy quay đã ghi $\vid_h$ hay rằng tải trọng là đúng sự thật.
\end{itemize}

\subsection{Các đối thủ lượng tử}
\label{sec:sec-qrom}

Mọi định lý ở trên đều là các phát biểu trong ROM cổ điển. Việc tua lại
(Định lý~\ref{thm:se}) và việc lập trình thích nghi (Bổ đề~\ref{lem:sim}) không
chuyển sang QROM như đã phát biểu. Barbosa và cộng sự~\cite{ref:barbosa} đưa ra
các chứng minh QROM của bước CMA-to-NMA cho Fiat--Shamir with aborts, và Don và
cộng sự~\cite{ref:dfms19} cùng Liu và Zhandry~\cite{ref:lz19} chỉ ra rằng phép
biến đổi Fiat--Shamir bảo toàn tính đúng đắn (tri thức) của các $\Sigma$-protocol
phù hợp trước các đối thủ lượng tử với tổn thất đa thức. Chúng tôi chưa kiểm
chứng rằng các giả thiết của họ đúng cho bộ trích xuất nới lỏng và bên chứng
minh có loại bỏ được dùng ở đây, nên chúng tôi không khẳng định định lý QROM
nào. Cách diễn đạt của chúng tôi là ``hậu lượng tử một cách hợp lý'': các phép
rút gọn nhắm tới M-SIS, M-LWE, ML-DSA và va chạm băm, không cái nào có một tấn
công lượng tử hiệu quả đã biết ở các tham số này.
, and after \appendices add
% % Full proofs for sections/security.tex and the specification/code
% cross-check. Input after \appendices in main.tex. Uses the macros that
% sections/security.tex provides.

\section{Các chứng minh an toàn}
\label{app:security-proofs}

\subsection{Quy ước}

Các oracle ngẫu nhiên được lấy mẫu lười (lazily). Một điểm của $\Hfs$ được
\emph{xác định} một khi nó đã được truy vấn hoặc lập trình; một truy vấn tại
một điểm chưa xác định là \emph{mới} (fresh). $Q_H$ chặn mọi truy vấn oracle
ngẫu nhiên của đối thủ, bao gồm các truy vấn của bên kiểm chứng trong
$\mathcal O_{\mathsf{Ver}}$ và một lần kiểm chứng cuối, còn $Q_P$ chặn số chứng
minh trung thực. Mọi phép rút gọn dưới đây tự chạy toàn bộ thí nghiệm, điều mà
nó làm được vì việc kiểm chứng là công khai và mọi cam kết trung thực đều do
bên thách thức sinh ra. Đối với oracle mở rộng khóa có một điểm duy nhất (thẻ
khóa theo sau bởi $\pid$); một phép rút gọn cần cấy một ma trận thách thức sẽ
lập trình điểm này trước khi đối thủ bắt đầu, điều này khiến $\mathsf{ck}$ được
phân phối đúng như một khóa đều.

\subsection{Chứng minh Bổ đề~\ref{lem:hvzk}}

Cố định $r\in S_1^k$ và $c\in\Cset$, và đặt $v=cr\in\mathbb Z^{kN}$ là tích
negacyclic nguyên. Vì $\norm{v}_\infty\le\norm{c}_1\norm{r}_\infty
\le\kappa<q/2$, việc tâm hóa tích theo modulo $q$ khôi phục $v$ một cách chính
xác, đúng như những gì phần hiện thực tính (\texttt{opening\_proof.py:79--80}).
Với một tọa độ và bất kỳ $a\in[-B,B]$, $\Pr[y_i+v_i=a]=\Pr[y_i=a-v_i]=1/2\gamma$,
bởi vì $a-v_i\in[-B-\kappa,B+\kappa]=[-\gamma+1,\gamma-1]$ nằm trong dải mặt nạ.
Các tọa độ độc lập, nên $\Pr[z=z_0]=(2\gamma)^{-kN}$ với mọi $z_0\in Z$. Lấy
tổng trên $Z$ cho xác suất chấp nhận $|Z|(2\gamma)^{-kN}=p_{\mathrm{acc}}$, và
với điều kiện được chấp nhận thì $z$ đều trên $Z$; không đại lượng nào phụ thuộc
vào $r$ hay $c$. \qed

\subsection{Chứng minh Bổ đề~\ref{lem:sim}}

\emph{Bên chứng minh thực} (\texttt{opening\_proof.py:86--109}). Trên $(x,t,r)$:
rút $e\leftarrow\{0,1\}^{256}$; $K=H_{\mathsf{hg}}(e\|x,t,r\text{ có tiền tố độ dài})$;
với $i=1,\dots,\Amax$: $y_i=H_{\mathsf y}(K\|i)\in Y$, $w_i=B_1y_i$,
$h_i=\Hfs(\pid,x,t,w_i)$, $c_i=G(h_i)$, $z_i=y_i+c_ir$; xuất ra $(h_i,z_i)$ tại
$i$ đầu tiên có $z_i\in Z$, hoặc $\bot$ sau $\Amax$ lần loại bỏ.

\emph{Bộ mô phỏng} $\mathcal S(x,t)$. Xuất ra $\bot$ với xác suất
$(1-p_{\mathrm{acc}})^{\Amax}$. Ngược lại rút $h\leftarrow\{0,1\}^{256}$ và
$z\leftarrow Z$, đặt $w=B_1z-G(h)t_1$; nếu $\Hfs$ đã xác định tại
$(\pid,x,t,w)$ thì loại bỏ; ngược lại lập trình $\Hfs(\pid,x,t,w):=h$ và xuất ra
$(h,z)$. Bộ mô phỏng không bao giờ dùng $r$.

Chúng tôi chặn độ thay đổi trên mỗi chứng minh; một cận hợp (union bound) trên
$Q_P$ chứng minh cho ta bổ đề. Viết $\eta_y=nN\log_2(2\gamma)$. Bởi vì
$B_1=[I_n|B_1']$, $n$ hàng đầu của $w=B_1y$ là $y_{[1..n]}+B_1'y_{[n+1..k]}$,
một ảnh song ánh của $y_{[1..n]}$ với $y_{[n+1..k]}$ cố định; do đó
$H_\infty(w)\ge\eta_y$ cho một mặt nạ đều,
$H_\infty(w)\ge\eta_y-\log_2\frac1{1-p_{\mathrm{acc}}}\ge\eta$ cho một mặt nạ
với điều kiện bị loại bỏ, và $H_\infty(w)\ge\eta$ cho một $w=B_1z-ct_1$ được
mô phỏng với $z$ đều trên $Z$.

\begin{itemize}
\item[$G_0$:] bên chứng minh thực.
\item[$G_1$:] $K$ được thay bằng một giá trị đều. Hai trò chơi chỉ khác nhau
  nếu một truy vấn nào đó tới $H_{\mathsf{hg}}$ khác với truy vấn của chính bên
  chứng minh bắt đầu bằng $e$: tối đa $(Q_H+Q_P)2^{-256}$.
\item[$G_2$:] mỗi $y_i$ được thay bằng một phần tử đều mới của $Y$. Hai trò
  chơi chỉ khác nhau nếu $H_{\mathsf y}$ bị truy vấn ở nơi khác tại $K\|i$: tối
  đa $(Q_H+Q_P)2^{-256}$.
\item[$G_3$:] trong một lần thử bị loại, bên chứng minh rút $h_i$ một cách đều
  và đặt $c_i=G(h_i)$ mà không truy vấn hay xác định $\Hfs$. Hai trò chơi chỉ
  khác nhau nếu điểm $(\pid,x,t,w_i)$, tại bất kỳ thời điểm nào, bị $\mathcal A$
  truy vấn hoặc được một chứng minh trung thực khác xác định: tối đa
  $\Amax(Q_H+\Amax Q_P)2^{-\eta}$. Đây là bước liên quan tới các truy vấn
  oracle ngẫu nhiên ẩn của các lần thử bị loại~\cite{ref:barbosa,ref:devevey}.
\item[$G_4$:] điểm của lần thử được chấp nhận được xác định bởi truy vấn của
  bên chứng minh; nó là mới ngoại trừ với xác suất $(Q_H+(\Amax+1)Q_P)2^{-\eta}$.
  Với điều kiện là mới, $h_i$ đều trong mọi lần thử và, theo
  Bổ đề~\ref{lem:hvzk}, lần thử $i$ được chấp nhận với xác suất đúng bằng
  $p_{\mathrm{acc}}$ bất kể $h_i$, $c_i$ và $r$. Vậy chỉ số của lần chấp nhận
  đầu tiên là hình học bị cắt cụt, và với điều kiện được chấp nhận, $h$ đều và
  $z$ đều trên $Z$, độc lập với nhau và với $r$. Các lần thử bị loại không còn
  liên kết với oracle nào sau $G_3$ và có thể được bỏ đi. Khi đó đầu ra và định
  nghĩa oracle $\Hfs(\pid,x,t,B_1z-G(h)t_1)=h$ được phân phối đúng như trong
  $\mathcal S$, mà sự kiện loại bỏ của nó chính là sự kiện không-mới vừa được
  chặn.
\end{itemize}
Tổng trên mỗi chứng minh là $2(Q_H+Q_P)2^{-256}+(\Amax+1)(Q_H+(\Amax+1)Q_P)2^{-\eta}$.

Đối với $\Pi_{\mathsf{eq}}$ (\texttt{equality\_proof.py:58--104}), bên chứng
minh rút hai mặt nạ mỗi lần thử, thông báo $(w_a,w_b,u)=(B_1y_a,B_1y_b,\langle b_2,y_a-y_b\rangle)$
và chấp nhận khi và chỉ khi cả hai đáp ứng đều nằm trong $Z$. Theo
Bổ đề~\ref{lem:hvzk} áp dụng cho mỗi đáp ứng, sự chấp nhận có xác suất đúng bằng
$p_{\mathrm{acc}}^2$ và $(z_a,z_b)$ được chấp nhận đều trên $Z^2$. Bộ mô phỏng
lấy mẫu $(h,z_a,z_b)$, đặt $w_a=B_1z_a-ct_{1,a}$, $w_b=B_1z_b-ct_{1,b}$,
$u=\langle b_2,z_a-z_b\rangle-c(t_{2,a}-t_{2,b})$ và lập trình $H_{\mathsf{eq}}$;
thông báo có min-entropy ít nhất $\eta$ thông qua $w_a$. Cùng các bước lai cho
cùng cận. \qed

\subsection{Một bổ đề trích xuất bằng rẽ nhánh}

\begin{lemma}[Trích xuất]
\label{lem:fork}
Gọi $\mathcal B$ là một thuật toán với các đồng xu $\rho$ thực hiện nhiều nhất
$Q$ truy vấn mới tới $\Hfs$, được trả lời bởi $h_1,\dots,h_Q\leftarrow\{0,1\}^{256}$
(oracle $G$ và mọi điểm được $\mathcal B$ lập trình đều do $\rho$ xác định).
$\mathcal B$ xuất ra $\bot$ hoặc $(J,x^*,t^*,\pi^*)$ sao cho
$\pi^*=(h_J,z^*)$ được chấp nhận cho $(x^*,t^*)$ và truy vấn mới thứ $J$ là
$(\pid,x^*,t^*,w^*)$ với $w^*=B_1z^*-G(h_J)t_1^*$. Đặt
$\mathrm{acc}=\Pr[\mathcal B\neq\bot]$. Tồn tại $\mathcal E$, chạy
$\mathcal B$ hai lần, xuất ra $x^*$ của lần chạy thứ nhất và một phần mở nới
lỏng của $t^*$ với xác suất $\varepsilon_{\mathsf{ext}}\ge\mathrm{acc}(\mathrm{acc}/Q-2^{-256})
-2^{-257}-Q2^{-256}$; tương đương
\[
\mathrm{acc}\le Q2^{-256}+\sqrt{Q\bigl(\varepsilon_{\mathsf{ext}}+2^{-257}+Q2^{-256}\bigr)}.
\]
\end{lemma}
\begin{proof}
$\mathcal E$ là thuật toán rẽ nhánh của Bellare và Neven~\cite{ref:bn06}
(Bổ đề~1, bổ đề rẽ nhánh tổng quát; xem thêm~\cite{ref:ps00}): chạy
$\mathcal B(\rho;h_1..h_Q)$; nếu nó trả về $J$, chạy lại nó trên
$(\rho;h_1..h_{J-1},h'_J..h'_Q)$ với $h'_J,\dots,h'_Q$ mới; thành công nếu lần
chạy thứ hai trả về cùng $J$ và $h'_J\neq h_J$. Bổ đề đó cho xác suất thành công
ít nhất $\mathrm{acc}(\mathrm{acc}/Q-2^{-256})$. Hai lần chạy giống hệt nhau cho
tới truy vấn mới thứ $J$, nên chúng chia sẻ $(x^*,t^*,w^*)$, và trả về $(h_J,z)$,
$(h'_J,z')$ với $B_1z-ct_1^*=w^*=B_1z'-c't_1^*$ cho $c=G(h_J)$, $c'=G(h'_J)$.
Giá trị $h'_J$ đều và độc lập với lần chạy thứ nhất; trừ khi nó là một trong
nhiều nhất $Q$ đầu vào mà tại đó $G$ đã được tính trong lần chạy thứ nhất
(xác suất $\le Q2^{-256}$), $G(h'_J)$ là một phần tử đều mới của $\Cset$ và
bằng $c$ với xác suất $1/|\Cset|\le2^{-257}$. Ngược lại
$f=c-c'\in\Cdiff$ và $\bar r=z-z'$ thỏa $\norm{\bar r}_\infty\le2B$ và
$B_1\bar r=ft_1^*$. Dạng thứ hai suy ra bằng cách giải bất đẳng thức bậc hai:
$\mathrm{acc}^2-Q2^{-256}\mathrm{acc}-Q\delta\le0$ với
$\delta=\varepsilon_{\mathsf{ext}}+2^{-257}+Q2^{-256}$ kéo theo
$\mathrm{acc}\le Q2^{-256}+\sqrt{Q\delta}$.
\end{proof}

\subsection{Chứng minh Định lý~\ref{thm:se}}

$G_0$ là $\mathsf{SE}$. Trong $G_1$, mọi chứng minh trung thực đều do bộ mô
phỏng tạo ra, nên $|\Pr[W_0]-\Pr[W_1]|\le\epszk(Q_P,Q_H)$, trong đó $W_i$ là sự
kiện thắng. Gọi $\mathsf{Prog}$ là sự kiện $\mathcal A$ thắng trong $G_1$ và
điểm $(\pid,x^*,t^*,w^*)$ của bên kiểm chứng đã được bộ mô phỏng lập trình, cho
một bộ ba $(x_i,t_i,\pi_i)\in L$. Khi đó $x^*=x_i$, $t^*=t_i$, $w^*=w_i$, và
sự chấp nhận buộc $h^*=h_i$, nên $c^*=c_i$. Vì $(x^*,t^*,\pi^*)\notin L$, ta
có $z^*\neq z_i$, và $B_1z^*-c_it_1^*=B_1z_i-c_it_1^*$ cho
$B_1(z^*-z_i)=0$ với $0<\norm{z^*-z_i}_\infty\le2B\le\beta$. Một đối thủ M-SIS
cấy ma trận của nó làm $B_1'$, tự rút $b_2'$, chạy $G_1$ và xuất ra $z^*-z_i$
thành công mỗi khi $\mathsf{Prog}$ xảy ra, nên
$\Pr[\mathsf{Prog}]\le\Adv^{\MSIS}$. Ngược lại, điểm của bên kiểm chứng được
xác định bởi một truy vấn mới, và $\mathcal B$ (chạy $G_1$ và xuất ra chỉ số
của truy vấn đó) thỏa giả thiết của Bổ đề~\ref{lem:fork} với
$Q=Q_H$ và $\mathrm{acc}\ge\varepsilon-\epszk-\Adv^{\MSIS}$. Thay vào cho ta
định lý.

Truy vấn bị rẽ nhánh chứa $x^*$, thứ mà bên kiểm chứng dựng từ các mục ghim của
nó và từ $\mathsf{Read}(\vid^*)$ (\texttt{video\_pipeline.py:278--284}); nó
không bao giờ lấy một trường phát biểu nào từ bên chứng minh ngoại trừ $N_s^*$,
thứ mà nó kiểm tra đối chiếu với kho của nó (hoặc với giá trị cố định) và bản
thân là một phần của $x^*$. Do đó phần mở trích xuất được được ràng buộc vào các
giá trị của bên kiểm chứng cho $\mathit{rid}$, $\pid$, gốc sổ đăng ký và epoch,
$N_s$, chính sách codec, hồ sơ vật mang, $|P|$, $\mathrm{SHA256}(P)$ và $V$. Hai
phát biểu khác nhau ở bất kỳ trường nào đều khác nhau khi là chuỗi, vì mọi
trường đều có độ rộng cố định (\texttt{video\_proof.py:68--77}). \qed

\subsection{Chứng minh Định lý~\ref{thm:binding}}

Cho $(\bar r,f,m)$ và $(\bar r',f',m')$ mở $t$ với $m\neq m'$, trong đó một
phần mở trung thực có $f=1$ và $\norm{\bar r}_\infty\le1$. Từ $ft_1=B_1\bar r$
và $f't_1=B_1\bar r'$ ta được $B_1x=0$ cho $x=f'\bar r-f\bar r'$. Từ
$ft_2=\langle b_2,\bar r\rangle+fm$ và $f't_2=\langle b_2,\bar r'\rangle+f'm'$,
\[
\langle b_2,x\rangle=f'(ft_2-fm)-f(f't_2-f'm')=ff'(m'-m).
\]
Nếu $x\equiv0\pmod q$ thì $ff'(m'-m)=0$ và, vì $f$ và $f'$ khả nghịch
(Bổ đề~\ref{lem:invertible}; $1$ hiển nhiên), $m=m'$, mâu thuẫn. Chuẩn
của $x$ tối đa là $\norm{f'}_1\norm{\bar r}_\infty+\norm{f}_1\norm{\bar r'}_\infty$,
tức $2\kappa\cdot2B+2\kappa\cdot2B=8\kappa B$ cho hai phần mở nới lỏng,
$2\kappa\cdot1+1\cdot2B$ cho một nới lỏng và một trung thực, và $2$ cho hai phần
mở trung thực. Tất cả đều tối đa $\beta<q/2$ (hàm khởi tạo tham số bắt buộc
$2\beta<q$, \texttt{params.py:53--54}), nên $x$ là một vector ngắn khác không
cả theo modulo $q$ lẫn trên $\mathbb Z$. Phép rút gọn cấy ma trận của nó làm
$B_1'$, lấy mẫu $b_2'$ và xuất ra $x$. \qed

\subsection{Proof of Lemma~\ref{lem:row}}

\begin{lemma}[Hiding instance]
\label{lem:mlwe-instance}
There is an efficient map that, for any $m$ chosen after seeing
$\mathsf{ck}$, sends an M-LWE instance $(A,u)$, $A\in\Rq^{(n+1)\times\lambda}$,
to $(\mathsf{ck},t)$ with $\mathsf{ck}$ uniform, such that $t=\Com(m;r)$ with
$r\leftarrow S_1^k$ if $u=As+e$, and $t$ is uniform if $u$ is uniform.
\end{lemma}
\begin{proof}
Write $A=\bigl[\begin{smallmatrix}A_1\\a_2\end{smallmatrix}\bigr]$ and
$u=(u_1,u_2)$. Draw $B_{1,2}'\leftarrow\Rq^{n\times1}$ and set $b_2'=a_2$,
$B_{1,3}'=A_1+B_{1,2}'a_2$, which makes $\mathsf{ck}$ uniform. Given $m$, set
$t_2=u_2+m$ and $t_1=u_1+B_{1,2}'u_2$. If $u=As+e$ with $e=(e_1,e_2)$ and
$r=(e_1,e_2,s)$, then $\langle b_2,r\rangle=e_2+a_2s=u_2$ and
$B_1r=e_1+B_{1,2}'e_2+(A_1+B_{1,2}'a_2)s=u_1+B_{1,2}'u_2$, so $t=\Com(m;r)$
with $r$ uniform ternary. If $u$ is uniform, so is $t$, because
$(u_1,u_2)\mapsto(u_1+B_{1,2}'u_2,u_2+m)$ is a bijection.
\end{proof}

Now let $\mathcal A$ output $(j,\bar r,f)$ with $f\in\Cdiff$,
$\norm{\bar r}_\infty\le2B$ and $ft^{(j)}_1=B_1\bar r$, and put $x=fr_j-\bar r$.
Then $B_1x=ft_1^{(j)}-ft_1^{(j)}=0$ and $\norm{x}_\infty\le2\kappa+2B<q/2$.

\emph{Case $x\neq0$.} An M-SIS adversary plants its matrix as $B_1'$, samples
$b_2'$, generates every $t^{(j)}$ itself (so it knows every $r_j$), and outputs
$x$. Hence $\Pr[x\neq0]\le\Adv^{\MSIS}$.

\emph{Case $x=0$.} Then $\bar r\equiv fr_j$, so $r_j$ is the centered
representative of $f^{-1}\bar r\bmod q$. An M-LWE distinguisher guesses
$j_0\leftarrow[Q]$, builds $\mathsf{ck}$ and $t^{(j_0)}$ from its instance by
Lemma~\ref{lem:mlwe-instance} (with the message $m_{j_0}$ that $\mathcal A$
chooses), generates the other commitments honestly, and outputs $1$ iff
$\mathcal A$ returns $j=j_0$ and $r'=\mathrm{center}(f^{-1}\bar r)$ is ternary
with $\Com(m_{j_0};r')=t^{(j_0)}$. On a real instance it outputs $1$ with
probability at least $\Pr[x=0]/Q$. On a uniform instance $t^{(j_0)}$ is uniform
and independent of $m_{j_0}$, so a ternary opening exists with probability at
most $|S_1^k|/q^{(n+1)N}=3^{kN}q^{-(n+1)N}$. Hence
$\Pr[x=0]\le Q(\Adv^{\MLWE}+3^{kN}q^{-(n+1)N})$. \qed

\subsection{Proof of Theorem~\ref{thm:hiding}}

$H_0$ is $\mathsf{IND}^0$. In $H_1$ all honest proofs, including the challenge
proof, are simulated: $|H_0-H_1|\le\epszk(Q_P+1,Q_H)$. In $H_2$ the challenge
commitment is replaced by a uniform pair; a distinguisher between $H_1$ and
$H_2$ is an M-LWE distinguisher by Lemma~\ref{lem:mlwe-instance} with
$m=\mathsf{enc}(p_0)$: it builds $\mathsf{ck}$ from its instance, generates the
other commitments honestly, answers $\mathcal O_{\mathsf{Disc}}$ and
$\mathcal O_{\mathsf{Link}}$ for them, and the challenge randomness is never
needed because the challenge proof is simulated and disclosure and linkage of
the challenge are excluded. So $|H_1-H_2|\le\Adv^{\MLWE}$. The stego video in
$H_2$ is $\mathsf{Write}(\vid,P,E)$ with $P$, $D$, $x$ independent of the
payload (fixed from a placeholder of the final length,
\texttt{video\_pipeline.py:216--227}), and $E$ contains only $\pid$, $N_s$, a
uniform $t$ and a simulated proof; $H_2$ is therefore independent of $b$. The
same three hybrids from $\mathsf{IND}^1$ end in the same $H_2$, which gives
the bound. The payload length is hidden because $\mathsf{enc}$ maps every
payload of at most 766 bytes to one polynomial and the envelope has a fixed
size. \qed

\subsection{Proof of Theorem~\ref{thm:replay}}

We model the store as a linearizable table keyed by
$d(N)=\mathrm{SHA256}(\text{tag}\|N)$ (\texttt{zkp\_sessions.py:113--117}).
Rows are inserted only by $\mathsf{Issue}$ with a fresh uniform $N$
(\texttt{:132--151}) and deleted only once expired (\texttt{:140--143}).
Acceptance requires the preflight \texttt{is\_active} and then
\texttt{consume}, an \texttt{UPDATE} with the conditions
\texttt{consumed\_at IS NULL}, matching binding and \texttt{expires\_at>now},
committed inside \texttt{BEGIN IMMEDIATE}, returning success iff exactly one
row changed (\texttt{:181--205}).

(i) After the first successful \texttt{consume} of a row, the condition
\texttt{consumed\_at IS NULL} fails for that row forever; a second success for
the same $N_s$ needs another row with key $d(N_s)$, i.e.\ a later issuance of
some $N'$ with $d(N')=d(N_s)$: either $N'=N_s$ (probability at most
$Q_I^22^{-257}$ over all pairs of issuances) or a SHA-256 collision.
(ii) Any row with key $d(N_s^*)$ comes from an issuance of some $N$ with
$d(N)=d(N_s^*)$, so $N_s^*$ was issued unless SHA-256 collides; the binding
and expiry are part of the \texttt{WHERE} clause, evaluated at consumption
time. (iii) The session field occupies bytes $106$--$137$ (0-indexed) of every statement,
and $x^*$ is a substring of the point $(\pid,x^*,t^*,w^*)$. A query made
before the issuance of $N_s^*$ fixes this field before $N_s^*$ is drawn, so it
matches with probability $2^{-256}$; a union bound over $Q_H$ queries and $Q_I$
sessions gives $Q_HQ_I2^{-256}$. \qed

\subsection{Proof of Theorem~\ref{thm:auth}}

Let $W_a$, $W_b$ be the two winning cases.

$W_a$: $\mathsf{Pins}$ verifies $\mathsf{Reg}^*$ under $\mathit{pk}_I$
(\texttt{video\_pipeline.py:146--147}). If its signing bytes were never signed
through $\mathcal O_{\mathsf{Reg}}$, it is an EUF-CMA forgery; the reduction
runs the whole experiment, using its signing oracle for
$\mathcal O_{\mathsf{Reg}}$. Otherwise $\mathsf{Reg}^*$ is a signed registry
and contains a descriptor equal to the registered descriptor of $t^*$
(\texttt{:150--153}). That descriptor has the registered constraint-module
hash, so it is not in any $\Delta$ (in-band descriptors have a different
module hash, unless SHA-256 collides on the two fixed module strings), and it
was computed from some $t'\in T_{\mathsf{reg}}$. Since $t^*\notin T_{\mathsf{reg}}$,
$t'\neq t^*$ and the canonical encodings differ, so
$\mathrm{SHA256}(B_1'\|b_2'\|t^*)=\mathrm{SHA256}(B_1'\|b_2'\|t')$
(\texttt{:119--128}) is a collision. Hence
$\Pr[W_a]\le\Adv^{\mathsf{euf\text{-}cma}}_{\mathsf{ML\text{-}DSA}}+\Adv^{\mathsf{cr}}_{\mathsf{SHA2}}$.

$W_b$: simulate all honest proofs ($\epszk$) and handle the programmed-point
case as in Theorem~\ref{thm:se} ($\Adv^{\MSIS}$). The remaining runner
$\mathcal B$ receives the honest credentials as a $Q_C$-instance challenge of
Lemma~\ref{lem:row}, never uses their randomness (all proofs are simulated,
and a disclosure query for an instance aborts the run, which does not affect
runs in which $\mathcal A$ wins with an undisclosed $j$), and outputs the
index of the verifier's fresh query. Lemma~\ref{lem:fork} extracts a relaxed
opening of $t^*=t_c^{(j)}$, which is a Lemma~\ref{lem:row} adversary, so
$\varepsilon_{\mathsf{ext}}\le\Adv^{\mathsf{row}}_{Q_C}$. Adding the cases
gives the theorem.

A tight alternative for $W_b$ follows the ML-DSA analysis: after replacing the
credential by a uniform pair (M-LWE), an accepted proof is a solution of the
SelfTargetMSIS problem for $[B_1\,|\,t_1]$ and $\Hfs$~\cite{ref:kls18}, without
rewinding. We state the rewinding bound because it uses only the estimated
M-SIS and M-LWE instances. \qed

\subsection{Proofs for linkage (Theorem~\ref{thm:link})}

\begin{lemma}[Special soundness of the equality proof]
\label{lem:eq-ss}
Two accepting transcripts $((w_a,w_b,u),c,z_a,z_b)$ and
$((w_a,w_b,u),c',z_a',z_b')$ with $c\neq c'$ yield relaxed openings
$(\bar r_a,f)$ of $t_a$ and $(\bar r_b,f)$ of $t_b$, $f=c-c'$, with equal
messages.
\end{lemma}
\begin{proof}
The verifier checks $B_1z_a=w_a+ct_{1,a}$, $B_1z_b=w_b+ct_{1,b}$ and
$\langle b_2,z_a-z_b\rangle=u+c(t_{2,a}-t_{2,b})$
(\texttt{equality\_proof.py:135--145}). Subtracting the two transcripts gives
$B_1\bar r_a=ft_{1,a}$, $B_1\bar r_b=ft_{1,b}$ and
$\langle b_2,\bar r_a-\bar r_b\rangle=f(t_{2,a}-t_{2,b})$ with
$\bar r_\ast=z_\ast-z_\ast'$, $\norm{\bar r_\ast}_\infty\le2B$. Then
$\bar m_a-\bar m_b=(t_{2,a}-t_{2,b})-f^{-1}\langle b_2,\bar r_a-\bar r_b\rangle=0$.
\end{proof}

\emph{Soundness.} Simulate $\mathcal O_{\mathsf{Link}}$ ($\epszk$). If the
verifier's $H_{\mathsf{eq}}$ point was programmed for the same context and
commitments, then $c$ is the programmed one and $(z_a,z_b)\neq(z_a^{(i)},z_b^{(i)})$
with $B_1(z_a-z_a^{(i)})=B_1(z_b-z_b^{(i)})=0$, an M-SIS solution of norm $2B$.
Otherwise Lemma~\ref{lem:fork}, with the forked oracle $H_{\mathsf{eq}}$ and
Lemma~\ref{lem:eq-ss} in place of the opening relation, extracts openings with
a common message $\bar m$. The adversary's first-run openings have messages
$m_a\neq m_b$, so $m_a\neq\bar m$ or $m_b\neq\bar m$, and
Theorem~\ref{thm:binding} applied to $t_a$ or $t_b$ gives an M-SIS solution.
Hence $\varepsilon_{\mathsf{ext}}\le\Adv^{\MSIS}$ in the bound of
Theorem~\ref{thm:se}.

\emph{Hiding.} Simulate the proof ($\epszk(Q_P+1,Q_H)$), then replace $t_a$ and
$t_b$ by uniform pairs, one M-LWE hybrid each (Lemma~\ref{lem:mlwe-instance};
the randomness is independent). The final game is independent of $b$; doing
the same from $b=1$ gives $2\epszk+4\Adv^{\MLWE}$.

\emph{Disclosure by proof.} The equality proof with $t_b=(0,\mathsf{enc}(p))$
and $r_b=0$ (a ternary opening of $t_b$) proves that $t_a$ hides $p$ without
revealing $r_a$: extraction gives a relaxed opening $(\bar r_b,f)$ of $t_b$
with $B_1\bar r_b=0$, so $\bar r_b=0$ unless M-SIS is broken, and then
$\bar m_b=\mathsf{enc}(p)=\bar m_a$. This is the fix proposed for item~D3.
\qed

\subsection{Proof of Lemma~\ref{lem:canon} and Theorem~\ref{thm:integrity}}

\emph{Unique decodability.} For a carrier segment the hashed string is the
domain tag followed, for each NAL unit, by \texttt{"NAL\textbackslash0"} and a
fixed 12-byte header, then either \texttt{"RAW\textbackslash0"}, a 64-bit
length and the RBSP, or a sequence of \texttt{"BITS\textbackslash0"}, a 64-bit
bit count and the packed bits, alternating with \texttt{"BLOCK\textbackslash0"},
the 16-byte key and length, and the signed 64-bit levels
(\texttt{video\_canonicalization.py:234--307}). The four tags differ in their
first two bytes, every field is fixed-width or length-prefixed, and the
padding of packed bits is determined by the bit count, so a left-to-right
parser recovers the whole sequence. The segment combiner hashes a tag, the
32-bit frames-per-segment and segment count, and fixed-width
$(i,f_i,d_i)$ triples (\texttt{blind\_sync\_streaming.py:225--233}); $V$ hashes a
tag and $D$. A difference in $\Canon$ is therefore a difference in the input
of one of the three SHA-256 layers whose outputs are equal, which is a
collision. Carrier segments are exactly those containing positions of $P$, so
for equal $P$ both sides hash each segment in the same mode.

\emph{What is determined.} Fix $\Canon$ and $P$. Segments without carriers
are fixed byte for byte. In a carrier segment, every NAL header field in the
serialization, every RBSP of a NAL unit without carriers, and every RBSP bit
outside carrier blocks up to the stop bit are fixed. Each carrier block is
fixed as a level vector except at carriers, where only the normalized value
$\nu=\mathrm{sgn}(\ell)2\lfloor|\ell|/2\rfloor$ is fixed. The pipeline enforces
the stable rule (\texttt{blind\_sync\_streaming.py:268--271}; carriers have
$|\ell|\ge4$ and the embedder flips the magnitude parity), so
$\ell=\nu+\mathrm{sgn}(\nu)b$ with $b$ the carrier bit; for the optional
trailing-one sign carrier, $\ell=\pm1$ with the sign as the bit. The carrier
bits are written in the order of $P$ as the chunk-framed encoding of $E$.
If that encoding is unique, $E$ fixes the bits and hence every carrier level.

\emph{Theorem~\ref{thm:integrity}.} Let $(x_h,t,\pi_h)\in L$ be the honest
triple. If $(x^*,t,\pi^*)\notin L$, the event is bounded as in
Theorem~\ref{thm:auth}, case $W_b$, with the $Q_P$ honest commitments as the
Lemma~\ref{lem:row} instances. Otherwise $x^*=x_h$ and $\pi^*=\pi_h$ (each
honest in-band commitment occurs in one triple except with negligible
probability; a credential occurs in several, all honest). Equal statements
give equal $|P|$ and $\mathrm{SHA256}(P)$, hence $P^*=P_h$ unless SHA-256
collides (the positions are hashed as canonical JSON of integer triples), and
equal $V$, hence equal $\Canon$ by the first part. The envelope is
$E^*=\texttt{magic}\|1\|\pid\|N_s\|t\|\pi^*$, and its parser accepts only this
exact layout (\texttt{video\_proof.py:142--153, 218--228}), so $E^*=E_h$. \qed

\subsection{Completeness}

An accepted attempt satisfies $B_1z-ct_1=B_1y+cB_1r-cB_1r=w$, and for the
equality proof $\langle b_2,z_a-z_b\rangle-c(t_{2,a}-t_{2,b})=\langle b_2,y_a-y_b\rangle$
because $t_{2,a}-t_{2,b}=\langle b_2,r_a-r_b\rangle$ for equal messages. By
Lemma~\ref{lem:hvzk} the number of attempts is geometric with parameter
$p_{\mathrm{acc}}$ (mean $5.23$) for the opening proof and $p_{\mathrm{acc}}^2$
(mean $27.4$) for the equality proof. With the caps $\Amax=512$
(\texttt{opening\_proof.py:33}) and $\Amax=4\,096$
(\texttt{equality\_proof.py:50}) the prover fails with probability
$2^{-156.7}$ and $2^{-220.0}$ respectively. Under the shared cap of 512 that
the reviewed code used, the equality proof failed with probability
$2^{-35.0}$ for \texttt{bdlop-open-128-v1} (means $4.65$ and $21.6$); this is
item~D4.

\section{Specification and implementation cross-check}
\label{app:code-check}

Each item lists the property relied on above and where the implementation
realizes it, at commit \texttt{ea96d60}. Paths are relative to
\texttt{VideoLevel/src}; \texttt{opening\_proof.py}, \texttt{equality\_proof.py},
\texttt{bdlop.py}, \texttt{params.py}, \texttt{ring.py}, \texttt{codec.py},
\texttt{sampling.py}, \texttt{video\_proof.py}, \texttt{video\_pipeline.py}
and \texttt{cli.py} are in \texttt{lattice\_zk/}.

\begin{enumerate}
\item \emph{Statement.} Field order and widths: \texttt{u16} version,
  $\mathit{rid}$, $\pid$, root, \texttt{u64} epoch, $N_s$, codec policy,
  carrier profile, \texttt{u64} carrier count, positions hash, $V$; every
  digest checked to be 32 bytes (\texttt{video\_proof.py:58--77}).
  Matches the specification, \S3.6. \textbf{Consistent.}
\item \emph{Fiat--Shamir input.} SHAKE-256 over the tag, $\pid$,
  \texttt{u64}$(|x|)\|x$, the full commitment (6\,720 bytes) and $w$ as
  big-endian 64-bit words; 32-byte output
  (\texttt{opening\_proof.py:42--50}). Only $x$ is length-prefixed; the other
  fields have lengths fixed by $\pid$. Equality proof: same with both
  commitments and $(w_a,w_b,u)$ (\texttt{equality\_proof.py:40--51}).
  \textbf{Consistent}; the specification's ``with length prefixes'' (\S4.7)
  was imprecise and has been corrected.
\item \emph{Hedging.} $K=\mathrm{SHA3}(e\|\text{length-prefixed }x,t,r)$
  (\texttt{opening\_proof.py:53--63,99--100}). \textbf{Consistent} (see D7).
\item \emph{Samplers.} Masks: exact uniform on $[-\gamma+1,\gamma]$ (18-bit
  draws below $2\gamma=2^{18}$, \texttt{sampling.py:60--66}); challenges:
  sign bits then Fisher--Yates with exact rejection, as in ML-DSA
  \textsf{SampleInBall} (\texttt{:69--81}); ternary randomness by rejection on
  2-bit draws (\texttt{:54--57}). \textbf{Consistent} with the uniformity used in
  Lemmas~\ref{lem:hvzk} and~\ref{lem:fork}.
\item \emph{Verification.} Exact proof length, packed responses below
  $2B+1$, zero padding, recomputed $w'$, constant-time seed comparison
  (\texttt{opening\_proof.py:118--138}). The packing has no padding bits at
  these parameters, so the proof encoding is canonical. \textbf{Consistent.}
\item \emph{Parameter checks.} Prime $q\equiv5\pmod8$, $2\sqrt2<\sqrt q$,
  $2\cdot8\kappa B<q$ (\texttt{params.py:37--54}, \texttt{ring.py:52--58});
  production pins require $|\Cset|\ge2^{256}$
  (\texttt{params.py:127--130}, \texttt{video\_pipeline.py:144--145}).
  \textbf{Consistent.}
\item \emph{Envelope.} In-band: exact length, magic \texttt{ZKLB}, version,
  pinned $\pid$, canonical commitment ($<q$ per coefficient), proof
  (\texttt{video\_proof.py:142--153}); the header session must equal the
  statement's (\texttt{:183}). Registered: \texttt{ZKLR}, header and proof only,
  pinned commitment (\texttt{:218--244}). \textbf{Consistent} with \S3.7.
\item \emph{Sessions.} One-use semantics as in the proof of
  Theorem~\ref{thm:replay}; binding $\mathsf{bind}(\mathsf{pins})$ hashes
  $\pid$, $\mathit{rid}$, root, epoch, codec policy and profile
  (\texttt{video\_pipeline.py:166--172}), so in registered mode it also
  depends on the credential through $\mathit{rid}$. Stateless mode requires
  an explicit \texttt{expected\_session} (\texttt{:262--263}).
  \textbf{Consistent.}
\item \emph{Registry.} ML-DSA-65 over a tagged canonical JSON encoding; the
  registered descriptor hashes the key and the credential
  (\texttt{video\_pipeline.py:109--131}, \texttt{zkp\_registry.py:128--147}).
  \textbf{Consistent}, but see D2.
\item \emph{Canonical digest.} As described in
  Section~\ref{sec:sec-integrity}; the specification (\S3.6) described $D$ as
  a single-stream digest and has been corrected to the segmented form.
\item \emph{Domain tags.} SHAKE-256 tags \texttt{.../bdlop-open/v1/\{key,
  randomness, fiat-shamir, mask\}\textbackslash0},
  \texttt{.../bdlop-equal/v1/\{fiat-shamir, mask\}\textbackslash0} and
  \texttt{.../challenge-poly/v1\textbackslash0} each end in the only NUL byte
  and are pairwise distinct, hence prefix-free. \textbf{Consistent.}
\end{enumerate}

\subsection{Discrepancies and gaps}

\begin{description}
\item[D1 (medium): carrier bits are malleable with a single chunk.]
  \texttt{blind\_payload\_chunks.py:12,84--137} accept a one-chunk set for any
  \texttt{chunk\_size} in $[|E|,65\,535]$, and the verifier never compares the
  extracted chunking with the pinned \texttt{profile.chunk\_bytes}
  (\texttt{video\_pipeline.py:} \texttt{267--303}). When the profile's chunk
  size is at least the envelope length, an adversary can rewrite the 16-bit
  field and the outer CRC-32: up to 48 carrier coefficients change by $\pm1$
  while $E$, $P$ and $D$ are unchanged, and the video is accepted. Rewriting
  8\,000 to 20\,000 changed 21 carrier bits in a direct test of the framing
  functions. This contradicts goal (iii) of the threat model (``any change to
  the carriers is rejected''). Fix: the verifier recomputes
  $\mathsf{Enc}_{\mathsf{prof}}(E)$, the chunks of $E$ under the pinned chunk
  size, and requires it to equal the extracted pieces (or requires the
  header field to equal the pinned size); add a negative test.
\item[D2 (medium): no holder identity in the signed registry.] The registered
  descriptor (\texttt{video\_pipeline.py:109--131}; field set
  \texttt{zkp\_registry.py:29--35}) binds a credential commitment, not a
  holder. ``Only the credential holder can produce the proof'' holds per
  commitment; which person that is is not covered by the ML-DSA signature.
  Fix: add a holder identifier (or its digest) to the registered descriptor.
\item[D3 (medium): disclosure hands over the commitment.]
  \texttt{verify\_disclosure} takes the randomness
  (\texttt{video\_proof.py:84--89,188--197}). The recipient can then produce
  fresh opening proofs for that commitment on new statements, i.e.\ accepted
  in-band videos with the creator's commitment, or impersonate a holder if a
  credential is disclosed. Fix: disclose $p$ with the equality proof against
  $(0,\mathsf{enc}(p))$ and zero randomness, bound to a disclosure context
  (proof above).
\item[D4 (low): equality-proof completeness $2^{-35}$.]
  \texttt{equality\_proof.py:29,83} reuses $\Amax=512$ with acceptance
  $p_{\mathrm{acc}}^2$ (figures for \texttt{bdlop-open-128-v1}, the set
  reviewed; $2^{-27.5}$ for \texttt{bdlop-open-128-v2}). Fix: a separate cap,
  e.g.\ 4\,096 attempts ($\approx2^{-280}$ for v1, $2^{-220}$ for v2); a larger cap does not affect zero knowledge because
  acceptance is witness-independent.
\item[D5 (low): scope of the canonical digest.] The digest is over
  ffmpeg-segmented files (\texttt{blind\_sync\_streaming.py:117--215}); segments
  without carriers are hashed as raw files without a domain tag
  (\texttt{:236--241,358,549}); RBSP bits after the stop bit of a carrier-bearing
  IDR NAL unit are excluded (\texttt{video\_canonicalization.py:}
  \texttt{279, 303, 310--317});
  carrier blocks are hashed as level vectors; the codec policy string
  (\texttt{video\_pipeline.py:25,98--99}) does not pin the segmenter version.
  None breaks Theorem~\ref{thm:integrity}, but each widens the set of
  unbound changes listed in Section~\ref{sec:sec-integrity}. Fix: tag the
  raw-segment digests, hash the full RBSP of carrier NAL units, and pin or
  replace the segmenter.
\item[D6 (low): API does not tie pins to the credential.]
  \texttt{verify\_video\_proof} accepts \texttt{pins} and \texttt{credential}
  separately (\texttt{video\_pipeline.py:250--265}) and does not check that
  \texttt{pins.relation\_id} is the relation of that credential. The CLI derives
  both together (\texttt{cli.py:98--107}), but a caller that pins one
  credential and passes another verifies an unregistered credential. Fix:
  recompute the descriptor of \texttt{credential} and compare relation
  identifiers inside \texttt{verify\_video\_proof}.
\item[D7 (info): hedging hash has no domain tag.]
  \texttt{opening\_proof.py:59} hashes $e\|\dots$ directly. Separation from the
  tagged SHA3-256 uses relies on $e$ being uniform; add a tag.
\item[D8 (info): unreduced commitments.] \texttt{Commitment} objects built
  without \texttt{from\_bytes} may hold coefficients in $[q,2^{30})$, which
  \texttt{pack\_uints} accepts (\texttt{codec.py:10--16}). Envelopes are
  unaffected (their decoder checks $<q$); pinned credentials loaded by other
  code paths are. Fix: range check in \texttt{Commitment}.
\item[D9 (info): one global key.] $\mathsf{ck}$ is expanded from a fixed public
  seed (\texttt{params.py:133--149}); all users share one M-SIS/M-LWE
  instance, so precomputation against it amortizes.
\item[D10 (info):] \texttt{verify\_statement\_binding} hard-codes the
  \texttt{lazer-v1} suite (\texttt{zkp\_registry.py:206--210}) and cannot serve
  \texttt{bdlop-open-v1}; the BDLOP path uses \texttt{pins\_from\_registry}.
\item[D11 (info): specification wording.] \S4.5 wrote the bijection as
  $(t_1,t_2)\mapsto(t_1-B_{1,2}'t_2,t_2)$; the correct map uses $t_2-m$
  (Lemma~\ref{lem:mlwe-instance}). Corrected. \S4.2 cites an ``ML-DSA
  specification, Appendix~B'' for the zero-knowledge argument; this reference
  should be checked.
\end{description}

\paragraph{Status after the review.} D1--D4 and D6 are fixed in the code
evaluated in Section~\ref{sec:evaluation}, each with a negative test: the
verifier compares the extracted chunks with the canonical chunking of $E$
under the pinned chunk size (D1); the registered descriptor can bind a
holder-identifier digest (D2); disclosure is the zero-randomness equality
proof of Section~\ref{sec:sec-soundness} and the randomness-revealing API was
removed (D3); the equality proof has its own 4\,096-attempt cap (D4); and
\texttt{verify\_video\_proof} recomputes the relation identifier of the
supplied parameters, profile, credential and holder identifier and refuses
mismatching pins (D6). D5, D7 and D8 and the reviewer question on
production-strength enforcement are closed in a second round, again each
with negative tests (\texttt{src/runtest/test\_hardening\_d5\_d7\_d8.py}).
D7: the hedged mask key hashes a domain tag, the protocol label
(opening, equality or binary proof) and the parameter identifier before the
entropy. Only \texttt{bdlop-open-128-v1} keeps the untagged derivation, so
its known-answer vectors are unchanged; the \texttt{bdlop-open-128-v2} and
binary-proof vectors were regenerated (commitments unchanged). The verifier
never recomputes a mask, so proofs from either derivation verify, which the
tests check. D8: a \texttt{Commitment} carries its parameter set and its
constructor rejects coefficients outside $[0,q)$ and wrong shapes, storing
read-only copies; verifiers still re-check objects built around the
constructor. Production strength: every application-level verifier
(\texttt{verify\_video\_proof}, the envelope, disclosure and linkage
verifiers, and \texttt{pins\_from\_registry}) refuses a parameter set with a
challenge space below $2^{256}$ unless the caller sets
\texttt{allow\_toy\_parameters}; hand-built pins with the toy set are refused.
D5 is closed by a new, opt-in codec policy (v2) with its own identifier, so
it changes profile, policy, relation and statement identifiers, while codec
policy v1 stays byte-identical for existing videos (recorded identifiers and
digests are tested). Under v2 (a) a segment without carriers is hashed with a
domain tag and its length and every segment digest is typed in the combined
digest; (b) a carrier-bearing IDR NAL unit is hashed over its complete RBSP
with its length, so zero bits after the stop bit are bound (under v1 a zero
byte appended there is accepted; under v2 it is rejected, through the real
H.264 transport); (c) the segmentation is pinned: the segments must be
exactly the access-unit partition of the input, compared byte for byte, so
the digest is a function of the input alone and a segmenter that cuts
elsewhere fails closed. A registered \texttt{bdlop-open-128-v2} run on a
200-frame CIF cut under v2 verified from the video alone, and the replayed
session, a one-bit change and the zero-extended carrier NAL unit were
rejected (\texttt{benchmark/results/q2/hardening\_e2e\_container\_cif\_200f\_v2.json}).
The ffmpeg version is not hashed, because that
would bind the verifier's installation rather than the video. Carrier blocks
are still hashed as level vectors, which is what lets the proof bind the
video without depending on its own bits. D9--D11 remain as described.
. Requires amsmath, amssymb, amsthm and the
% theorem environments definition, lemma, theorem defined in main.tex. Extra
% bibliography entries are in sections/security_refs.bib
% (\bibliography{references,sections/security_refs}).
%
% Every code reference is to commit ea96d60 of feature/lattice-pq-stack.

\providecommand{\Rq}{\mathcal{R}_q}
\providecommand{\norm}[1]{\lVert #1 \rVert}
\providecommand{\Adv}{\mathsf{Adv}}
\providecommand{\pid}{\mathit{pid}}
\providecommand{\Cset}{\mathcal{C}}
\providecommand{\Cdiff}{\bar{\mathcal{C}}}
\providecommand{\Com}{\mathsf{Com}}
\providecommand{\vid}{\mathsf{vid}}
\providecommand{\Hfs}{H_{\mathsf{fs}}}
\providecommand{\Stmt}{\mathsf{Stmt}}
\providecommand{\Canon}{\mathsf{Canon}}
\providecommand{\MSIS}{\textsf{M-SIS}}
\providecommand{\MLWE}{\textsf{M-LWE}}
\providecommand{\epszk}{\varepsilon_{\mathsf{zk}}}
\providecommand{\Amax}{A_{\max}}

\section{Mô hình và phân tích an toàn}
\label{sec:security-model}

Mục này phát biểu hệ thống chứng minh được điều gì và dưới những giả thiết nào.
Mọi kết quả đều nằm trong mô hình oracle ngẫu nhiên cổ điển (ROM). Không có định
lý nào trong mô hình oracle ngẫu nhiên lượng tử (QROM) được khẳng định
(Mục~\ref{sec:sec-qrom}). Các chứng minh đầy đủ nằm trong
Phụ lục~\ref{app:security-proofs}; việc kiểm tra chéo từng định nghĩa đối chiếu
với phần hiện thực nằm trong Phụ lục~\ref{app:code-check}.

\subsection{Ký hiệu, tham số và giả thiết}
\label{sec:sec-notation}

Đặt $\mathcal R=\mathbb Z[X]/(X^N+1)$ và $\Rq=\mathcal R/q\mathcal R$, với
$N=256$ và số nguyên tố $q=2^{30}-35\equiv5\pmod 8$. Khóa cam kết là
$\mathsf{ck}=(B_1,b_2)$ với $B_1=[I_n\,|\,B_1']\in\Rq^{n\times k}$,
$b_2=[0^n\,|\,1\,|\,b_2']\in\Rq^{1\times k}$, $k=n+\lambda+1$, $n=6$, $\lambda=7$ (nên $k=14$).
$S_1$ là tập các đa thức có hệ số trong $\{-1,0,1\}$. Với một thông điệp
$m\in\Rq$ và $r\in S_1^k$,
$\Com(m;r)=(t_1,t_2)=(B_1r,\ \langle b_2,r\rangle+m)$. Tập thách thức là
$\Cset=\{c\in\mathcal R:\norm{c}_\infty=1,\ \norm{c}_1=\kappa\}$ với
$\kappa=60$, nên $|\Cset|=\binom{256}{60}2^{60}=2^{257.01}$; tập hiệu của nó
là $\Cdiff=\{c-c':c\neq c'\in\Cset\}$. Các mặt nạ đều (uniform) trên
$Y=[-\gamma+1,\gamma]^{kN}$ với $\gamma=2^{17}$; các đáp ứng được chấp nhận khi
và chỉ khi chúng nằm trong $Z=[-B,B]^{kN}$, $B=\gamma-\kappa-1$. Chúng tôi dùng
$\norm{fv}_\infty\le\norm{f}_1\norm{v}_\infty$ xuyên suốt. Các đại lượng cụ thể
dùng trong các cận được tập hợp trong Bảng~\ref{tab:sec-numbers}.

\begin{table}[t]
\centering
\caption{Các đại lượng cụ thể cho \texttt{bdlop-open-128-v2}}
\label{tab:sec-numbers}
\footnotesize
\begin{tabular}{@{}lr@{}}
\toprule
Quantity & Value \\
\midrule
$|\Cset|$ & $2^{257.01}$\\
Fiat--Shamir seed length & 256 bits\\
$p_{\mathrm{acc}}=((2B+1)/2\gamma)^{kN}$ (opening) & $0.1912$\\
$p_{\mathrm{acc}}^2$ (equality proof) & $0.0365$\\
Attempt cap $\Amax$ (opening; equality) & $512$; $4\,096$\\
$(1-p_{\mathrm{acc}})^{512}$; $(1-p_{\mathrm{acc}}^2)^{4096}$ & $2^{-156.7}$; $2^{-220.0}$\\
$\eta=nN\log_2(2B+1)$ (min-entropy of $w$) & $27\,646.98$ bits\\
$\beta=8\kappa B$ (M-SIS bound) & $62\,885\,280\approx2^{25.91}$\\
$2\kappa+2B$ (one honest, one relaxed opening) & $262\,142\approx2^{18.0}$\\
$\sqrt{q/2}$ (invertibility threshold) & $23\,170.5$\\
$3^{kN}q^{-(n+1)N}$ & $\approx2^{-48\,079}$\\
Statement $x$ & 274 bytes\\
Envelope, in-band / registered & 14\,885 / 8\,165 bytes\\
\bottomrule
\end{tabular}
\end{table}

\begin{lemma}[Hiệu của các thách thức là khả nghịch]
\label{lem:invertible}
Mọi $f\in\Cdiff$ đều khả nghịch trong $\Rq$.
\end{lemma}
\begin{proof}
Vì $q\equiv5\pmod8$, $X^N+1$ tách theo modulo $q$ thành đúng hai nhân tử bất
khả quy, và mọi $y\in\Rq$ khác không với $\norm{y}_\infty<q^{1/2}/\sqrt2$ đều
khả nghịch (Lyubashevsky--Seiler~\cite{ref:ls18}, Hệ quả~1.2 với $k=2$). Mọi
$f\in\Cdiff$ đều khác không với $\norm{f}_\infty\le2<23\,170$. Hàm khởi tạo của
bộ tham số bắt buộc cả hai điều kiện (\texttt{params.py:49--51},
\texttt{ring.py:52--58}).
\end{proof}

\begin{definition}[Các phần mở]
\label{def:relaxed}
Một \emph{phần mở nới lỏng} của $t=(t_1,t_2)$ là một cặp $(\bar r,f)$ với
$f\in\Cdiff$, $\norm{\bar r}_\infty\le2B$ và $ft_1=B_1\bar r$. Nó xác định thông
điệp $\bar m=t_2-f^{-1}\langle b_2,\bar r\rangle$, là $\bar m$ duy nhất thỏa
$ft_2=\langle b_2,\bar r\rangle+f\bar m$ (Bổ đề~\ref{lem:invertible}). Một
\emph{phần mở trung thực} là $(r,1)$ với $r\in S_1^k$ và $t_1=B_1r$, có thông
điệp $t_2-\langle b_2,r\rangle$; đó là thứ mà một lần tiết lộ bộc lộ ra.
\end{definition}

\paragraph{Các giả thiết}
$\Adv^{\MSIS}$ là lợi thế của việc tìm $x\in\mathcal R^k$ với
$B_1x\equiv0\pmod q$, $0<\norm{x}_\infty\le\beta=8\kappa B$, cho $B_1'$ đều.
$\Adv^{\MLWE}$ là lợi thế của việc phân biệt $(A,As+e)$ với đều, cho
$A\in\Rq^{(n+1)\times\lambda}$ đều và $s\in S_1^\lambda$, $e\in S_1^{n+1}$.
Đây là hai thể hiện được ước lượng ở 140/127 và 164/150 bit (core-SVP cổ
điển/lượng tử, Bảng~\ref{tab:params}); ước lượng M-SIS lượng tử chỉ vừa dưới
128 bit. Mọi nghiệm M-SIS trong mục này đều có chuẩn tối đa $\beta$, nên không
cần thêm giả thiết lưới nào nữa.
$\Adv^{\mathsf{cr}}_{\mathsf{SHA2}}$ và $\Adv^{\mathsf{cr}}_{\mathsf{SHA3}}$ là
các lợi thế tìm va chạm đối với SHA-256 và SHA3-256, và
$\Adv^{\mathsf{euf\text{-}cma}}_{\mathsf{ML\text{-}DSA}}$ là lợi thế EUF-CMA đối
với ML-DSA-65~\cite{ref:fips204}.

\paragraph{Các oracle ngẫu nhiên}
Chúng tôi mô hình hóa như các oracle ngẫu nhiên độc lập các thể hiện SHAKE-256
được dùng cho hạt giống thách thức Fiat--Shamir $\Hfs$ (và $H_{\mathsf{eq}}$ cho
chứng minh đẳng thức), phép mở rộng thách thức $G:\{0,1\}^{256}\to\Cset$
(\textsf{SampleInBall}), luồng mặt nạ, phép mở rộng khóa, và hàm băm phòng hộ
(hedging) SHA3-256. Các thẻ miền của chúng là không-tiền-tố đôi một
(Phụ lục~\ref{app:code-check}), nên một oracle SHAKE-256 tách thành các oracle
độc lập. $G$ đều chính xác trên $\Cset$ với một đầu vào mới: các dấu và các chỉ
số Fisher--Yates được rút bằng lấy mẫu loại bỏ chính xác. Phép mở rộng khóa
dưới dạng một oracle ngẫu nhiên làm cho $\mathsf{ck}$ đều; đó là một khóa toàn
cục duy nhất được dẫn xuất từ một hạt giống công khai (không có khóa riêng cho
từng người dùng). $Q_H$ ký hiệu tổng số truy vấn oracle ngẫu nhiên của đối thủ,
bao gồm cả những truy vấn được thực hiện thay mặt nó bởi các truy vấn kiểm
chứng, cộng thêm một.

\subsection{Cú pháp của lược đồ hợp thành}
\label{sec:sec-syntax}

Tầng vận chuyển là một tập các hàm tất định của một hồ sơ vật mang đã ghim
$\mathsf{prof}$ và một khóa đồng bộ $k_s$: $\mathsf{Place}(\vid,\ell)$ trả về
các vị trí vật mang có thứ tự $P$ cho một envelope $\ell$ byte;
$\Canon(\vid,P)$ trả về digest đã chuẩn hóa theo vật mang $D$;
$\mathsf{Write}(\vid,P,E)$ ghi envelope $E$ vào các vật mang;
$\mathsf{Read}(\vid)$ trả về $(E,P,D)$ hoặc $\bot$. \emph{Tính đúng đắn vận
chuyển} nghĩa là $\mathsf{Read}(\mathsf{Write}(\vid,P,E))=(E,P,\Canon(\vid,P))$
với $P=\mathsf{Place}(\vid,|E|)$; đó là một tính chất hàm, và việc nó thất bại
chỉ gây ra từ chối.

\begin{itemize}
\item $\mathsf{Setup}$: cố định bộ tham số, $\pid$ (một digest SHA3-256 của tất
  cả các tham số và hạt giống thiết lập) và $\mathsf{ck}$.
\item $\mathsf{RegKeyGen}$, $\mathsf{RegSign}$: bên phát hành nắm một cặp khóa
  ML-DSA-65 $(\mathit{pk}_I,\mathit{sk}_I)$ và ký một sổ đăng ký
  $\mathsf{Reg}=(e,\Delta)$: một epoch và một tập các mô tả quan hệ. Một mô tả
  băm $\pid$, khóa, module quan hệ và chính sách (chính sách codec và hồ sơ vật
  mang); trong chế độ đã đăng ký nó còn băm một cam kết chứng chỉ $t_c$.
\item $\mathsf{Pins}$: cho $\mathsf{Reg}$, $\mathit{pk}_I$, $\mathsf{prof}$,
  một ngưỡng chống quay lui $e_{\min}$ và (chế độ đã đăng ký) $t_c$, kiểm tra
  chữ ký, tính lại mô tả kỳ vọng, đòi hỏi đúng một kết quả khớp, và xuất ra
  $\mathsf{pins}=(\mathit{rid},\mathit{root},e,\mathit{cpid},
  \mathit{prof\_id})$, hoặc thất bại.
\item $\mathsf{Issue}$: rút một thách thức đều 32 byte $N_s$ và lưu nó cùng ràng
  buộc $\mathsf{bind}(\mathsf{pins})$ và một thời điểm hết hạn.
\item $\mathsf{Embed}$: trên vật phủ $\vid$, tải trọng $p$ (in-band) hoặc chứng
  chỉ $(t_c,r_c)$ (đã đăng ký), $\mathsf{pins}$ và $N_s$: tính
  $P=\mathsf{Place}(\vid,\ell)$ từ một placeholder bằng không có độ dài cuối
  cùng $\ell$, $D=\Canon(\vid,P)$ và $x=\Stmt(\mathsf{pins},N_s,P,D)$; với
  in-band, lấy mẫu $r\leftarrow S_1^k$, đặt $t=\Com(\mathsf{enc}(p);r)$; chứng
  minh $\pi\leftarrow\mathsf{P}(x,t,r)$; lập $E=\texttt{"ZKLB"}\|1\|\pid\|N_s\|t\|\pi$
  (hoặc $\texttt{"ZKLR"}\|1\|\pid\|N_s\|\pi$) và xuất ra
  $\mathsf{Write}(\vid,P,E)$ và, với in-band, phần mở $(p,r)$.
\item $\mathsf{Verify}$: trên $\vid^*$: $(E^*,P^*,D^*)=\mathsf{Read}(\vid^*)$;
  phân tích $E^*$ một cách nghiêm ngặt; lấy $N_s^*$ từ $E^*$ (chế độ có lưu
  trữ) hoặc dùng một $N_s$ cố định (chế độ không trạng thái);
  $x^*=\Stmt(\mathsf{pins},N_s^*,P^*,D^*)$; chấp nhận khi và chỉ khi
  $\mathsf{V}(x^*,t^*,\pi^*)=1$ với $t^*$ từ $E^*$ (in-band) hoặc $t_c$ đã ghim
  (đã đăng ký) và, trong chế độ có lưu trữ, $N_s^*$ còn sống đối với
  $\mathsf{bind}(\mathsf{pins})$ và được tiêu thụ một cách nguyên tử.
\item $\mathsf{VerDisc}$: chấp nhận $(E,p,r)$ khi và chỉ khi $(r,1)$ là một phần
  mở trung thực của cam kết trong $E$ tới $\mathsf{enc}(p)$.
\item $\mathsf{Link}$, $\mathsf{VerLink}$: chứng minh đẳng thức
  $\Pi_{\mathsf{eq}}$ trên hai cam kết, được ràng buộc vào một digest ngữ cảnh
  $\mathit{lx}=\mathrm{SHA3}(\text{tag}\|\pid\|\text{ngữ cảnh có tiền tố độ dài})$.
\end{itemize}

Phát biểu là chuỗi 274 byte
\begin{multline*}
\Stmt=\mathtt{u16}(1)\|\mathit{rid}\|\pid\|\mathit{root}\|\mathtt{u64}(e)\|N_s\|
\mathit{cpid}\|\mathit{prof\_id}\\
\|\mathtt{u64}(|P|)\|\mathrm{SHA256}(P)\|V,\quad
V=\mathrm{SHA256}(\text{tag}\|D),
\end{multline*}
với $\mathrm{SHA256}(P)$ được lấy trên một mã hóa JSON chuẩn tắc của các vị trí
có thứ tự. Chứng minh mở $\Pi_{\mathsf{open}}$ xuất ra $\pi=(h,z)$ với
$h=\Hfs(\text{tag}\|\pid\|\mathtt{u64}(|x|)\|x\|\mathsf{enc}(t)\|\mathsf{enc}(w))$,
$w=B_1y$, $c=G(h)$, $z=y+cr\in Z$; bên kiểm chứng tính lại $w'=B_1z-ct_1$ và
kiểm tra $h$. Mọi trường sau $x$ đều có độ dài được cố định bởi $\pid$, nên đầu
vào của hàm băm giải mã được một cách duy nhất.

\subsection{Mô hình đối thủ}
\label{sec:sec-adversary}

Đối thủ $\mathcal A$ là một thuật toán xác suất cổ điển. Nó kiểm soát mọi kênh,
chọn các vật phủ và các tải trọng, biết $k_s$ (không tính chất nào dưới đây phụ
thuộc vào tính bí mật của nó), và có các oracle sau. $\mathcal O_{\mathsf{Iss}}$
trả về một $N_s$ mới. $\mathcal O_{\mathsf{Emb}}(\vid,p,N_s)$ chạy một
$\mathsf{Embed}$ in-band trung thực và trả về video stego;
$\mathcal O_{\mathsf{EmbR}}(j,\vid,N_s)$ làm điều tương tự cho chứng chỉ đã đăng
ký của bên nắm giữ trung thực $j$. Chúng tôi cho phép các oracle này chứng minh
các phát biểu $x$ tùy ý do $\mathcal A$ chọn, điều này chỉ làm mạnh thêm các kết
quả. Các đầu ra của chúng tạo thành danh sách $L$ các bộ ba $(x,t,\pi)$.
$\mathcal O_{\mathsf{Ver}}$ chạy $\mathsf{Verify}$ với kho dùng chung.
$\mathcal O_{\mathsf{Link}}$ trả về các chứng minh liên kết trung thực giữa các
cam kết đã ghi. $\mathcal O_{\mathsf{Disc}}(i)$ trả về phần mở của cam kết trung
thực thứ $i$ và đánh dấu nó là \emph{đã tiết lộ}.
$\mathcal O_{\mathsf{Reg}}(e,\Delta,T)$ khiến bên phát hành ký một sổ đăng ký
với epoch $e$ mà các mô tả của nó là các mô tả in-band $\Delta$ và các mô tả đã
đăng ký mà bên phát hành tính từ các cam kết chứng chỉ trong $T$, tất cả do
$\mathcal A$ chọn; $T_{\mathsf{reg}}$ là hợp của mọi $T$ đã truy vấn. (Một bên
phát hành chỉ ký một mô tả đã đăng ký cho một chứng chỉ mà nó đã thẩm định, nên
các mô tả đã đăng ký chỉ đi vào qua $T$.) $Q_P$ chặn số chứng minh được tạo bởi
các bên trung thực.

Bên phát hành và bên kiểm chứng là trung thực; kho phiên là tuyến tính hóa được
(linearizable) và bền vững. Việc làm hỏng bên nắm giữ là tĩnh (đối thủ đăng ký
các chứng chỉ của riêng nó qua $\mathcal O_{\mathsf{Reg}}$); tiết lộ là hình
thức làm hỏng thích nghi duy nhất, và các cam kết đã tiết lộ bị loại khỏi các
điều kiện thắng. Các kênh phụ, lỗi codec, và tính đúng đắn của tầng vận chuyển
nằm ngoài phạm vi.

\subsection{Tính không tiết lộ tri thức của chứng minh Fiat--Shamir-with-aborts}
\label{sec:sec-zk}

\begin{lemma}[Chấp nhận độc lập với nhân chứng]
\label{lem:hvzk}
Với mọi $r\in S_1^k$ và $c\in\Cset$, một lần thử với $y\leftarrow Y$ được chấp
nhận với xác suất đúng bằng $p_{\mathrm{acc}}$, và với điều kiện được chấp nhận
thì $z=y+cr$ đều trên $Z$.
\end{lemma}
\noindent
Mỗi hệ số $v$ của $cr$ thỏa $|v|\le\kappa$, nên $2B+1$ giá trị
$z_i\in[-B,B]$ tương ứng với các mặt nạ $z_i-v\in[-\gamma+1,\gamma-1]\subset
[-\gamma+1,\gamma]$, mỗi cái có xác suất $1/2\gamma$.

\begin{lemma}[Mô phỏng]
\label{lem:sim}
Tồn tại một bộ mô phỏng mà, chỉ với $(x,t)$ và quyền lập trình $\Hfs$, xuất ra
các chứng minh có phân phối hợp với các oracle ngẫu nhiên nằm trong khoảng
\begin{multline*}
\epszk(Q_P,Q_H)=2Q_P(Q_H+Q_P)2^{-256}\\
+(\Amax+1)Q_P\bigl(Q_H+(\Amax+1)Q_P\bigr)2^{-\eta}
\end{multline*}
so với $Q_P$ chứng minh thực, với mọi đối thủ thực hiện $Q_H$ truy vấn. Điều
tương tự đúng cho $\Pi_{\mathsf{eq}}$ với $p_{\mathrm{acc}}^2$ thay cho $p_{\mathrm{acc}}$.
\end{lemma}
\noindent
Bộ mô phỏng lấy mẫu số lần thử từ phân phối hình học bị cắt cụt với tham số
$p_{\mathrm{acc}}$, xuất ra $\bot$ với xác suất $(1-p_{\mathrm{acc}})^{\Amax}$,
và ngược lại lấy mẫu $h\leftarrow\{0,1\}^{256}$, $z\leftarrow Z$, đặt
$w=B_1z-G(h)t_1$ và lập trình $\Hfs$ tại $(\pid,x,t,w)$. Bước tế nhị là các truy
vấn oracle ngẫu nhiên mà bên chứng minh thực thực hiện trong các lần thử bị loại
và không bao giờ tiết lộ: các phân tích trước đây về Fiat--Shamir with aborts đã
xử lý sai việc tái lập trình tương ứng trong bước CMA-to-NMA, như được Barbosa
và cộng sự~\cite{ref:barbosa} chỉ ra và sửa chữa, và được Devevey và cộng
sự~\cite{ref:devevey} phân tích chi tiết. Ở đây bước này được giải quyết trong
ROM cổ điển nhờ hai sự kiện đặc thù của hiện thực này: sự chấp nhận là
\emph{hoàn toàn} độc lập với nhân chứng và với $c$ (Bổ đề~\ref{lem:hvzk}), và
mọi điểm truy vấn ẩn đều chứa $w=B_1y$ có min-entropy ít nhất $\eta$ (bởi vì
$B_1=[I_n|B_1']$). Tính độc lập chính xác cũng có nghĩa là số lần thử không rò
rỉ gì, dù rằng ngoài ra phần hiện thực không chạy trong thời gian hằng.

\subsection{Tính đúng đắn tri thức và ràng buộc vào video}
\label{sec:sec-soundness}

\begin{definition}[Ràng buộc video trích xuất được qua mô phỏng]
\label{def:se}
Trong thí nghiệm $\mathsf{SE}$, $\mathcal A$ tương tác với mọi oracle và xuất
ra $\vid^*$. Gọi $x^*$ là phát biểu mà bên kiểm chứng tính lại từ $\vid^*$ và
các mục ghim của chính nó, và $(t^*,\pi^*)$ là cam kết và chứng minh mà nó kiểm
chứng. $\mathcal A$ \emph{thắng} nếu $\mathsf{Verify}$ chấp nhận và
$(x^*,t^*,\pi^*)\notin L$. Một bộ trích xuất $\mathcal E^{\mathcal A}$
\emph{thành công} nếu nó xuất ra $x^*$ và một phần mở nới lỏng của $t^*$, trong
đó $x^*$ là phát biểu của bên kiểm chứng trong một lần chạy thắng của $\mathcal A$.
\end{definition}

\begin{theorem}[Tính đúng đắn tri thức ràng buộc vào phát biểu]
\label{thm:se}
Giả sử $\mathcal A$ thắng $\mathsf{SE}$ với xác suất $\varepsilon$. Tồn tại một
bộ trích xuất $\mathcal E$ chạy $\mathcal A$ hai lần với các oracle được mô
phỏng và thành công với xác suất $\varepsilon_{\mathsf{ext}}$ sao cho
\begin{multline*}
\varepsilon\le\epszk(Q_P,Q_H)+\Adv^{\MSIS}+Q_H2^{-256}\\
+\sqrt{Q_H\bigl(\varepsilon_{\mathsf{ext}}+Q_H2^{-256}+2^{-257}\bigr)}.
\end{multline*}
\end{theorem}
\noindent
\emph{Ý tưởng chứng minh.} Thay các chứng minh trung thực bằng các chứng minh
được mô phỏng (Bổ đề~\ref{lem:sim}). Nếu $\mathcal A$ tái dùng một điểm đã lập
trình, tức $(x^*,t^*)=(x_i,t_i)$ và $w^*=w_i$, thì $h^*=h_i$, do đó $c^*=c_i$ và
$B_1(z^*-z_i)=0$ với $0<\norm{z^*-z_i}_\infty\le2B$: một nghiệm M-SIS. Ngược
lại, truy vấn $(\pid,x^*,t^*,w^*)$ của bên kiểm chứng được trả lời bằng một giá
trị ngẫu nhiên mới, và bổ đề rẽ nhánh tổng quát của Bellare và
Neven~\cite{ref:bn06} (Bổ đề~1) với không gian đầu ra $2^{256}$ cho hai chứng
minh được chấp nhận có cùng $(x^*,t^*,w^*)$ và các hạt giống $h\neq h'$. Khi đó
$c=G(h)\neq G(h')=c'$ ngoại trừ với xác suất $2^{-257}+Q_H2^{-256}$, và
$B_1(z-z')=(c-c')t_1^*$ cho phần mở nới lỏng $(z-z',c-c')$. Hiệu thách thức là
khả nghịch theo Bổ đề~\ref{lem:invertible}, chính điều này khiến phần mở xác
định một thông điệp duy nhất. Bởi vì truy vấn bị rẽ nhánh chứa $x^*$, phần mở
được ràng buộc vào đúng phát biểu do bên kiểm chứng tính lại: cam kết video, các
vị trí vật mang, thách thức phiên, gốc sổ đăng ký và epoch, quan hệ và chính
sách của nó. \qed

Tổn thất căn bậc hai là cái giá thường thấy của việc tua lại (rewinding). Với
$Q_H=2^{64}$, cận biến một xác suất trích xuất $2^{-128}$ thành $\varepsilon
\lesssim2^{-32}$. Đối với tính xác thực (Định lý~\ref{thm:auth}), có một cận
chặt nhờ hướng SelfTargetMSIS của Kiltz, Lyubashevsky và
Schaffner~\cite{ref:kls18}, với cái giá là giả thiết phi chuẩn đó.

\begin{theorem}[Tính ràng buộc đối với các phần mở nới lỏng]
\label{thm:binding}
Không đối thủ nào có thể xuất ra $t$ và hai phần mở của $t$, mỗi phần mở trung
thực hoặc nới lỏng, với các thông điệp khác nhau, ngoại trừ với xác suất
$\Adv^{\MSIS}$.
\end{theorem}
\noindent
Cho $(\bar r,f,m)$ và $(\bar r',f',m')$, vector
$x=f'\bar r-f\bar r'$ thỏa $B_1x=0$; nếu $x=0$ thì
$\langle b_2,x\rangle=ff'(m-m')=0$, mâu thuẫn với tính khả nghịch. Chuẩn của nó
tối đa là $2\cdot2\kappa\cdot2B=\beta$ cho hai phần mở nới lỏng, $2\kappa+2B$
cho một trung thực và một nới lỏng, và $2$ cho hai trung thực, tất cả đều dưới
$q/2$ (\texttt{params.py:53}), nên $x\not\equiv0$ là một nghiệm M-SIS.

\noindent
Gộp lại, các Định lý~\ref{thm:se} và~\ref{thm:binding} nói rằng: ai tạo ra một
video in-band được chấp nhận đều bị ràng buộc vào thông điệp $\bar m$ mà bộ
trích xuất tính ra, và bất kỳ lần tiết lộ về sau nào mà $\mathsf{VerDisc}$ chấp
nhận đều bộc lộ một tải trọng $p$ với $\mathsf{enc}(p)=\bar m$, trừ khi M-SIS
bị phá.

\emph{Việc tiết lộ giữ cho tính ngẫu nhiên riêng tư.} Một lần tiết lộ là tải
trọng $p$ cùng với một chứng minh đẳng thức giữa $t$ và cam kết tất định
$\mathsf{Com}(\mathsf{enc}(p);0)$, được Fiat--Shamir ràng buộc vào một ngữ cảnh
do bên nhận chọn (định danh phát biểu và nonce). Bên nhận biết $p$ nhưng không
biết phần mở nào của $t$, nên nó không thể tạo các chứng minh mở mới cho $t$;
theo tính đúng đắn của chứng minh đẳng thức, một lần tiết lộ được chấp nhận vẫn
thỏa $\mathsf{enc}(p)=\bar m$ trừ khi M-SIS bị phá. Một phiên bản trước đây bộc
lộ chính $r$; đây là mục~D3 trong Phụ lục~\ref{app:code-check} và đã được sửa.
Các định lý dưới đây vẫn coi các cam kết mà \emph{phần mở} của chúng đã bị rò rỉ
(ví dụ do lộ khóa) là đã tiết lộ.

\begin{lemma}[Tính một chiều của các cam kết trung thực]
\label{lem:row}
Cho $\mathsf{ck}$ và $t^{(j)}=\Com(m_j;r_j)$, $j\le Q$, với
$r_j\leftarrow S_1^k$ và các thông điệp $m_j$ được chọn thích nghi trước khi mỗi
$t^{(j)}$ được rút, không đối thủ nào xuất ra một phần mở nới lỏng của bất kỳ
$t^{(j)}$ nào ngoại trừ với xác suất
$\Adv^{\mathsf{row}}_Q=\Adv^{\MSIS}+Q\bigl(\Adv^{\MLWE}+3^{kN}q^{-(n+1)N}\bigr)$.
\end{lemma}
\noindent
Đặt $x=fr_j-\bar r$, nên $B_1x=0$ và $\norm{x}_\infty\le2\kappa+2B$. Nếu
$x\neq0$ thì nó là một nghiệm M-SIS (phép rút gọn đã sinh mọi $t^{(j)}$ và biết
$r_j$). Nếu $x=0$ thì đối thủ đã tìm được $r_j=f^{-1}\bar r$, điều mà một bộ
phân biệt cấy mẫu M-LWE của nó tại một $j$ được đoán sẽ kiểm tra đối chiếu với
$t^{(j)}$; với $t^{(j)}$ đều, một phần mở bậc ba tồn tại với xác suất tối đa
$3^{kN}q^{-(n+1)N}$. Chỉ riêng tính ràng buộc là không đủ ở đây: một đối thủ
khôi phục được $r_j$ sẽ mở $t^{(j)}$ tới đúng thông điệp.

\subsection{Che giấu tải trọng}
\label{sec:sec-hiding}

\begin{definition}[Tính không phân biệt được tải trọng]
\label{def:hiding}
Trong $\mathsf{IND}^b$, $\mathcal A$ có mọi oracle và nộp một lần
$(\vid,p_0,p_1,N_s)$ với $|p_0|,|p_1|\le766$ byte; bên thách thức trả về
$\mathcal O_{\mathsf{Emb}}(\vid,p_b,N_s)$. $\mathcal A$ không được truy vấn
$\mathcal O_{\mathsf{Disc}}$ hoặc $\mathcal O_{\mathsf{Link}}$ trên cam kết
thách thức. $\Adv^{\mathsf{ind}}=|\Pr[\mathsf{IND}^1=1]-\Pr[\mathsf{IND}^0=1]|$.
\end{definition}

\begin{theorem}[Che giấu]
\label{thm:hiding}
$\Adv^{\mathsf{ind}}\le2\,\epszk(Q_P+1,Q_H)+2\Adv^{\MLWE}$.
\end{theorem}
\noindent
Video stego chỉ phụ thuộc vào $b$ thông qua $E$: $P$ và $D$ được cố định từ một
placeholder trước khi chứng minh, và envelope có kích thước cố định, nên độ dài
tải trọng cũng được che giấu. Mô phỏng chứng minh thách thức
(Bổ đề~\ref{lem:sim}), rồi thay $t$ bằng một cặp đều: với
$B_1'=[B_{1,2}'\,|\,B_{1,3}']$ và $r=(r_1,r_2,r_3)$,
$(t_1-B_{1,2}'(t_2-m),\,t_2-m)=A r_3+(r_1,r_2)$ với
$A=\bigl[\begin{smallmatrix}B_{1,3}'-B_{1,2}'b_2'\\ b_2'\end{smallmatrix}\bigr]$
đều, một mẫu M-LWE hạng $\lambda$ với $n+1$ mẫu vành; ánh xạ là một song ánh với
mỗi $m$ cố định. Khi đó thí nghiệm độc lập với $b$. Tính che giấu mang tính tính
toán và là tính chất duy nhất dựa trên M-LWE; nó đúng cho tải trọng, chứ không
cho sự kiện rằng một tải trọng tồn tại, vốn là một tính chất giấu tin được đánh
giá bằng thực nghiệm.

\subsection{Khả năng chống phát lại}
\label{sec:sec-replay}

\begin{definition}[Phát lại]
\label{def:replay}
Trong $\mathsf{RPL}$, $\mathcal A$ có mọi oracle đối với một bên kiểm chứng có
kho phiên và thắng nếu (i) hai lời gọi $\mathcal O_{\mathsf{Ver}}$ chấp nhận với
cùng $N_s$; hoặc (ii) một lời gọi chấp nhận với một $N_s$ không được phát hành
cho $\mathsf{bind}(\mathsf{pins})$, hoặc đã hết hạn; hoặc (iii) một lời gọi chấp
nhận $\vid^*$ có bộ ba $(x^*,t^*,\pi^*)\notin L$ và $\mathcal A$ đã truy vấn
$\Hfs$ tại điểm $(\pid,x^*,t^*,w^*)$ của bên kiểm chứng trước khi $N_s^*$ được
phát hành.
\end{definition}

\begin{theorem}[Khả năng chống phát lại]
\label{thm:replay}
$\Pr[\text{(i) hoặc (ii)}]\le\Adv^{\mathsf{cr}}_{\mathsf{SHA2}}+Q_I^22^{-257}$ và
$\Pr[\text{(iii)}]\le Q_HQ_I2^{-256}$, với $Q_I$ phiên được phát hành.
\end{theorem}
\noindent
(i) và (ii) suy ra từ ngữ nghĩa của kho: các hàng chỉ tồn tại cho các thách
thức đã phát hành (khóa theo một digest SHA-256 của $N_s$), và việc tiêu thụ là
một \texttt{UPDATE \dots\ WHERE consumed\_at IS NULL AND expires\_at>now} có
điều kiện đơn nhất bên trong một giao dịch ghi, nên nhiều nhất một bên gọi quan
sát được một hàng đã cập nhật (\texttt{zkp\_sessions.py:181--241}). Với (iii),
$N_s^*$ là một chuỗi con của $x^*$, nên một truy vấn thực hiện trước đó dự đoán
một giá trị 256 bit đều. Cùng với Định lý~\ref{thm:se}, một phiên được chấp nhận
tương ứng với một bên đã biết một phần mở và đã tính chứng minh sau khi phiên
được phát hành.

\emph{Điều mà một bên kiểm chứng không trạng thái không thể đem lại.} Với một
$N_s$ cố định (\texttt{expected\_session}), (i) thất bại: bất kỳ bản sao nào của
một video được chấp nhận đều lại được chấp nhận, bởi bất kỳ ai và bao nhiêu lần
tùy ý, vì một chứng minh không tương tác có thể chuyển nhượng. Nếu $N_s$ bị tái
dùng qua nhiều lần kiểm chứng, (iii) cũng thất bại và chỉ còn lại ràng buộc
video (Định lý~\ref{thm:se}). Không biến thể nào ngăn được việc tiếp chuyển
(relay): chế độ in-band không ràng buộc $N_s$ vào một bên chứng minh, nên ai
trình ra một video hợp lệ trước thì tiêu thụ phiên.

\subsection{Tính xác thực chế độ đã đăng ký}
\label{sec:sec-auth}

\begin{definition}[Tính xác thực]
\label{def:auth}
Trong $\mathsf{AUTH}$, bên thách thức tạo các chứng chỉ trung thực
$t_c^{(j)}=\Com(m_j;r_j)$, $j\le Q_C$, với $m_j$ do $\mathcal A$ chọn, và
$\mathcal A$ có mọi oracle. $\mathcal A$ xuất ra $(\mathsf{Reg}^*,\vid^*,t^*)$
và thắng nếu một bên kiểm chứng ghim $\mathit{pk}_I$ và $t^*$ chấp nhận $\vid^*$
với các mục ghim từ $\mathsf{Reg}^*$, bộ ba $(x^*,t^*,\pi^*)\notin L$, và hoặc
(a) $t^*\notin T_{\mathsf{reg}}$ (chưa bao giờ được bên phát hành đăng ký), hoặc
(b) $t^*=t_c^{(j)}$ cho một $j$ trung thực chưa tiết lộ.
\end{definition}

\begin{theorem}[Tính xác thực]
\label{thm:auth}
\begin{multline*}
\Adv^{\mathsf{auth}}\le\Adv^{\mathsf{euf\text{-}cma}}_{\mathsf{ML\text{-}DSA}}
+\Adv^{\mathsf{cr}}_{\mathsf{SHA2}}+\epszk(Q_P,Q_H)\\
+\Adv^{\MSIS}+Q_H2^{-256}\\
+\sqrt{Q_H\bigl(\Adv^{\mathsf{row}}_{Q_C}+Q_H2^{-256}+2^{-257}\bigr)}.
\end{multline*}
\end{theorem}
\noindent
Trường hợp (a) cho một chữ ký giả mạo ML-DSA trên các byte ký của sổ đăng ký
hoặc, nếu $\mathsf{Reg}^*$ đã được ký, một va chạm SHA-256 giữa các đầu vào có
độ dài cố định $B_1'\|b_2'\|t^*$ và $B_1'\|b_2'\|t'$ với một
$t'\in T_{\mathsf{reg}}$ nào đó, vì mô tả khớp chứa digest đó. Trong trường hợp
(b), Định lý~\ref{thm:se} áp dụng cho trò chơi được mô phỏng trích xuất một phần
mở nới lỏng của một chứng chỉ trung thực, điều mà Bổ đề~\ref{lem:row} loại trừ.
Định lý này nói về các chứng chỉ; sự liên kết tới một con người chỉ mạnh ngang
với việc bên phát hành kiểm tra bên nắm giữ trước khi ký. Khi bên phát hành ràng
buộc một định danh của bên nắm giữ, mô tả đã ký chứa một digest có tách miền của
nó bên cạnh cam kết, nên quan hệ đã ghim của bên kiểm chứng nêu tên cả hai
(Phụ lục~\ref{app:code-check}, mục~D2, đã sửa).

\subsection{Liên kết}
\label{sec:sec-linkage}

\begin{definition}[Tính đúng đắn và tính che giấu của liên kết]
\label{def:link}
Trong $\mathsf{LNK\text{-}S}$, $\mathcal A$ xuất ra một ngữ cảnh, $t_a,t_b$, một
$\pi_{\mathsf{eq}}$ được chấp nhận không do $\mathcal O_{\mathsf{Link}}$ xuất ra
cho bộ ba đó, và các phần mở (trung thực hoặc nới lỏng) của $t_a$ và $t_b$ với
các thông điệp khác nhau. Trong $\mathsf{LNK\text{-}H}^b$, $\mathcal A$ nộp
$(p_0,p_1,\mathit{ctx})$ và nhận $t_a=\Com(p_b;r_a)$,
$t_b=\Com(p_b;r_b)$ và $\pi_{\mathsf{eq}}$ cho $\mathit{ctx}$.
\end{definition}

\begin{theorem}[Liên kết]
\label{thm:link}
$\Adv^{\mathsf{lnk\text{-}s}}$ tuân theo cận của Định lý~\ref{thm:se} với
$\varepsilon_{\mathsf{ext}}\le\Adv^{\MSIS}$, và
$\Adv^{\mathsf{lnk\text{-}h}}\le2\,\epszk(Q_P+1,Q_H)+4\Adv^{\MLWE}$.
\end{theorem}
\noindent
Hai bản ghi bị rẽ nhánh với một thông báo $(w_a,w_b,u)$ cho các phần mở nới lỏng
$(\bar r_a,f)$, $(\bar r_b,f)$ với một $f$ \emph{chung} và
$f(t_{2,a}-t_{2,b})=\langle b_2,\bar r_a-\bar r_b\rangle$, do đó
$\bar m_a=\bar m_b$; hai phần mở của đối thủ khi đó mâu thuẫn với
Định lý~\ref{thm:binding} cho $t_a$ hoặc cho $t_b$. Việc liên kết tới một chứng
chỉ đã đăng ký còn xác thực thêm bên nắm giữ của nó, như trong
Định lý~\ref{thm:auth}, bởi vì việc trích xuất cũng cho ra một phần mở nới lỏng
của $t_c$. Tính che giấu dùng bộ mô phỏng cho $\Pi_{\mathsf{eq}}$ và hai bước
lai M-LWE; sự bằng nhau của hai tải trọng được bộc lộ theo thiết kế.

\subsection{Tính toàn vẹn của nội dung video}
\label{sec:sec-integrity}

$\Canon$ chia luồng bằng bộ phân đoạn đã ghim (\texttt{ffmpeg -c:v copy
-f segment}, $F$ khung mỗi đoạn) và kết hợp các digest của từng đoạn với
$F$ và các số đếm khung. Một đoạn mang các chunk envelope được băm qua một phép
tuần tự hóa có gắn thẻ, phân định bằng độ dài: với mỗi đơn vị NAL là chỉ số,
kiểu, \texttt{forbidden\_zero\_bit}, \texttt{nal\_ref\_idc} và kích thước mã
khởi đầu (start-code); RBSP nguyên văn cho các đơn vị NAL không có vật mang; và,
bên trong một đơn vị NAL IDR có vật mang, các bit RBSP giữa các khối vật mang
cho tới bit dừng (stop bit), với mỗi khối vật mang được thay bằng vector hệ số
của nó trong đó mỗi vật mang được chuẩn hóa (độ lớn làm tròn xuống số chẵn, giữ
dấu). Các đoạn khác được băm từng byte.

\begin{lemma}[Digest và envelope xác định những gì]
\label{lem:canon}
Phép tuần tự hóa giải mã được một cách duy nhất, nên $V$ bằng nhau kéo theo
$\Canon$ bằng nhau ngoại trừ với xác suất $3\Adv^{\mathsf{cr}}_{\mathsf{SHA2}}$.
Dưới quy tắc vật mang ổn định ($|\ell|\ge4$, một bit chẵn-lẻ), một mức vật mang
là $\mathrm{sgn}(\ell)(2\lfloor|\ell|/2\rfloor+b)$, nên $(\Canon,P)$ và các bit
vật mang xác định toàn bộ vector hệ số của mọi khối vật mang. Nếu chuỗi bit vật
mang là mã hóa duy nhất của $E$ dưới hồ sơ đã ghim, thì $(\Canon,P,E)$ xác định:
sự phân đoạn và các số đếm khung; mọi byte của các đoạn không có vật mang; mọi
trường header NAL và bit RBSP bên ngoài các khối vật mang; và các vector hệ số
của các khối vật mang.
\end{lemma}

\begin{theorem}[Tính toàn vẹn]
\label{thm:integrity}
Giả sử một lần nhúng trung thực tạo ra $\vid_h$ cho phiên $N_s$ với cam kết $t$
(hoặc chứng chỉ $t_c$), và $\mathcal A$ xuất ra $\vid^*$ được chấp nhận cho
$N_s$ với cùng cam kết (hoặc chứng chỉ), không tiết lộ. Khi đó
$\Canon(\vid^*)=\Canon(\vid_h)$, $P^*=P_h$ và $E^*=E_h$, ngoại trừ với xác suất
tối đa
\begin{multline*}
4\Adv^{\mathsf{cr}}_{\mathsf{SHA2}}+\epszk(Q_P,Q_H)+\Adv^{\MSIS}+Q_H2^{-256}\\
+\sqrt{Q_H\bigl(\Adv^{\mathsf{row}}_{Q_P}+Q_H2^{-256}+2^{-257}\bigr)}.
\end{multline*}
\end{theorem}
\noindent
Nếu $x^*\neq x_h$, bộ ba là mới và việc trích xuất cho ra một phần mở nới lỏng
của một cam kết trung thực. Nếu $x^*=x_h$, giá trị băm vị trí và $V$ khớp nhau,
và Bổ đề~\ref{lem:canon} được áp dụng; khi đó chứng minh chính là chứng minh
trung thực, ngược lại trường hợp điểm-đã-lập-trình của Định lý~\ref{thm:se} cho
một nghiệm M-SIS.

\paragraph{Các giới hạn của khẳng định toàn vẹn}
\begin{itemize}
\item Chế độ in-band ràng buộc nội dung vào một \emph{cam kết}, chứ không vào
  một con người: bất kỳ ai nắm một $N_s$ còn sống đều có thể nhúng lại bất kỳ
  video đã chỉnh sửa nào với một cam kết mới của riêng nó, và nó sẽ được chấp
  nhận.
\item Các vật mang chỉ được bảo vệ như các bit envelope. Định lý~\ref{thm:integrity}
  cần các bit vật mang là mã hóa duy nhất của $E$. Bên kiểm chứng chia lại chunk
  của $E$ theo kích thước chunk đã ghim và từ chối trừ khi các mảnh trích xuất
  được đúng là mã hóa đó; nếu không có phép kiểm này, một envelope một-chunk sẽ
  để trống trường \texttt{chunk\_size} 16 bit của nó (mục~D1 trong
  Phụ lục~\ref{app:code-check}, đã sửa).
\item Không được ràng buộc: các byte không sống sót qua bộ phân đoạn đã ghim (dữ
  liệu bộ chứa, các luồng khác, dữ liệu ngoài các đơn vị NAL); các bit căn lề
  RBSP sau bit dừng trong các đơn vị NAL IDR có mang vật mang; các từ mã CAVLC
  thay thế cho cùng một vector hệ số của một khối vật mang; tầng chống giả lập
  (digest bao phủ RBSP); và phiên bản của bộ phân đoạn, vốn được chính sách
  codec nêu tên nhưng không ghim.
\item Không có khẳng định về tính đúng đắn codec: digest bao phủ cú pháp, không
  bao phủ các điểm ảnh đã giải mã. Không có khẳng định về nguồn gốc: không gì
  nói rằng một máy quay đã ghi $\vid_h$ hay rằng tải trọng là đúng sự thật.
\end{itemize}

\subsection{Các đối thủ lượng tử}
\label{sec:sec-qrom}

Mọi định lý ở trên đều là các phát biểu trong ROM cổ điển. Việc tua lại
(Định lý~\ref{thm:se}) và việc lập trình thích nghi (Bổ đề~\ref{lem:sim}) không
chuyển sang QROM như đã phát biểu. Barbosa và cộng sự~\cite{ref:barbosa} đưa ra
các chứng minh QROM của bước CMA-to-NMA cho Fiat--Shamir with aborts, và Don và
cộng sự~\cite{ref:dfms19} cùng Liu và Zhandry~\cite{ref:lz19} chỉ ra rằng phép
biến đổi Fiat--Shamir bảo toàn tính đúng đắn (tri thức) của các $\Sigma$-protocol
phù hợp trước các đối thủ lượng tử với tổn thất đa thức. Chúng tôi chưa kiểm
chứng rằng các giả thiết của họ đúng cho bộ trích xuất nới lỏng và bên chứng
minh có loại bỏ được dùng ở đây, nên chúng tôi không khẳng định định lý QROM
nào. Cách diễn đạt của chúng tôi là ``hậu lượng tử một cách hợp lý'': các phép
rút gọn nhắm tới M-SIS, M-LWE, ML-DSA và va chạm băm, không cái nào có một tấn
công lượng tử hiệu quả đã biết ở các tham số này.
, and after \appendices add
% % Full proofs for sections/security.tex and the specification/code
% cross-check. Input after \appendices in main.tex. Uses the macros that
% sections/security.tex provides.

\section{Các chứng minh an toàn}
\label{app:security-proofs}

\subsection{Quy ước}

Các oracle ngẫu nhiên được lấy mẫu lười (lazily). Một điểm của $\Hfs$ được
\emph{xác định} một khi nó đã được truy vấn hoặc lập trình; một truy vấn tại
một điểm chưa xác định là \emph{mới} (fresh). $Q_H$ chặn mọi truy vấn oracle
ngẫu nhiên của đối thủ, bao gồm các truy vấn của bên kiểm chứng trong
$\mathcal O_{\mathsf{Ver}}$ và một lần kiểm chứng cuối, còn $Q_P$ chặn số chứng
minh trung thực. Mọi phép rút gọn dưới đây tự chạy toàn bộ thí nghiệm, điều mà
nó làm được vì việc kiểm chứng là công khai và mọi cam kết trung thực đều do
bên thách thức sinh ra. Đối với oracle mở rộng khóa có một điểm duy nhất (thẻ
khóa theo sau bởi $\pid$); một phép rút gọn cần cấy một ma trận thách thức sẽ
lập trình điểm này trước khi đối thủ bắt đầu, điều này khiến $\mathsf{ck}$ được
phân phối đúng như một khóa đều.

\subsection{Chứng minh Bổ đề~\ref{lem:hvzk}}

Cố định $r\in S_1^k$ và $c\in\Cset$, và đặt $v=cr\in\mathbb Z^{kN}$ là tích
negacyclic nguyên. Vì $\norm{v}_\infty\le\norm{c}_1\norm{r}_\infty
\le\kappa<q/2$, việc tâm hóa tích theo modulo $q$ khôi phục $v$ một cách chính
xác, đúng như những gì phần hiện thực tính (\texttt{opening\_proof.py:79--80}).
Với một tọa độ và bất kỳ $a\in[-B,B]$, $\Pr[y_i+v_i=a]=\Pr[y_i=a-v_i]=1/2\gamma$,
bởi vì $a-v_i\in[-B-\kappa,B+\kappa]=[-\gamma+1,\gamma-1]$ nằm trong dải mặt nạ.
Các tọa độ độc lập, nên $\Pr[z=z_0]=(2\gamma)^{-kN}$ với mọi $z_0\in Z$. Lấy
tổng trên $Z$ cho xác suất chấp nhận $|Z|(2\gamma)^{-kN}=p_{\mathrm{acc}}$, và
với điều kiện được chấp nhận thì $z$ đều trên $Z$; không đại lượng nào phụ thuộc
vào $r$ hay $c$. \qed

\subsection{Chứng minh Bổ đề~\ref{lem:sim}}

\emph{Bên chứng minh thực} (\texttt{opening\_proof.py:86--109}). Trên $(x,t,r)$:
rút $e\leftarrow\{0,1\}^{256}$; $K=H_{\mathsf{hg}}(e\|x,t,r\text{ có tiền tố độ dài})$;
với $i=1,\dots,\Amax$: $y_i=H_{\mathsf y}(K\|i)\in Y$, $w_i=B_1y_i$,
$h_i=\Hfs(\pid,x,t,w_i)$, $c_i=G(h_i)$, $z_i=y_i+c_ir$; xuất ra $(h_i,z_i)$ tại
$i$ đầu tiên có $z_i\in Z$, hoặc $\bot$ sau $\Amax$ lần loại bỏ.

\emph{Bộ mô phỏng} $\mathcal S(x,t)$. Xuất ra $\bot$ với xác suất
$(1-p_{\mathrm{acc}})^{\Amax}$. Ngược lại rút $h\leftarrow\{0,1\}^{256}$ và
$z\leftarrow Z$, đặt $w=B_1z-G(h)t_1$; nếu $\Hfs$ đã xác định tại
$(\pid,x,t,w)$ thì loại bỏ; ngược lại lập trình $\Hfs(\pid,x,t,w):=h$ và xuất ra
$(h,z)$. Bộ mô phỏng không bao giờ dùng $r$.

Chúng tôi chặn độ thay đổi trên mỗi chứng minh; một cận hợp (union bound) trên
$Q_P$ chứng minh cho ta bổ đề. Viết $\eta_y=nN\log_2(2\gamma)$. Bởi vì
$B_1=[I_n|B_1']$, $n$ hàng đầu của $w=B_1y$ là $y_{[1..n]}+B_1'y_{[n+1..k]}$,
một ảnh song ánh của $y_{[1..n]}$ với $y_{[n+1..k]}$ cố định; do đó
$H_\infty(w)\ge\eta_y$ cho một mặt nạ đều,
$H_\infty(w)\ge\eta_y-\log_2\frac1{1-p_{\mathrm{acc}}}\ge\eta$ cho một mặt nạ
với điều kiện bị loại bỏ, và $H_\infty(w)\ge\eta$ cho một $w=B_1z-ct_1$ được
mô phỏng với $z$ đều trên $Z$.

\begin{itemize}
\item[$G_0$:] bên chứng minh thực.
\item[$G_1$:] $K$ được thay bằng một giá trị đều. Hai trò chơi chỉ khác nhau
  nếu một truy vấn nào đó tới $H_{\mathsf{hg}}$ khác với truy vấn của chính bên
  chứng minh bắt đầu bằng $e$: tối đa $(Q_H+Q_P)2^{-256}$.
\item[$G_2$:] mỗi $y_i$ được thay bằng một phần tử đều mới của $Y$. Hai trò
  chơi chỉ khác nhau nếu $H_{\mathsf y}$ bị truy vấn ở nơi khác tại $K\|i$: tối
  đa $(Q_H+Q_P)2^{-256}$.
\item[$G_3$:] trong một lần thử bị loại, bên chứng minh rút $h_i$ một cách đều
  và đặt $c_i=G(h_i)$ mà không truy vấn hay xác định $\Hfs$. Hai trò chơi chỉ
  khác nhau nếu điểm $(\pid,x,t,w_i)$, tại bất kỳ thời điểm nào, bị $\mathcal A$
  truy vấn hoặc được một chứng minh trung thực khác xác định: tối đa
  $\Amax(Q_H+\Amax Q_P)2^{-\eta}$. Đây là bước liên quan tới các truy vấn
  oracle ngẫu nhiên ẩn của các lần thử bị loại~\cite{ref:barbosa,ref:devevey}.
\item[$G_4$:] điểm của lần thử được chấp nhận được xác định bởi truy vấn của
  bên chứng minh; nó là mới ngoại trừ với xác suất $(Q_H+(\Amax+1)Q_P)2^{-\eta}$.
  Với điều kiện là mới, $h_i$ đều trong mọi lần thử và, theo
  Bổ đề~\ref{lem:hvzk}, lần thử $i$ được chấp nhận với xác suất đúng bằng
  $p_{\mathrm{acc}}$ bất kể $h_i$, $c_i$ và $r$. Vậy chỉ số của lần chấp nhận
  đầu tiên là hình học bị cắt cụt, và với điều kiện được chấp nhận, $h$ đều và
  $z$ đều trên $Z$, độc lập với nhau và với $r$. Các lần thử bị loại không còn
  liên kết với oracle nào sau $G_3$ và có thể được bỏ đi. Khi đó đầu ra và định
  nghĩa oracle $\Hfs(\pid,x,t,B_1z-G(h)t_1)=h$ được phân phối đúng như trong
  $\mathcal S$, mà sự kiện loại bỏ của nó chính là sự kiện không-mới vừa được
  chặn.
\end{itemize}
Tổng trên mỗi chứng minh là $2(Q_H+Q_P)2^{-256}+(\Amax+1)(Q_H+(\Amax+1)Q_P)2^{-\eta}$.

Đối với $\Pi_{\mathsf{eq}}$ (\texttt{equality\_proof.py:58--104}), bên chứng
minh rút hai mặt nạ mỗi lần thử, thông báo $(w_a,w_b,u)=(B_1y_a,B_1y_b,\langle b_2,y_a-y_b\rangle)$
và chấp nhận khi và chỉ khi cả hai đáp ứng đều nằm trong $Z$. Theo
Bổ đề~\ref{lem:hvzk} áp dụng cho mỗi đáp ứng, sự chấp nhận có xác suất đúng bằng
$p_{\mathrm{acc}}^2$ và $(z_a,z_b)$ được chấp nhận đều trên $Z^2$. Bộ mô phỏng
lấy mẫu $(h,z_a,z_b)$, đặt $w_a=B_1z_a-ct_{1,a}$, $w_b=B_1z_b-ct_{1,b}$,
$u=\langle b_2,z_a-z_b\rangle-c(t_{2,a}-t_{2,b})$ và lập trình $H_{\mathsf{eq}}$;
thông báo có min-entropy ít nhất $\eta$ thông qua $w_a$. Cùng các bước lai cho
cùng cận. \qed

\subsection{Một bổ đề trích xuất bằng rẽ nhánh}

\begin{lemma}[Trích xuất]
\label{lem:fork}
Gọi $\mathcal B$ là một thuật toán với các đồng xu $\rho$ thực hiện nhiều nhất
$Q$ truy vấn mới tới $\Hfs$, được trả lời bởi $h_1,\dots,h_Q\leftarrow\{0,1\}^{256}$
(oracle $G$ và mọi điểm được $\mathcal B$ lập trình đều do $\rho$ xác định).
$\mathcal B$ xuất ra $\bot$ hoặc $(J,x^*,t^*,\pi^*)$ sao cho
$\pi^*=(h_J,z^*)$ được chấp nhận cho $(x^*,t^*)$ và truy vấn mới thứ $J$ là
$(\pid,x^*,t^*,w^*)$ với $w^*=B_1z^*-G(h_J)t_1^*$. Đặt
$\mathrm{acc}=\Pr[\mathcal B\neq\bot]$. Tồn tại $\mathcal E$, chạy
$\mathcal B$ hai lần, xuất ra $x^*$ của lần chạy thứ nhất và một phần mở nới
lỏng của $t^*$ với xác suất $\varepsilon_{\mathsf{ext}}\ge\mathrm{acc}(\mathrm{acc}/Q-2^{-256})
-2^{-257}-Q2^{-256}$; tương đương
\[
\mathrm{acc}\le Q2^{-256}+\sqrt{Q\bigl(\varepsilon_{\mathsf{ext}}+2^{-257}+Q2^{-256}\bigr)}.
\]
\end{lemma}
\begin{proof}
$\mathcal E$ là thuật toán rẽ nhánh của Bellare và Neven~\cite{ref:bn06}
(Bổ đề~1, bổ đề rẽ nhánh tổng quát; xem thêm~\cite{ref:ps00}): chạy
$\mathcal B(\rho;h_1..h_Q)$; nếu nó trả về $J$, chạy lại nó trên
$(\rho;h_1..h_{J-1},h'_J..h'_Q)$ với $h'_J,\dots,h'_Q$ mới; thành công nếu lần
chạy thứ hai trả về cùng $J$ và $h'_J\neq h_J$. Bổ đề đó cho xác suất thành công
ít nhất $\mathrm{acc}(\mathrm{acc}/Q-2^{-256})$. Hai lần chạy giống hệt nhau cho
tới truy vấn mới thứ $J$, nên chúng chia sẻ $(x^*,t^*,w^*)$, và trả về $(h_J,z)$,
$(h'_J,z')$ với $B_1z-ct_1^*=w^*=B_1z'-c't_1^*$ cho $c=G(h_J)$, $c'=G(h'_J)$.
Giá trị $h'_J$ đều và độc lập với lần chạy thứ nhất; trừ khi nó là một trong
nhiều nhất $Q$ đầu vào mà tại đó $G$ đã được tính trong lần chạy thứ nhất
(xác suất $\le Q2^{-256}$), $G(h'_J)$ là một phần tử đều mới của $\Cset$ và
bằng $c$ với xác suất $1/|\Cset|\le2^{-257}$. Ngược lại
$f=c-c'\in\Cdiff$ và $\bar r=z-z'$ thỏa $\norm{\bar r}_\infty\le2B$ và
$B_1\bar r=ft_1^*$. Dạng thứ hai suy ra bằng cách giải bất đẳng thức bậc hai:
$\mathrm{acc}^2-Q2^{-256}\mathrm{acc}-Q\delta\le0$ với
$\delta=\varepsilon_{\mathsf{ext}}+2^{-257}+Q2^{-256}$ kéo theo
$\mathrm{acc}\le Q2^{-256}+\sqrt{Q\delta}$.
\end{proof}

\subsection{Chứng minh Định lý~\ref{thm:se}}

$G_0$ là $\mathsf{SE}$. Trong $G_1$, mọi chứng minh trung thực đều do bộ mô
phỏng tạo ra, nên $|\Pr[W_0]-\Pr[W_1]|\le\epszk(Q_P,Q_H)$, trong đó $W_i$ là sự
kiện thắng. Gọi $\mathsf{Prog}$ là sự kiện $\mathcal A$ thắng trong $G_1$ và
điểm $(\pid,x^*,t^*,w^*)$ của bên kiểm chứng đã được bộ mô phỏng lập trình, cho
một bộ ba $(x_i,t_i,\pi_i)\in L$. Khi đó $x^*=x_i$, $t^*=t_i$, $w^*=w_i$, và
sự chấp nhận buộc $h^*=h_i$, nên $c^*=c_i$. Vì $(x^*,t^*,\pi^*)\notin L$, ta
có $z^*\neq z_i$, và $B_1z^*-c_it_1^*=B_1z_i-c_it_1^*$ cho
$B_1(z^*-z_i)=0$ với $0<\norm{z^*-z_i}_\infty\le2B\le\beta$. Một đối thủ M-SIS
cấy ma trận của nó làm $B_1'$, tự rút $b_2'$, chạy $G_1$ và xuất ra $z^*-z_i$
thành công mỗi khi $\mathsf{Prog}$ xảy ra, nên
$\Pr[\mathsf{Prog}]\le\Adv^{\MSIS}$. Ngược lại, điểm của bên kiểm chứng được
xác định bởi một truy vấn mới, và $\mathcal B$ (chạy $G_1$ và xuất ra chỉ số
của truy vấn đó) thỏa giả thiết của Bổ đề~\ref{lem:fork} với
$Q=Q_H$ và $\mathrm{acc}\ge\varepsilon-\epszk-\Adv^{\MSIS}$. Thay vào cho ta
định lý.

Truy vấn bị rẽ nhánh chứa $x^*$, thứ mà bên kiểm chứng dựng từ các mục ghim của
nó và từ $\mathsf{Read}(\vid^*)$ (\texttt{video\_pipeline.py:278--284}); nó
không bao giờ lấy một trường phát biểu nào từ bên chứng minh ngoại trừ $N_s^*$,
thứ mà nó kiểm tra đối chiếu với kho của nó (hoặc với giá trị cố định) và bản
thân là một phần của $x^*$. Do đó phần mở trích xuất được được ràng buộc vào các
giá trị của bên kiểm chứng cho $\mathit{rid}$, $\pid$, gốc sổ đăng ký và epoch,
$N_s$, chính sách codec, hồ sơ vật mang, $|P|$, $\mathrm{SHA256}(P)$ và $V$. Hai
phát biểu khác nhau ở bất kỳ trường nào đều khác nhau khi là chuỗi, vì mọi
trường đều có độ rộng cố định (\texttt{video\_proof.py:68--77}). \qed

\subsection{Chứng minh Định lý~\ref{thm:binding}}

Cho $(\bar r,f,m)$ và $(\bar r',f',m')$ mở $t$ với $m\neq m'$, trong đó một
phần mở trung thực có $f=1$ và $\norm{\bar r}_\infty\le1$. Từ $ft_1=B_1\bar r$
và $f't_1=B_1\bar r'$ ta được $B_1x=0$ cho $x=f'\bar r-f\bar r'$. Từ
$ft_2=\langle b_2,\bar r\rangle+fm$ và $f't_2=\langle b_2,\bar r'\rangle+f'm'$,
\[
\langle b_2,x\rangle=f'(ft_2-fm)-f(f't_2-f'm')=ff'(m'-m).
\]
Nếu $x\equiv0\pmod q$ thì $ff'(m'-m)=0$ và, vì $f$ và $f'$ khả nghịch
(Bổ đề~\ref{lem:invertible}; $1$ hiển nhiên), $m=m'$, mâu thuẫn. Chuẩn
của $x$ tối đa là $\norm{f'}_1\norm{\bar r}_\infty+\norm{f}_1\norm{\bar r'}_\infty$,
tức $2\kappa\cdot2B+2\kappa\cdot2B=8\kappa B$ cho hai phần mở nới lỏng,
$2\kappa\cdot1+1\cdot2B$ cho một nới lỏng và một trung thực, và $2$ cho hai phần
mở trung thực. Tất cả đều tối đa $\beta<q/2$ (hàm khởi tạo tham số bắt buộc
$2\beta<q$, \texttt{params.py:53--54}), nên $x$ là một vector ngắn khác không
cả theo modulo $q$ lẫn trên $\mathbb Z$. Phép rút gọn cấy ma trận của nó làm
$B_1'$, lấy mẫu $b_2'$ và xuất ra $x$. \qed

\subsection{Proof of Lemma~\ref{lem:row}}

\begin{lemma}[Hiding instance]
\label{lem:mlwe-instance}
There is an efficient map that, for any $m$ chosen after seeing
$\mathsf{ck}$, sends an M-LWE instance $(A,u)$, $A\in\Rq^{(n+1)\times\lambda}$,
to $(\mathsf{ck},t)$ with $\mathsf{ck}$ uniform, such that $t=\Com(m;r)$ with
$r\leftarrow S_1^k$ if $u=As+e$, and $t$ is uniform if $u$ is uniform.
\end{lemma}
\begin{proof}
Write $A=\bigl[\begin{smallmatrix}A_1\\a_2\end{smallmatrix}\bigr]$ and
$u=(u_1,u_2)$. Draw $B_{1,2}'\leftarrow\Rq^{n\times1}$ and set $b_2'=a_2$,
$B_{1,3}'=A_1+B_{1,2}'a_2$, which makes $\mathsf{ck}$ uniform. Given $m$, set
$t_2=u_2+m$ and $t_1=u_1+B_{1,2}'u_2$. If $u=As+e$ with $e=(e_1,e_2)$ and
$r=(e_1,e_2,s)$, then $\langle b_2,r\rangle=e_2+a_2s=u_2$ and
$B_1r=e_1+B_{1,2}'e_2+(A_1+B_{1,2}'a_2)s=u_1+B_{1,2}'u_2$, so $t=\Com(m;r)$
with $r$ uniform ternary. If $u$ is uniform, so is $t$, because
$(u_1,u_2)\mapsto(u_1+B_{1,2}'u_2,u_2+m)$ is a bijection.
\end{proof}

Now let $\mathcal A$ output $(j,\bar r,f)$ with $f\in\Cdiff$,
$\norm{\bar r}_\infty\le2B$ and $ft^{(j)}_1=B_1\bar r$, and put $x=fr_j-\bar r$.
Then $B_1x=ft_1^{(j)}-ft_1^{(j)}=0$ and $\norm{x}_\infty\le2\kappa+2B<q/2$.

\emph{Case $x\neq0$.} An M-SIS adversary plants its matrix as $B_1'$, samples
$b_2'$, generates every $t^{(j)}$ itself (so it knows every $r_j$), and outputs
$x$. Hence $\Pr[x\neq0]\le\Adv^{\MSIS}$.

\emph{Case $x=0$.} Then $\bar r\equiv fr_j$, so $r_j$ is the centered
representative of $f^{-1}\bar r\bmod q$. An M-LWE distinguisher guesses
$j_0\leftarrow[Q]$, builds $\mathsf{ck}$ and $t^{(j_0)}$ from its instance by
Lemma~\ref{lem:mlwe-instance} (with the message $m_{j_0}$ that $\mathcal A$
chooses), generates the other commitments honestly, and outputs $1$ iff
$\mathcal A$ returns $j=j_0$ and $r'=\mathrm{center}(f^{-1}\bar r)$ is ternary
with $\Com(m_{j_0};r')=t^{(j_0)}$. On a real instance it outputs $1$ with
probability at least $\Pr[x=0]/Q$. On a uniform instance $t^{(j_0)}$ is uniform
and independent of $m_{j_0}$, so a ternary opening exists with probability at
most $|S_1^k|/q^{(n+1)N}=3^{kN}q^{-(n+1)N}$. Hence
$\Pr[x=0]\le Q(\Adv^{\MLWE}+3^{kN}q^{-(n+1)N})$. \qed

\subsection{Proof of Theorem~\ref{thm:hiding}}

$H_0$ is $\mathsf{IND}^0$. In $H_1$ all honest proofs, including the challenge
proof, are simulated: $|H_0-H_1|\le\epszk(Q_P+1,Q_H)$. In $H_2$ the challenge
commitment is replaced by a uniform pair; a distinguisher between $H_1$ and
$H_2$ is an M-LWE distinguisher by Lemma~\ref{lem:mlwe-instance} with
$m=\mathsf{enc}(p_0)$: it builds $\mathsf{ck}$ from its instance, generates the
other commitments honestly, answers $\mathcal O_{\mathsf{Disc}}$ and
$\mathcal O_{\mathsf{Link}}$ for them, and the challenge randomness is never
needed because the challenge proof is simulated and disclosure and linkage of
the challenge are excluded. So $|H_1-H_2|\le\Adv^{\MLWE}$. The stego video in
$H_2$ is $\mathsf{Write}(\vid,P,E)$ with $P$, $D$, $x$ independent of the
payload (fixed from a placeholder of the final length,
\texttt{video\_pipeline.py:216--227}), and $E$ contains only $\pid$, $N_s$, a
uniform $t$ and a simulated proof; $H_2$ is therefore independent of $b$. The
same three hybrids from $\mathsf{IND}^1$ end in the same $H_2$, which gives
the bound. The payload length is hidden because $\mathsf{enc}$ maps every
payload of at most 766 bytes to one polynomial and the envelope has a fixed
size. \qed

\subsection{Proof of Theorem~\ref{thm:replay}}

We model the store as a linearizable table keyed by
$d(N)=\mathrm{SHA256}(\text{tag}\|N)$ (\texttt{zkp\_sessions.py:113--117}).
Rows are inserted only by $\mathsf{Issue}$ with a fresh uniform $N$
(\texttt{:132--151}) and deleted only once expired (\texttt{:140--143}).
Acceptance requires the preflight \texttt{is\_active} and then
\texttt{consume}, an \texttt{UPDATE} with the conditions
\texttt{consumed\_at IS NULL}, matching binding and \texttt{expires\_at>now},
committed inside \texttt{BEGIN IMMEDIATE}, returning success iff exactly one
row changed (\texttt{:181--205}).

(i) After the first successful \texttt{consume} of a row, the condition
\texttt{consumed\_at IS NULL} fails for that row forever; a second success for
the same $N_s$ needs another row with key $d(N_s)$, i.e.\ a later issuance of
some $N'$ with $d(N')=d(N_s)$: either $N'=N_s$ (probability at most
$Q_I^22^{-257}$ over all pairs of issuances) or a SHA-256 collision.
(ii) Any row with key $d(N_s^*)$ comes from an issuance of some $N$ with
$d(N)=d(N_s^*)$, so $N_s^*$ was issued unless SHA-256 collides; the binding
and expiry are part of the \texttt{WHERE} clause, evaluated at consumption
time. (iii) The session field occupies bytes $106$--$137$ (0-indexed) of every statement,
and $x^*$ is a substring of the point $(\pid,x^*,t^*,w^*)$. A query made
before the issuance of $N_s^*$ fixes this field before $N_s^*$ is drawn, so it
matches with probability $2^{-256}$; a union bound over $Q_H$ queries and $Q_I$
sessions gives $Q_HQ_I2^{-256}$. \qed

\subsection{Proof of Theorem~\ref{thm:auth}}

Let $W_a$, $W_b$ be the two winning cases.

$W_a$: $\mathsf{Pins}$ verifies $\mathsf{Reg}^*$ under $\mathit{pk}_I$
(\texttt{video\_pipeline.py:146--147}). If its signing bytes were never signed
through $\mathcal O_{\mathsf{Reg}}$, it is an EUF-CMA forgery; the reduction
runs the whole experiment, using its signing oracle for
$\mathcal O_{\mathsf{Reg}}$. Otherwise $\mathsf{Reg}^*$ is a signed registry
and contains a descriptor equal to the registered descriptor of $t^*$
(\texttt{:150--153}). That descriptor has the registered constraint-module
hash, so it is not in any $\Delta$ (in-band descriptors have a different
module hash, unless SHA-256 collides on the two fixed module strings), and it
was computed from some $t'\in T_{\mathsf{reg}}$. Since $t^*\notin T_{\mathsf{reg}}$,
$t'\neq t^*$ and the canonical encodings differ, so
$\mathrm{SHA256}(B_1'\|b_2'\|t^*)=\mathrm{SHA256}(B_1'\|b_2'\|t')$
(\texttt{:119--128}) is a collision. Hence
$\Pr[W_a]\le\Adv^{\mathsf{euf\text{-}cma}}_{\mathsf{ML\text{-}DSA}}+\Adv^{\mathsf{cr}}_{\mathsf{SHA2}}$.

$W_b$: simulate all honest proofs ($\epszk$) and handle the programmed-point
case as in Theorem~\ref{thm:se} ($\Adv^{\MSIS}$). The remaining runner
$\mathcal B$ receives the honest credentials as a $Q_C$-instance challenge of
Lemma~\ref{lem:row}, never uses their randomness (all proofs are simulated,
and a disclosure query for an instance aborts the run, which does not affect
runs in which $\mathcal A$ wins with an undisclosed $j$), and outputs the
index of the verifier's fresh query. Lemma~\ref{lem:fork} extracts a relaxed
opening of $t^*=t_c^{(j)}$, which is a Lemma~\ref{lem:row} adversary, so
$\varepsilon_{\mathsf{ext}}\le\Adv^{\mathsf{row}}_{Q_C}$. Adding the cases
gives the theorem.

A tight alternative for $W_b$ follows the ML-DSA analysis: after replacing the
credential by a uniform pair (M-LWE), an accepted proof is a solution of the
SelfTargetMSIS problem for $[B_1\,|\,t_1]$ and $\Hfs$~\cite{ref:kls18}, without
rewinding. We state the rewinding bound because it uses only the estimated
M-SIS and M-LWE instances. \qed

\subsection{Proofs for linkage (Theorem~\ref{thm:link})}

\begin{lemma}[Special soundness of the equality proof]
\label{lem:eq-ss}
Two accepting transcripts $((w_a,w_b,u),c,z_a,z_b)$ and
$((w_a,w_b,u),c',z_a',z_b')$ with $c\neq c'$ yield relaxed openings
$(\bar r_a,f)$ of $t_a$ and $(\bar r_b,f)$ of $t_b$, $f=c-c'$, with equal
messages.
\end{lemma}
\begin{proof}
The verifier checks $B_1z_a=w_a+ct_{1,a}$, $B_1z_b=w_b+ct_{1,b}$ and
$\langle b_2,z_a-z_b\rangle=u+c(t_{2,a}-t_{2,b})$
(\texttt{equality\_proof.py:135--145}). Subtracting the two transcripts gives
$B_1\bar r_a=ft_{1,a}$, $B_1\bar r_b=ft_{1,b}$ and
$\langle b_2,\bar r_a-\bar r_b\rangle=f(t_{2,a}-t_{2,b})$ with
$\bar r_\ast=z_\ast-z_\ast'$, $\norm{\bar r_\ast}_\infty\le2B$. Then
$\bar m_a-\bar m_b=(t_{2,a}-t_{2,b})-f^{-1}\langle b_2,\bar r_a-\bar r_b\rangle=0$.
\end{proof}

\emph{Soundness.} Simulate $\mathcal O_{\mathsf{Link}}$ ($\epszk$). If the
verifier's $H_{\mathsf{eq}}$ point was programmed for the same context and
commitments, then $c$ is the programmed one and $(z_a,z_b)\neq(z_a^{(i)},z_b^{(i)})$
with $B_1(z_a-z_a^{(i)})=B_1(z_b-z_b^{(i)})=0$, an M-SIS solution of norm $2B$.
Otherwise Lemma~\ref{lem:fork}, with the forked oracle $H_{\mathsf{eq}}$ and
Lemma~\ref{lem:eq-ss} in place of the opening relation, extracts openings with
a common message $\bar m$. The adversary's first-run openings have messages
$m_a\neq m_b$, so $m_a\neq\bar m$ or $m_b\neq\bar m$, and
Theorem~\ref{thm:binding} applied to $t_a$ or $t_b$ gives an M-SIS solution.
Hence $\varepsilon_{\mathsf{ext}}\le\Adv^{\MSIS}$ in the bound of
Theorem~\ref{thm:se}.

\emph{Hiding.} Simulate the proof ($\epszk(Q_P+1,Q_H)$), then replace $t_a$ and
$t_b$ by uniform pairs, one M-LWE hybrid each (Lemma~\ref{lem:mlwe-instance};
the randomness is independent). The final game is independent of $b$; doing
the same from $b=1$ gives $2\epszk+4\Adv^{\MLWE}$.

\emph{Disclosure by proof.} The equality proof with $t_b=(0,\mathsf{enc}(p))$
and $r_b=0$ (a ternary opening of $t_b$) proves that $t_a$ hides $p$ without
revealing $r_a$: extraction gives a relaxed opening $(\bar r_b,f)$ of $t_b$
with $B_1\bar r_b=0$, so $\bar r_b=0$ unless M-SIS is broken, and then
$\bar m_b=\mathsf{enc}(p)=\bar m_a$. This is the fix proposed for item~D3.
\qed

\subsection{Proof of Lemma~\ref{lem:canon} and Theorem~\ref{thm:integrity}}

\emph{Unique decodability.} For a carrier segment the hashed string is the
domain tag followed, for each NAL unit, by \texttt{"NAL\textbackslash0"} and a
fixed 12-byte header, then either \texttt{"RAW\textbackslash0"}, a 64-bit
length and the RBSP, or a sequence of \texttt{"BITS\textbackslash0"}, a 64-bit
bit count and the packed bits, alternating with \texttt{"BLOCK\textbackslash0"},
the 16-byte key and length, and the signed 64-bit levels
(\texttt{video\_canonicalization.py:234--307}). The four tags differ in their
first two bytes, every field is fixed-width or length-prefixed, and the
padding of packed bits is determined by the bit count, so a left-to-right
parser recovers the whole sequence. The segment combiner hashes a tag, the
32-bit frames-per-segment and segment count, and fixed-width
$(i,f_i,d_i)$ triples (\texttt{blind\_sync\_streaming.py:225--233}); $V$ hashes a
tag and $D$. A difference in $\Canon$ is therefore a difference in the input
of one of the three SHA-256 layers whose outputs are equal, which is a
collision. Carrier segments are exactly those containing positions of $P$, so
for equal $P$ both sides hash each segment in the same mode.

\emph{What is determined.} Fix $\Canon$ and $P$. Segments without carriers
are fixed byte for byte. In a carrier segment, every NAL header field in the
serialization, every RBSP of a NAL unit without carriers, and every RBSP bit
outside carrier blocks up to the stop bit are fixed. Each carrier block is
fixed as a level vector except at carriers, where only the normalized value
$\nu=\mathrm{sgn}(\ell)2\lfloor|\ell|/2\rfloor$ is fixed. The pipeline enforces
the stable rule (\texttt{blind\_sync\_streaming.py:268--271}; carriers have
$|\ell|\ge4$ and the embedder flips the magnitude parity), so
$\ell=\nu+\mathrm{sgn}(\nu)b$ with $b$ the carrier bit; for the optional
trailing-one sign carrier, $\ell=\pm1$ with the sign as the bit. The carrier
bits are written in the order of $P$ as the chunk-framed encoding of $E$.
If that encoding is unique, $E$ fixes the bits and hence every carrier level.

\emph{Theorem~\ref{thm:integrity}.} Let $(x_h,t,\pi_h)\in L$ be the honest
triple. If $(x^*,t,\pi^*)\notin L$, the event is bounded as in
Theorem~\ref{thm:auth}, case $W_b$, with the $Q_P$ honest commitments as the
Lemma~\ref{lem:row} instances. Otherwise $x^*=x_h$ and $\pi^*=\pi_h$ (each
honest in-band commitment occurs in one triple except with negligible
probability; a credential occurs in several, all honest). Equal statements
give equal $|P|$ and $\mathrm{SHA256}(P)$, hence $P^*=P_h$ unless SHA-256
collides (the positions are hashed as canonical JSON of integer triples), and
equal $V$, hence equal $\Canon$ by the first part. The envelope is
$E^*=\texttt{magic}\|1\|\pid\|N_s\|t\|\pi^*$, and its parser accepts only this
exact layout (\texttt{video\_proof.py:142--153, 218--228}), so $E^*=E_h$. \qed

\subsection{Completeness}

An accepted attempt satisfies $B_1z-ct_1=B_1y+cB_1r-cB_1r=w$, and for the
equality proof $\langle b_2,z_a-z_b\rangle-c(t_{2,a}-t_{2,b})=\langle b_2,y_a-y_b\rangle$
because $t_{2,a}-t_{2,b}=\langle b_2,r_a-r_b\rangle$ for equal messages. By
Lemma~\ref{lem:hvzk} the number of attempts is geometric with parameter
$p_{\mathrm{acc}}$ (mean $5.23$) for the opening proof and $p_{\mathrm{acc}}^2$
(mean $27.4$) for the equality proof. With the caps $\Amax=512$
(\texttt{opening\_proof.py:33}) and $\Amax=4\,096$
(\texttt{equality\_proof.py:50}) the prover fails with probability
$2^{-156.7}$ and $2^{-220.0}$ respectively. Under the shared cap of 512 that
the reviewed code used, the equality proof failed with probability
$2^{-35.0}$ for \texttt{bdlop-open-128-v1} (means $4.65$ and $21.6$); this is
item~D4.

\section{Specification and implementation cross-check}
\label{app:code-check}

Each item lists the property relied on above and where the implementation
realizes it, at commit \texttt{ea96d60}. Paths are relative to
\texttt{VideoLevel/src}; \texttt{opening\_proof.py}, \texttt{equality\_proof.py},
\texttt{bdlop.py}, \texttt{params.py}, \texttt{ring.py}, \texttt{codec.py},
\texttt{sampling.py}, \texttt{video\_proof.py}, \texttt{video\_pipeline.py}
and \texttt{cli.py} are in \texttt{lattice\_zk/}.

\begin{enumerate}
\item \emph{Statement.} Field order and widths: \texttt{u16} version,
  $\mathit{rid}$, $\pid$, root, \texttt{u64} epoch, $N_s$, codec policy,
  carrier profile, \texttt{u64} carrier count, positions hash, $V$; every
  digest checked to be 32 bytes (\texttt{video\_proof.py:58--77}).
  Matches the specification, \S3.6. \textbf{Consistent.}
\item \emph{Fiat--Shamir input.} SHAKE-256 over the tag, $\pid$,
  \texttt{u64}$(|x|)\|x$, the full commitment (6\,720 bytes) and $w$ as
  big-endian 64-bit words; 32-byte output
  (\texttt{opening\_proof.py:42--50}). Only $x$ is length-prefixed; the other
  fields have lengths fixed by $\pid$. Equality proof: same with both
  commitments and $(w_a,w_b,u)$ (\texttt{equality\_proof.py:40--51}).
  \textbf{Consistent}; the specification's ``with length prefixes'' (\S4.7)
  was imprecise and has been corrected.
\item \emph{Hedging.} $K=\mathrm{SHA3}(e\|\text{length-prefixed }x,t,r)$
  (\texttt{opening\_proof.py:53--63,99--100}). \textbf{Consistent} (see D7).
\item \emph{Samplers.} Masks: exact uniform on $[-\gamma+1,\gamma]$ (18-bit
  draws below $2\gamma=2^{18}$, \texttt{sampling.py:60--66}); challenges:
  sign bits then Fisher--Yates with exact rejection, as in ML-DSA
  \textsf{SampleInBall} (\texttt{:69--81}); ternary randomness by rejection on
  2-bit draws (\texttt{:54--57}). \textbf{Consistent} with the uniformity used in
  Lemmas~\ref{lem:hvzk} and~\ref{lem:fork}.
\item \emph{Verification.} Exact proof length, packed responses below
  $2B+1$, zero padding, recomputed $w'$, constant-time seed comparison
  (\texttt{opening\_proof.py:118--138}). The packing has no padding bits at
  these parameters, so the proof encoding is canonical. \textbf{Consistent.}
\item \emph{Parameter checks.} Prime $q\equiv5\pmod8$, $2\sqrt2<\sqrt q$,
  $2\cdot8\kappa B<q$ (\texttt{params.py:37--54}, \texttt{ring.py:52--58});
  production pins require $|\Cset|\ge2^{256}$
  (\texttt{params.py:127--130}, \texttt{video\_pipeline.py:144--145}).
  \textbf{Consistent.}
\item \emph{Envelope.} In-band: exact length, magic \texttt{ZKLB}, version,
  pinned $\pid$, canonical commitment ($<q$ per coefficient), proof
  (\texttt{video\_proof.py:142--153}); the header session must equal the
  statement's (\texttt{:183}). Registered: \texttt{ZKLR}, header and proof only,
  pinned commitment (\texttt{:218--244}). \textbf{Consistent} with \S3.7.
\item \emph{Sessions.} One-use semantics as in the proof of
  Theorem~\ref{thm:replay}; binding $\mathsf{bind}(\mathsf{pins})$ hashes
  $\pid$, $\mathit{rid}$, root, epoch, codec policy and profile
  (\texttt{video\_pipeline.py:166--172}), so in registered mode it also
  depends on the credential through $\mathit{rid}$. Stateless mode requires
  an explicit \texttt{expected\_session} (\texttt{:262--263}).
  \textbf{Consistent.}
\item \emph{Registry.} ML-DSA-65 over a tagged canonical JSON encoding; the
  registered descriptor hashes the key and the credential
  (\texttt{video\_pipeline.py:109--131}, \texttt{zkp\_registry.py:128--147}).
  \textbf{Consistent}, but see D2.
\item \emph{Canonical digest.} As described in
  Section~\ref{sec:sec-integrity}; the specification (\S3.6) described $D$ as
  a single-stream digest and has been corrected to the segmented form.
\item \emph{Domain tags.} SHAKE-256 tags \texttt{.../bdlop-open/v1/\{key,
  randomness, fiat-shamir, mask\}\textbackslash0},
  \texttt{.../bdlop-equal/v1/\{fiat-shamir, mask\}\textbackslash0} and
  \texttt{.../challenge-poly/v1\textbackslash0} each end in the only NUL byte
  and are pairwise distinct, hence prefix-free. \textbf{Consistent.}
\end{enumerate}

\subsection{Discrepancies and gaps}

\begin{description}
\item[D1 (medium): carrier bits are malleable with a single chunk.]
  \texttt{blind\_payload\_chunks.py:12,84--137} accept a one-chunk set for any
  \texttt{chunk\_size} in $[|E|,65\,535]$, and the verifier never compares the
  extracted chunking with the pinned \texttt{profile.chunk\_bytes}
  (\texttt{video\_pipeline.py:} \texttt{267--303}). When the profile's chunk
  size is at least the envelope length, an adversary can rewrite the 16-bit
  field and the outer CRC-32: up to 48 carrier coefficients change by $\pm1$
  while $E$, $P$ and $D$ are unchanged, and the video is accepted. Rewriting
  8\,000 to 20\,000 changed 21 carrier bits in a direct test of the framing
  functions. This contradicts goal (iii) of the threat model (``any change to
  the carriers is rejected''). Fix: the verifier recomputes
  $\mathsf{Enc}_{\mathsf{prof}}(E)$, the chunks of $E$ under the pinned chunk
  size, and requires it to equal the extracted pieces (or requires the
  header field to equal the pinned size); add a negative test.
\item[D2 (medium): no holder identity in the signed registry.] The registered
  descriptor (\texttt{video\_pipeline.py:109--131}; field set
  \texttt{zkp\_registry.py:29--35}) binds a credential commitment, not a
  holder. ``Only the credential holder can produce the proof'' holds per
  commitment; which person that is is not covered by the ML-DSA signature.
  Fix: add a holder identifier (or its digest) to the registered descriptor.
\item[D3 (medium): disclosure hands over the commitment.]
  \texttt{verify\_disclosure} takes the randomness
  (\texttt{video\_proof.py:84--89,188--197}). The recipient can then produce
  fresh opening proofs for that commitment on new statements, i.e.\ accepted
  in-band videos with the creator's commitment, or impersonate a holder if a
  credential is disclosed. Fix: disclose $p$ with the equality proof against
  $(0,\mathsf{enc}(p))$ and zero randomness, bound to a disclosure context
  (proof above).
\item[D4 (low): equality-proof completeness $2^{-35}$.]
  \texttt{equality\_proof.py:29,83} reuses $\Amax=512$ with acceptance
  $p_{\mathrm{acc}}^2$ (figures for \texttt{bdlop-open-128-v1}, the set
  reviewed; $2^{-27.5}$ for \texttt{bdlop-open-128-v2}). Fix: a separate cap,
  e.g.\ 4\,096 attempts ($\approx2^{-280}$ for v1, $2^{-220}$ for v2); a larger cap does not affect zero knowledge because
  acceptance is witness-independent.
\item[D5 (low): scope of the canonical digest.] The digest is over
  ffmpeg-segmented files (\texttt{blind\_sync\_streaming.py:117--215}); segments
  without carriers are hashed as raw files without a domain tag
  (\texttt{:236--241,358,549}); RBSP bits after the stop bit of a carrier-bearing
  IDR NAL unit are excluded (\texttt{video\_canonicalization.py:}
  \texttt{279, 303, 310--317});
  carrier blocks are hashed as level vectors; the codec policy string
  (\texttt{video\_pipeline.py:25,98--99}) does not pin the segmenter version.
  None breaks Theorem~\ref{thm:integrity}, but each widens the set of
  unbound changes listed in Section~\ref{sec:sec-integrity}. Fix: tag the
  raw-segment digests, hash the full RBSP of carrier NAL units, and pin or
  replace the segmenter.
\item[D6 (low): API does not tie pins to the credential.]
  \texttt{verify\_video\_proof} accepts \texttt{pins} and \texttt{credential}
  separately (\texttt{video\_pipeline.py:250--265}) and does not check that
  \texttt{pins.relation\_id} is the relation of that credential. The CLI derives
  both together (\texttt{cli.py:98--107}), but a caller that pins one
  credential and passes another verifies an unregistered credential. Fix:
  recompute the descriptor of \texttt{credential} and compare relation
  identifiers inside \texttt{verify\_video\_proof}.
\item[D7 (info): hedging hash has no domain tag.]
  \texttt{opening\_proof.py:59} hashes $e\|\dots$ directly. Separation from the
  tagged SHA3-256 uses relies on $e$ being uniform; add a tag.
\item[D8 (info): unreduced commitments.] \texttt{Commitment} objects built
  without \texttt{from\_bytes} may hold coefficients in $[q,2^{30})$, which
  \texttt{pack\_uints} accepts (\texttt{codec.py:10--16}). Envelopes are
  unaffected (their decoder checks $<q$); pinned credentials loaded by other
  code paths are. Fix: range check in \texttt{Commitment}.
\item[D9 (info): one global key.] $\mathsf{ck}$ is expanded from a fixed public
  seed (\texttt{params.py:133--149}); all users share one M-SIS/M-LWE
  instance, so precomputation against it amortizes.
\item[D10 (info):] \texttt{verify\_statement\_binding} hard-codes the
  \texttt{lazer-v1} suite (\texttt{zkp\_registry.py:206--210}) and cannot serve
  \texttt{bdlop-open-v1}; the BDLOP path uses \texttt{pins\_from\_registry}.
\item[D11 (info): specification wording.] \S4.5 wrote the bijection as
  $(t_1,t_2)\mapsto(t_1-B_{1,2}'t_2,t_2)$; the correct map uses $t_2-m$
  (Lemma~\ref{lem:mlwe-instance}). Corrected. \S4.2 cites an ``ML-DSA
  specification, Appendix~B'' for the zero-knowledge argument; this reference
  should be checked.
\end{description}

\paragraph{Status after the review.} D1--D4 and D6 are fixed in the code
evaluated in Section~\ref{sec:evaluation}, each with a negative test: the
verifier compares the extracted chunks with the canonical chunking of $E$
under the pinned chunk size (D1); the registered descriptor can bind a
holder-identifier digest (D2); disclosure is the zero-randomness equality
proof of Section~\ref{sec:sec-soundness} and the randomness-revealing API was
removed (D3); the equality proof has its own 4\,096-attempt cap (D4); and
\texttt{verify\_video\_proof} recomputes the relation identifier of the
supplied parameters, profile, credential and holder identifier and refuses
mismatching pins (D6). D5, D7 and D8 and the reviewer question on
production-strength enforcement are closed in a second round, again each
with negative tests (\texttt{src/runtest/test\_hardening\_d5\_d7\_d8.py}).
D7: the hedged mask key hashes a domain tag, the protocol label
(opening, equality or binary proof) and the parameter identifier before the
entropy. Only \texttt{bdlop-open-128-v1} keeps the untagged derivation, so
its known-answer vectors are unchanged; the \texttt{bdlop-open-128-v2} and
binary-proof vectors were regenerated (commitments unchanged). The verifier
never recomputes a mask, so proofs from either derivation verify, which the
tests check. D8: a \texttt{Commitment} carries its parameter set and its
constructor rejects coefficients outside $[0,q)$ and wrong shapes, storing
read-only copies; verifiers still re-check objects built around the
constructor. Production strength: every application-level verifier
(\texttt{verify\_video\_proof}, the envelope, disclosure and linkage
verifiers, and \texttt{pins\_from\_registry}) refuses a parameter set with a
challenge space below $2^{256}$ unless the caller sets
\texttt{allow\_toy\_parameters}; hand-built pins with the toy set are refused.
D5 is closed by a new, opt-in codec policy (v2) with its own identifier, so
it changes profile, policy, relation and statement identifiers, while codec
policy v1 stays byte-identical for existing videos (recorded identifiers and
digests are tested). Under v2 (a) a segment without carriers is hashed with a
domain tag and its length and every segment digest is typed in the combined
digest; (b) a carrier-bearing IDR NAL unit is hashed over its complete RBSP
with its length, so zero bits after the stop bit are bound (under v1 a zero
byte appended there is accepted; under v2 it is rejected, through the real
H.264 transport); (c) the segmentation is pinned: the segments must be
exactly the access-unit partition of the input, compared byte for byte, so
the digest is a function of the input alone and a segmenter that cuts
elsewhere fails closed. A registered \texttt{bdlop-open-128-v2} run on a
200-frame CIF cut under v2 verified from the video alone, and the replayed
session, a one-bit change and the zero-extended carrier NAL unit were
rejected (\texttt{benchmark/results/q2/hardening\_e2e\_container\_cif\_200f\_v2.json}).
The ffmpeg version is not hashed, because that
would bind the verifier's installation rather than the video. Carrier blocks
are still hashed as level vectors, which is what lets the proof bind the
video without depending on its own bits. D9--D11 remain as described.
. Requires amsmath, amssymb, amsthm and the
% theorem environments definition, lemma, theorem defined in main.tex. Extra
% bibliography entries are in sections/security_refs.bib
% (\bibliography{references,sections/security_refs}).
%
% Every code reference is to commit ea96d60 of feature/lattice-pq-stack.

\providecommand{\Rq}{\mathcal{R}_q}
\providecommand{\norm}[1]{\lVert #1 \rVert}
\providecommand{\Adv}{\mathsf{Adv}}
\providecommand{\pid}{\mathit{pid}}
\providecommand{\Cset}{\mathcal{C}}
\providecommand{\Cdiff}{\bar{\mathcal{C}}}
\providecommand{\Com}{\mathsf{Com}}
\providecommand{\vid}{\mathsf{vid}}
\providecommand{\Hfs}{H_{\mathsf{fs}}}
\providecommand{\Stmt}{\mathsf{Stmt}}
\providecommand{\Canon}{\mathsf{Canon}}
\providecommand{\MSIS}{\textsf{M-SIS}}
\providecommand{\MLWE}{\textsf{M-LWE}}
\providecommand{\epszk}{\varepsilon_{\mathsf{zk}}}
\providecommand{\Amax}{A_{\max}}

\section{Mô hình và phân tích an toàn}
\label{sec:security-model}

Mục này phát biểu hệ thống chứng minh được điều gì và dưới những giả thiết nào.
Mọi kết quả đều nằm trong mô hình oracle ngẫu nhiên cổ điển (ROM). Không có định
lý nào trong mô hình oracle ngẫu nhiên lượng tử (QROM) được khẳng định
(Mục~\ref{sec:sec-qrom}). Các chứng minh đầy đủ nằm trong
Phụ lục~\ref{app:security-proofs}; việc kiểm tra chéo từng định nghĩa đối chiếu
với phần hiện thực nằm trong Phụ lục~\ref{app:code-check}.

\subsection{Ký hiệu, tham số và giả thiết}
\label{sec:sec-notation}

Đặt $\mathcal R=\mathbb Z[X]/(X^N+1)$ và $\Rq=\mathcal R/q\mathcal R$, với
$N=256$ và số nguyên tố $q=2^{30}-35\equiv5\pmod 8$. Khóa cam kết là
$\mathsf{ck}=(B_1,b_2)$ với $B_1=[I_n\,|\,B_1']\in\Rq^{n\times k}$,
$b_2=[0^n\,|\,1\,|\,b_2']\in\Rq^{1\times k}$, $k=n+\lambda+1$, $n=6$, $\lambda=7$ (nên $k=14$).
$S_1$ là tập các đa thức có hệ số trong $\{-1,0,1\}$. Với một thông điệp
$m\in\Rq$ và $r\in S_1^k$,
$\Com(m;r)=(t_1,t_2)=(B_1r,\ \langle b_2,r\rangle+m)$. Tập thách thức là
$\Cset=\{c\in\mathcal R:\norm{c}_\infty=1,\ \norm{c}_1=\kappa\}$ với
$\kappa=60$, nên $|\Cset|=\binom{256}{60}2^{60}=2^{257.01}$; tập hiệu của nó
là $\Cdiff=\{c-c':c\neq c'\in\Cset\}$. Các mặt nạ đều (uniform) trên
$Y=[-\gamma+1,\gamma]^{kN}$ với $\gamma=2^{17}$; các đáp ứng được chấp nhận khi
và chỉ khi chúng nằm trong $Z=[-B,B]^{kN}$, $B=\gamma-\kappa-1$. Chúng tôi dùng
$\norm{fv}_\infty\le\norm{f}_1\norm{v}_\infty$ xuyên suốt. Các đại lượng cụ thể
dùng trong các cận được tập hợp trong Bảng~\ref{tab:sec-numbers}.

\begin{table}[t]
\centering
\caption{Các đại lượng cụ thể cho \texttt{bdlop-open-128-v2}}
\label{tab:sec-numbers}
\footnotesize
\begin{tabular}{@{}lr@{}}
\toprule
Quantity & Value \\
\midrule
$|\Cset|$ & $2^{257.01}$\\
Fiat--Shamir seed length & 256 bits\\
$p_{\mathrm{acc}}=((2B+1)/2\gamma)^{kN}$ (opening) & $0.1912$\\
$p_{\mathrm{acc}}^2$ (equality proof) & $0.0365$\\
Attempt cap $\Amax$ (opening; equality) & $512$; $4\,096$\\
$(1-p_{\mathrm{acc}})^{512}$; $(1-p_{\mathrm{acc}}^2)^{4096}$ & $2^{-156.7}$; $2^{-220.0}$\\
$\eta=nN\log_2(2B+1)$ (min-entropy of $w$) & $27\,646.98$ bits\\
$\beta=8\kappa B$ (M-SIS bound) & $62\,885\,280\approx2^{25.91}$\\
$2\kappa+2B$ (one honest, one relaxed opening) & $262\,142\approx2^{18.0}$\\
$\sqrt{q/2}$ (invertibility threshold) & $23\,170.5$\\
$3^{kN}q^{-(n+1)N}$ & $\approx2^{-48\,079}$\\
Statement $x$ & 274 bytes\\
Envelope, in-band / registered & 14\,885 / 8\,165 bytes\\
\bottomrule
\end{tabular}
\end{table}

\begin{lemma}[Hiệu của các thách thức là khả nghịch]
\label{lem:invertible}
Mọi $f\in\Cdiff$ đều khả nghịch trong $\Rq$.
\end{lemma}
\begin{proof}
Vì $q\equiv5\pmod8$, $X^N+1$ tách theo modulo $q$ thành đúng hai nhân tử bất
khả quy, và mọi $y\in\Rq$ khác không với $\norm{y}_\infty<q^{1/2}/\sqrt2$ đều
khả nghịch (Lyubashevsky--Seiler~\cite{ref:ls18}, Hệ quả~1.2 với $k=2$). Mọi
$f\in\Cdiff$ đều khác không với $\norm{f}_\infty\le2<23\,170$. Hàm khởi tạo của
bộ tham số bắt buộc cả hai điều kiện (\texttt{params.py:49--51},
\texttt{ring.py:52--58}).
\end{proof}

\begin{definition}[Các phần mở]
\label{def:relaxed}
Một \emph{phần mở nới lỏng} của $t=(t_1,t_2)$ là một cặp $(\bar r,f)$ với
$f\in\Cdiff$, $\norm{\bar r}_\infty\le2B$ và $ft_1=B_1\bar r$. Nó xác định thông
điệp $\bar m=t_2-f^{-1}\langle b_2,\bar r\rangle$, là $\bar m$ duy nhất thỏa
$ft_2=\langle b_2,\bar r\rangle+f\bar m$ (Bổ đề~\ref{lem:invertible}). Một
\emph{phần mở trung thực} là $(r,1)$ với $r\in S_1^k$ và $t_1=B_1r$, có thông
điệp $t_2-\langle b_2,r\rangle$; đó là thứ mà một lần tiết lộ bộc lộ ra.
\end{definition}

\paragraph{Các giả thiết}
$\Adv^{\MSIS}$ là lợi thế của việc tìm $x\in\mathcal R^k$ với
$B_1x\equiv0\pmod q$, $0<\norm{x}_\infty\le\beta=8\kappa B$, cho $B_1'$ đều.
$\Adv^{\MLWE}$ là lợi thế của việc phân biệt $(A,As+e)$ với đều, cho
$A\in\Rq^{(n+1)\times\lambda}$ đều và $s\in S_1^\lambda$, $e\in S_1^{n+1}$.
Đây là hai thể hiện được ước lượng ở 140/127 và 164/150 bit (core-SVP cổ
điển/lượng tử, Bảng~\ref{tab:params}); ước lượng M-SIS lượng tử chỉ vừa dưới
128 bit. Mọi nghiệm M-SIS trong mục này đều có chuẩn tối đa $\beta$, nên không
cần thêm giả thiết lưới nào nữa.
$\Adv^{\mathsf{cr}}_{\mathsf{SHA2}}$ và $\Adv^{\mathsf{cr}}_{\mathsf{SHA3}}$ là
các lợi thế tìm va chạm đối với SHA-256 và SHA3-256, và
$\Adv^{\mathsf{euf\text{-}cma}}_{\mathsf{ML\text{-}DSA}}$ là lợi thế EUF-CMA đối
với ML-DSA-65~\cite{ref:fips204}.

\paragraph{Các oracle ngẫu nhiên}
Chúng tôi mô hình hóa như các oracle ngẫu nhiên độc lập các thể hiện SHAKE-256
được dùng cho hạt giống thách thức Fiat--Shamir $\Hfs$ (và $H_{\mathsf{eq}}$ cho
chứng minh đẳng thức), phép mở rộng thách thức $G:\{0,1\}^{256}\to\Cset$
(\textsf{SampleInBall}), luồng mặt nạ, phép mở rộng khóa, và hàm băm phòng hộ
(hedging) SHA3-256. Các thẻ miền của chúng là không-tiền-tố đôi một
(Phụ lục~\ref{app:code-check}), nên một oracle SHAKE-256 tách thành các oracle
độc lập. $G$ đều chính xác trên $\Cset$ với một đầu vào mới: các dấu và các chỉ
số Fisher--Yates được rút bằng lấy mẫu loại bỏ chính xác. Phép mở rộng khóa
dưới dạng một oracle ngẫu nhiên làm cho $\mathsf{ck}$ đều; đó là một khóa toàn
cục duy nhất được dẫn xuất từ một hạt giống công khai (không có khóa riêng cho
từng người dùng). $Q_H$ ký hiệu tổng số truy vấn oracle ngẫu nhiên của đối thủ,
bao gồm cả những truy vấn được thực hiện thay mặt nó bởi các truy vấn kiểm
chứng, cộng thêm một.

\subsection{Cú pháp của lược đồ hợp thành}
\label{sec:sec-syntax}

Tầng vận chuyển là một tập các hàm tất định của một hồ sơ vật mang đã ghim
$\mathsf{prof}$ và một khóa đồng bộ $k_s$: $\mathsf{Place}(\vid,\ell)$ trả về
các vị trí vật mang có thứ tự $P$ cho một envelope $\ell$ byte;
$\Canon(\vid,P)$ trả về digest đã chuẩn hóa theo vật mang $D$;
$\mathsf{Write}(\vid,P,E)$ ghi envelope $E$ vào các vật mang;
$\mathsf{Read}(\vid)$ trả về $(E,P,D)$ hoặc $\bot$. \emph{Tính đúng đắn vận
chuyển} nghĩa là $\mathsf{Read}(\mathsf{Write}(\vid,P,E))=(E,P,\Canon(\vid,P))$
với $P=\mathsf{Place}(\vid,|E|)$; đó là một tính chất hàm, và việc nó thất bại
chỉ gây ra từ chối.

\begin{itemize}
\item $\mathsf{Setup}$: cố định bộ tham số, $\pid$ (một digest SHA3-256 của tất
  cả các tham số và hạt giống thiết lập) và $\mathsf{ck}$.
\item $\mathsf{RegKeyGen}$, $\mathsf{RegSign}$: bên phát hành nắm một cặp khóa
  ML-DSA-65 $(\mathit{pk}_I,\mathit{sk}_I)$ và ký một sổ đăng ký
  $\mathsf{Reg}=(e,\Delta)$: một epoch và một tập các mô tả quan hệ. Một mô tả
  băm $\pid$, khóa, module quan hệ và chính sách (chính sách codec và hồ sơ vật
  mang); trong chế độ đã đăng ký nó còn băm một cam kết chứng chỉ $t_c$.
\item $\mathsf{Pins}$: cho $\mathsf{Reg}$, $\mathit{pk}_I$, $\mathsf{prof}$,
  một ngưỡng chống quay lui $e_{\min}$ và (chế độ đã đăng ký) $t_c$, kiểm tra
  chữ ký, tính lại mô tả kỳ vọng, đòi hỏi đúng một kết quả khớp, và xuất ra
  $\mathsf{pins}=(\mathit{rid},\mathit{root},e,\mathit{cpid},
  \mathit{prof\_id})$, hoặc thất bại.
\item $\mathsf{Issue}$: rút một thách thức đều 32 byte $N_s$ và lưu nó cùng ràng
  buộc $\mathsf{bind}(\mathsf{pins})$ và một thời điểm hết hạn.
\item $\mathsf{Embed}$: trên vật phủ $\vid$, tải trọng $p$ (in-band) hoặc chứng
  chỉ $(t_c,r_c)$ (đã đăng ký), $\mathsf{pins}$ và $N_s$: tính
  $P=\mathsf{Place}(\vid,\ell)$ từ một placeholder bằng không có độ dài cuối
  cùng $\ell$, $D=\Canon(\vid,P)$ và $x=\Stmt(\mathsf{pins},N_s,P,D)$; với
  in-band, lấy mẫu $r\leftarrow S_1^k$, đặt $t=\Com(\mathsf{enc}(p);r)$; chứng
  minh $\pi\leftarrow\mathsf{P}(x,t,r)$; lập $E=\texttt{"ZKLB"}\|1\|\pid\|N_s\|t\|\pi$
  (hoặc $\texttt{"ZKLR"}\|1\|\pid\|N_s\|\pi$) và xuất ra
  $\mathsf{Write}(\vid,P,E)$ và, với in-band, phần mở $(p,r)$.
\item $\mathsf{Verify}$: trên $\vid^*$: $(E^*,P^*,D^*)=\mathsf{Read}(\vid^*)$;
  phân tích $E^*$ một cách nghiêm ngặt; lấy $N_s^*$ từ $E^*$ (chế độ có lưu
  trữ) hoặc dùng một $N_s$ cố định (chế độ không trạng thái);
  $x^*=\Stmt(\mathsf{pins},N_s^*,P^*,D^*)$; chấp nhận khi và chỉ khi
  $\mathsf{V}(x^*,t^*,\pi^*)=1$ với $t^*$ từ $E^*$ (in-band) hoặc $t_c$ đã ghim
  (đã đăng ký) và, trong chế độ có lưu trữ, $N_s^*$ còn sống đối với
  $\mathsf{bind}(\mathsf{pins})$ và được tiêu thụ một cách nguyên tử.
\item $\mathsf{VerDisc}$: chấp nhận $(E,p,r)$ khi và chỉ khi $(r,1)$ là một phần
  mở trung thực của cam kết trong $E$ tới $\mathsf{enc}(p)$.
\item $\mathsf{Link}$, $\mathsf{VerLink}$: chứng minh đẳng thức
  $\Pi_{\mathsf{eq}}$ trên hai cam kết, được ràng buộc vào một digest ngữ cảnh
  $\mathit{lx}=\mathrm{SHA3}(\text{tag}\|\pid\|\text{ngữ cảnh có tiền tố độ dài})$.
\end{itemize}

Phát biểu là chuỗi 274 byte
\begin{multline*}
\Stmt=\mathtt{u16}(1)\|\mathit{rid}\|\pid\|\mathit{root}\|\mathtt{u64}(e)\|N_s\|
\mathit{cpid}\|\mathit{prof\_id}\\
\|\mathtt{u64}(|P|)\|\mathrm{SHA256}(P)\|V,\quad
V=\mathrm{SHA256}(\text{tag}\|D),
\end{multline*}
với $\mathrm{SHA256}(P)$ được lấy trên một mã hóa JSON chuẩn tắc của các vị trí
có thứ tự. Chứng minh mở $\Pi_{\mathsf{open}}$ xuất ra $\pi=(h,z)$ với
$h=\Hfs(\text{tag}\|\pid\|\mathtt{u64}(|x|)\|x\|\mathsf{enc}(t)\|\mathsf{enc}(w))$,
$w=B_1y$, $c=G(h)$, $z=y+cr\in Z$; bên kiểm chứng tính lại $w'=B_1z-ct_1$ và
kiểm tra $h$. Mọi trường sau $x$ đều có độ dài được cố định bởi $\pid$, nên đầu
vào của hàm băm giải mã được một cách duy nhất.

\subsection{Mô hình đối thủ}
\label{sec:sec-adversary}

Đối thủ $\mathcal A$ là một thuật toán xác suất cổ điển. Nó kiểm soát mọi kênh,
chọn các vật phủ và các tải trọng, biết $k_s$ (không tính chất nào dưới đây phụ
thuộc vào tính bí mật của nó), và có các oracle sau. $\mathcal O_{\mathsf{Iss}}$
trả về một $N_s$ mới. $\mathcal O_{\mathsf{Emb}}(\vid,p,N_s)$ chạy một
$\mathsf{Embed}$ in-band trung thực và trả về video stego;
$\mathcal O_{\mathsf{EmbR}}(j,\vid,N_s)$ làm điều tương tự cho chứng chỉ đã đăng
ký của bên nắm giữ trung thực $j$. Chúng tôi cho phép các oracle này chứng minh
các phát biểu $x$ tùy ý do $\mathcal A$ chọn, điều này chỉ làm mạnh thêm các kết
quả. Các đầu ra của chúng tạo thành danh sách $L$ các bộ ba $(x,t,\pi)$.
$\mathcal O_{\mathsf{Ver}}$ chạy $\mathsf{Verify}$ với kho dùng chung.
$\mathcal O_{\mathsf{Link}}$ trả về các chứng minh liên kết trung thực giữa các
cam kết đã ghi. $\mathcal O_{\mathsf{Disc}}(i)$ trả về phần mở của cam kết trung
thực thứ $i$ và đánh dấu nó là \emph{đã tiết lộ}.
$\mathcal O_{\mathsf{Reg}}(e,\Delta,T)$ khiến bên phát hành ký một sổ đăng ký
với epoch $e$ mà các mô tả của nó là các mô tả in-band $\Delta$ và các mô tả đã
đăng ký mà bên phát hành tính từ các cam kết chứng chỉ trong $T$, tất cả do
$\mathcal A$ chọn; $T_{\mathsf{reg}}$ là hợp của mọi $T$ đã truy vấn. (Một bên
phát hành chỉ ký một mô tả đã đăng ký cho một chứng chỉ mà nó đã thẩm định, nên
các mô tả đã đăng ký chỉ đi vào qua $T$.) $Q_P$ chặn số chứng minh được tạo bởi
các bên trung thực.

Bên phát hành và bên kiểm chứng là trung thực; kho phiên là tuyến tính hóa được
(linearizable) và bền vững. Việc làm hỏng bên nắm giữ là tĩnh (đối thủ đăng ký
các chứng chỉ của riêng nó qua $\mathcal O_{\mathsf{Reg}}$); tiết lộ là hình
thức làm hỏng thích nghi duy nhất, và các cam kết đã tiết lộ bị loại khỏi các
điều kiện thắng. Các kênh phụ, lỗi codec, và tính đúng đắn của tầng vận chuyển
nằm ngoài phạm vi.

\subsection{Tính không tiết lộ tri thức của chứng minh Fiat--Shamir-with-aborts}
\label{sec:sec-zk}

\begin{lemma}[Chấp nhận độc lập với nhân chứng]
\label{lem:hvzk}
Với mọi $r\in S_1^k$ và $c\in\Cset$, một lần thử với $y\leftarrow Y$ được chấp
nhận với xác suất đúng bằng $p_{\mathrm{acc}}$, và với điều kiện được chấp nhận
thì $z=y+cr$ đều trên $Z$.
\end{lemma}
\noindent
Mỗi hệ số $v$ của $cr$ thỏa $|v|\le\kappa$, nên $2B+1$ giá trị
$z_i\in[-B,B]$ tương ứng với các mặt nạ $z_i-v\in[-\gamma+1,\gamma-1]\subset
[-\gamma+1,\gamma]$, mỗi cái có xác suất $1/2\gamma$.

\begin{lemma}[Mô phỏng]
\label{lem:sim}
Tồn tại một bộ mô phỏng mà, chỉ với $(x,t)$ và quyền lập trình $\Hfs$, xuất ra
các chứng minh có phân phối hợp với các oracle ngẫu nhiên nằm trong khoảng
\begin{multline*}
\epszk(Q_P,Q_H)=2Q_P(Q_H+Q_P)2^{-256}\\
+(\Amax+1)Q_P\bigl(Q_H+(\Amax+1)Q_P\bigr)2^{-\eta}
\end{multline*}
so với $Q_P$ chứng minh thực, với mọi đối thủ thực hiện $Q_H$ truy vấn. Điều
tương tự đúng cho $\Pi_{\mathsf{eq}}$ với $p_{\mathrm{acc}}^2$ thay cho $p_{\mathrm{acc}}$.
\end{lemma}
\noindent
Bộ mô phỏng lấy mẫu số lần thử từ phân phối hình học bị cắt cụt với tham số
$p_{\mathrm{acc}}$, xuất ra $\bot$ với xác suất $(1-p_{\mathrm{acc}})^{\Amax}$,
và ngược lại lấy mẫu $h\leftarrow\{0,1\}^{256}$, $z\leftarrow Z$, đặt
$w=B_1z-G(h)t_1$ và lập trình $\Hfs$ tại $(\pid,x,t,w)$. Bước tế nhị là các truy
vấn oracle ngẫu nhiên mà bên chứng minh thực thực hiện trong các lần thử bị loại
và không bao giờ tiết lộ: các phân tích trước đây về Fiat--Shamir with aborts đã
xử lý sai việc tái lập trình tương ứng trong bước CMA-to-NMA, như được Barbosa
và cộng sự~\cite{ref:barbosa} chỉ ra và sửa chữa, và được Devevey và cộng
sự~\cite{ref:devevey} phân tích chi tiết. Ở đây bước này được giải quyết trong
ROM cổ điển nhờ hai sự kiện đặc thù của hiện thực này: sự chấp nhận là
\emph{hoàn toàn} độc lập với nhân chứng và với $c$ (Bổ đề~\ref{lem:hvzk}), và
mọi điểm truy vấn ẩn đều chứa $w=B_1y$ có min-entropy ít nhất $\eta$ (bởi vì
$B_1=[I_n|B_1']$). Tính độc lập chính xác cũng có nghĩa là số lần thử không rò
rỉ gì, dù rằng ngoài ra phần hiện thực không chạy trong thời gian hằng.

\subsection{Tính đúng đắn tri thức và ràng buộc vào video}
\label{sec:sec-soundness}

\begin{definition}[Ràng buộc video trích xuất được qua mô phỏng]
\label{def:se}
Trong thí nghiệm $\mathsf{SE}$, $\mathcal A$ tương tác với mọi oracle và xuất
ra $\vid^*$. Gọi $x^*$ là phát biểu mà bên kiểm chứng tính lại từ $\vid^*$ và
các mục ghim của chính nó, và $(t^*,\pi^*)$ là cam kết và chứng minh mà nó kiểm
chứng. $\mathcal A$ \emph{thắng} nếu $\mathsf{Verify}$ chấp nhận và
$(x^*,t^*,\pi^*)\notin L$. Một bộ trích xuất $\mathcal E^{\mathcal A}$
\emph{thành công} nếu nó xuất ra $x^*$ và một phần mở nới lỏng của $t^*$, trong
đó $x^*$ là phát biểu của bên kiểm chứng trong một lần chạy thắng của $\mathcal A$.
\end{definition}

\begin{theorem}[Tính đúng đắn tri thức ràng buộc vào phát biểu]
\label{thm:se}
Giả sử $\mathcal A$ thắng $\mathsf{SE}$ với xác suất $\varepsilon$. Tồn tại một
bộ trích xuất $\mathcal E$ chạy $\mathcal A$ hai lần với các oracle được mô
phỏng và thành công với xác suất $\varepsilon_{\mathsf{ext}}$ sao cho
\begin{multline*}
\varepsilon\le\epszk(Q_P,Q_H)+\Adv^{\MSIS}+Q_H2^{-256}\\
+\sqrt{Q_H\bigl(\varepsilon_{\mathsf{ext}}+Q_H2^{-256}+2^{-257}\bigr)}.
\end{multline*}
\end{theorem}
\noindent
\emph{Ý tưởng chứng minh.} Thay các chứng minh trung thực bằng các chứng minh
được mô phỏng (Bổ đề~\ref{lem:sim}). Nếu $\mathcal A$ tái dùng một điểm đã lập
trình, tức $(x^*,t^*)=(x_i,t_i)$ và $w^*=w_i$, thì $h^*=h_i$, do đó $c^*=c_i$ và
$B_1(z^*-z_i)=0$ với $0<\norm{z^*-z_i}_\infty\le2B$: một nghiệm M-SIS. Ngược
lại, truy vấn $(\pid,x^*,t^*,w^*)$ của bên kiểm chứng được trả lời bằng một giá
trị ngẫu nhiên mới, và bổ đề rẽ nhánh tổng quát của Bellare và
Neven~\cite{ref:bn06} (Bổ đề~1) với không gian đầu ra $2^{256}$ cho hai chứng
minh được chấp nhận có cùng $(x^*,t^*,w^*)$ và các hạt giống $h\neq h'$. Khi đó
$c=G(h)\neq G(h')=c'$ ngoại trừ với xác suất $2^{-257}+Q_H2^{-256}$, và
$B_1(z-z')=(c-c')t_1^*$ cho phần mở nới lỏng $(z-z',c-c')$. Hiệu thách thức là
khả nghịch theo Bổ đề~\ref{lem:invertible}, chính điều này khiến phần mở xác
định một thông điệp duy nhất. Bởi vì truy vấn bị rẽ nhánh chứa $x^*$, phần mở
được ràng buộc vào đúng phát biểu do bên kiểm chứng tính lại: cam kết video, các
vị trí vật mang, thách thức phiên, gốc sổ đăng ký và epoch, quan hệ và chính
sách của nó. \qed

Tổn thất căn bậc hai là cái giá thường thấy của việc tua lại (rewinding). Với
$Q_H=2^{64}$, cận biến một xác suất trích xuất $2^{-128}$ thành $\varepsilon
\lesssim2^{-32}$. Đối với tính xác thực (Định lý~\ref{thm:auth}), có một cận
chặt nhờ hướng SelfTargetMSIS của Kiltz, Lyubashevsky và
Schaffner~\cite{ref:kls18}, với cái giá là giả thiết phi chuẩn đó.

\begin{theorem}[Tính ràng buộc đối với các phần mở nới lỏng]
\label{thm:binding}
Không đối thủ nào có thể xuất ra $t$ và hai phần mở của $t$, mỗi phần mở trung
thực hoặc nới lỏng, với các thông điệp khác nhau, ngoại trừ với xác suất
$\Adv^{\MSIS}$.
\end{theorem}
\noindent
Cho $(\bar r,f,m)$ và $(\bar r',f',m')$, vector
$x=f'\bar r-f\bar r'$ thỏa $B_1x=0$; nếu $x=0$ thì
$\langle b_2,x\rangle=ff'(m-m')=0$, mâu thuẫn với tính khả nghịch. Chuẩn của nó
tối đa là $2\cdot2\kappa\cdot2B=\beta$ cho hai phần mở nới lỏng, $2\kappa+2B$
cho một trung thực và một nới lỏng, và $2$ cho hai trung thực, tất cả đều dưới
$q/2$ (\texttt{params.py:53}), nên $x\not\equiv0$ là một nghiệm M-SIS.

\noindent
Gộp lại, các Định lý~\ref{thm:se} và~\ref{thm:binding} nói rằng: ai tạo ra một
video in-band được chấp nhận đều bị ràng buộc vào thông điệp $\bar m$ mà bộ
trích xuất tính ra, và bất kỳ lần tiết lộ về sau nào mà $\mathsf{VerDisc}$ chấp
nhận đều bộc lộ một tải trọng $p$ với $\mathsf{enc}(p)=\bar m$, trừ khi M-SIS
bị phá.

\emph{Việc tiết lộ giữ cho tính ngẫu nhiên riêng tư.} Một lần tiết lộ là tải
trọng $p$ cùng với một chứng minh đẳng thức giữa $t$ và cam kết tất định
$\mathsf{Com}(\mathsf{enc}(p);0)$, được Fiat--Shamir ràng buộc vào một ngữ cảnh
do bên nhận chọn (định danh phát biểu và nonce). Bên nhận biết $p$ nhưng không
biết phần mở nào của $t$, nên nó không thể tạo các chứng minh mở mới cho $t$;
theo tính đúng đắn của chứng minh đẳng thức, một lần tiết lộ được chấp nhận vẫn
thỏa $\mathsf{enc}(p)=\bar m$ trừ khi M-SIS bị phá. Một phiên bản trước đây bộc
lộ chính $r$; đây là mục~D3 trong Phụ lục~\ref{app:code-check} và đã được sửa.
Các định lý dưới đây vẫn coi các cam kết mà \emph{phần mở} của chúng đã bị rò rỉ
(ví dụ do lộ khóa) là đã tiết lộ.

\begin{lemma}[Tính một chiều của các cam kết trung thực]
\label{lem:row}
Cho $\mathsf{ck}$ và $t^{(j)}=\Com(m_j;r_j)$, $j\le Q$, với
$r_j\leftarrow S_1^k$ và các thông điệp $m_j$ được chọn thích nghi trước khi mỗi
$t^{(j)}$ được rút, không đối thủ nào xuất ra một phần mở nới lỏng của bất kỳ
$t^{(j)}$ nào ngoại trừ với xác suất
$\Adv^{\mathsf{row}}_Q=\Adv^{\MSIS}+Q\bigl(\Adv^{\MLWE}+3^{kN}q^{-(n+1)N}\bigr)$.
\end{lemma}
\noindent
Đặt $x=fr_j-\bar r$, nên $B_1x=0$ và $\norm{x}_\infty\le2\kappa+2B$. Nếu
$x\neq0$ thì nó là một nghiệm M-SIS (phép rút gọn đã sinh mọi $t^{(j)}$ và biết
$r_j$). Nếu $x=0$ thì đối thủ đã tìm được $r_j=f^{-1}\bar r$, điều mà một bộ
phân biệt cấy mẫu M-LWE của nó tại một $j$ được đoán sẽ kiểm tra đối chiếu với
$t^{(j)}$; với $t^{(j)}$ đều, một phần mở bậc ba tồn tại với xác suất tối đa
$3^{kN}q^{-(n+1)N}$. Chỉ riêng tính ràng buộc là không đủ ở đây: một đối thủ
khôi phục được $r_j$ sẽ mở $t^{(j)}$ tới đúng thông điệp.

\subsection{Che giấu tải trọng}
\label{sec:sec-hiding}

\begin{definition}[Tính không phân biệt được tải trọng]
\label{def:hiding}
Trong $\mathsf{IND}^b$, $\mathcal A$ có mọi oracle và nộp một lần
$(\vid,p_0,p_1,N_s)$ với $|p_0|,|p_1|\le766$ byte; bên thách thức trả về
$\mathcal O_{\mathsf{Emb}}(\vid,p_b,N_s)$. $\mathcal A$ không được truy vấn
$\mathcal O_{\mathsf{Disc}}$ hoặc $\mathcal O_{\mathsf{Link}}$ trên cam kết
thách thức. $\Adv^{\mathsf{ind}}=|\Pr[\mathsf{IND}^1=1]-\Pr[\mathsf{IND}^0=1]|$.
\end{definition}

\begin{theorem}[Che giấu]
\label{thm:hiding}
$\Adv^{\mathsf{ind}}\le2\,\epszk(Q_P+1,Q_H)+2\Adv^{\MLWE}$.
\end{theorem}
\noindent
Video stego chỉ phụ thuộc vào $b$ thông qua $E$: $P$ và $D$ được cố định từ một
placeholder trước khi chứng minh, và envelope có kích thước cố định, nên độ dài
tải trọng cũng được che giấu. Mô phỏng chứng minh thách thức
(Bổ đề~\ref{lem:sim}), rồi thay $t$ bằng một cặp đều: với
$B_1'=[B_{1,2}'\,|\,B_{1,3}']$ và $r=(r_1,r_2,r_3)$,
$(t_1-B_{1,2}'(t_2-m),\,t_2-m)=A r_3+(r_1,r_2)$ với
$A=\bigl[\begin{smallmatrix}B_{1,3}'-B_{1,2}'b_2'\\ b_2'\end{smallmatrix}\bigr]$
đều, một mẫu M-LWE hạng $\lambda$ với $n+1$ mẫu vành; ánh xạ là một song ánh với
mỗi $m$ cố định. Khi đó thí nghiệm độc lập với $b$. Tính che giấu mang tính tính
toán và là tính chất duy nhất dựa trên M-LWE; nó đúng cho tải trọng, chứ không
cho sự kiện rằng một tải trọng tồn tại, vốn là một tính chất giấu tin được đánh
giá bằng thực nghiệm.

\subsection{Khả năng chống phát lại}
\label{sec:sec-replay}

\begin{definition}[Phát lại]
\label{def:replay}
Trong $\mathsf{RPL}$, $\mathcal A$ có mọi oracle đối với một bên kiểm chứng có
kho phiên và thắng nếu (i) hai lời gọi $\mathcal O_{\mathsf{Ver}}$ chấp nhận với
cùng $N_s$; hoặc (ii) một lời gọi chấp nhận với một $N_s$ không được phát hành
cho $\mathsf{bind}(\mathsf{pins})$, hoặc đã hết hạn; hoặc (iii) một lời gọi chấp
nhận $\vid^*$ có bộ ba $(x^*,t^*,\pi^*)\notin L$ và $\mathcal A$ đã truy vấn
$\Hfs$ tại điểm $(\pid,x^*,t^*,w^*)$ của bên kiểm chứng trước khi $N_s^*$ được
phát hành.
\end{definition}

\begin{theorem}[Khả năng chống phát lại]
\label{thm:replay}
$\Pr[\text{(i) hoặc (ii)}]\le\Adv^{\mathsf{cr}}_{\mathsf{SHA2}}+Q_I^22^{-257}$ và
$\Pr[\text{(iii)}]\le Q_HQ_I2^{-256}$, với $Q_I$ phiên được phát hành.
\end{theorem}
\noindent
(i) và (ii) suy ra từ ngữ nghĩa của kho: các hàng chỉ tồn tại cho các thách
thức đã phát hành (khóa theo một digest SHA-256 của $N_s$), và việc tiêu thụ là
một \texttt{UPDATE \dots\ WHERE consumed\_at IS NULL AND expires\_at>now} có
điều kiện đơn nhất bên trong một giao dịch ghi, nên nhiều nhất một bên gọi quan
sát được một hàng đã cập nhật (\texttt{zkp\_sessions.py:181--241}). Với (iii),
$N_s^*$ là một chuỗi con của $x^*$, nên một truy vấn thực hiện trước đó dự đoán
một giá trị 256 bit đều. Cùng với Định lý~\ref{thm:se}, một phiên được chấp nhận
tương ứng với một bên đã biết một phần mở và đã tính chứng minh sau khi phiên
được phát hành.

\emph{Điều mà một bên kiểm chứng không trạng thái không thể đem lại.} Với một
$N_s$ cố định (\texttt{expected\_session}), (i) thất bại: bất kỳ bản sao nào của
một video được chấp nhận đều lại được chấp nhận, bởi bất kỳ ai và bao nhiêu lần
tùy ý, vì một chứng minh không tương tác có thể chuyển nhượng. Nếu $N_s$ bị tái
dùng qua nhiều lần kiểm chứng, (iii) cũng thất bại và chỉ còn lại ràng buộc
video (Định lý~\ref{thm:se}). Không biến thể nào ngăn được việc tiếp chuyển
(relay): chế độ in-band không ràng buộc $N_s$ vào một bên chứng minh, nên ai
trình ra một video hợp lệ trước thì tiêu thụ phiên.

\subsection{Tính xác thực chế độ đã đăng ký}
\label{sec:sec-auth}

\begin{definition}[Tính xác thực]
\label{def:auth}
Trong $\mathsf{AUTH}$, bên thách thức tạo các chứng chỉ trung thực
$t_c^{(j)}=\Com(m_j;r_j)$, $j\le Q_C$, với $m_j$ do $\mathcal A$ chọn, và
$\mathcal A$ có mọi oracle. $\mathcal A$ xuất ra $(\mathsf{Reg}^*,\vid^*,t^*)$
và thắng nếu một bên kiểm chứng ghim $\mathit{pk}_I$ và $t^*$ chấp nhận $\vid^*$
với các mục ghim từ $\mathsf{Reg}^*$, bộ ba $(x^*,t^*,\pi^*)\notin L$, và hoặc
(a) $t^*\notin T_{\mathsf{reg}}$ (chưa bao giờ được bên phát hành đăng ký), hoặc
(b) $t^*=t_c^{(j)}$ cho một $j$ trung thực chưa tiết lộ.
\end{definition}

\begin{theorem}[Tính xác thực]
\label{thm:auth}
\begin{multline*}
\Adv^{\mathsf{auth}}\le\Adv^{\mathsf{euf\text{-}cma}}_{\mathsf{ML\text{-}DSA}}
+\Adv^{\mathsf{cr}}_{\mathsf{SHA2}}+\epszk(Q_P,Q_H)\\
+\Adv^{\MSIS}+Q_H2^{-256}\\
+\sqrt{Q_H\bigl(\Adv^{\mathsf{row}}_{Q_C}+Q_H2^{-256}+2^{-257}\bigr)}.
\end{multline*}
\end{theorem}
\noindent
Trường hợp (a) cho một chữ ký giả mạo ML-DSA trên các byte ký của sổ đăng ký
hoặc, nếu $\mathsf{Reg}^*$ đã được ký, một va chạm SHA-256 giữa các đầu vào có
độ dài cố định $B_1'\|b_2'\|t^*$ và $B_1'\|b_2'\|t'$ với một
$t'\in T_{\mathsf{reg}}$ nào đó, vì mô tả khớp chứa digest đó. Trong trường hợp
(b), Định lý~\ref{thm:se} áp dụng cho trò chơi được mô phỏng trích xuất một phần
mở nới lỏng của một chứng chỉ trung thực, điều mà Bổ đề~\ref{lem:row} loại trừ.
Định lý này nói về các chứng chỉ; sự liên kết tới một con người chỉ mạnh ngang
với việc bên phát hành kiểm tra bên nắm giữ trước khi ký. Khi bên phát hành ràng
buộc một định danh của bên nắm giữ, mô tả đã ký chứa một digest có tách miền của
nó bên cạnh cam kết, nên quan hệ đã ghim của bên kiểm chứng nêu tên cả hai
(Phụ lục~\ref{app:code-check}, mục~D2, đã sửa).

\subsection{Liên kết}
\label{sec:sec-linkage}

\begin{definition}[Tính đúng đắn và tính che giấu của liên kết]
\label{def:link}
Trong $\mathsf{LNK\text{-}S}$, $\mathcal A$ xuất ra một ngữ cảnh, $t_a,t_b$, một
$\pi_{\mathsf{eq}}$ được chấp nhận không do $\mathcal O_{\mathsf{Link}}$ xuất ra
cho bộ ba đó, và các phần mở (trung thực hoặc nới lỏng) của $t_a$ và $t_b$ với
các thông điệp khác nhau. Trong $\mathsf{LNK\text{-}H}^b$, $\mathcal A$ nộp
$(p_0,p_1,\mathit{ctx})$ và nhận $t_a=\Com(p_b;r_a)$,
$t_b=\Com(p_b;r_b)$ và $\pi_{\mathsf{eq}}$ cho $\mathit{ctx}$.
\end{definition}

\begin{theorem}[Liên kết]
\label{thm:link}
$\Adv^{\mathsf{lnk\text{-}s}}$ tuân theo cận của Định lý~\ref{thm:se} với
$\varepsilon_{\mathsf{ext}}\le\Adv^{\MSIS}$, và
$\Adv^{\mathsf{lnk\text{-}h}}\le2\,\epszk(Q_P+1,Q_H)+4\Adv^{\MLWE}$.
\end{theorem}
\noindent
Hai bản ghi bị rẽ nhánh với một thông báo $(w_a,w_b,u)$ cho các phần mở nới lỏng
$(\bar r_a,f)$, $(\bar r_b,f)$ với một $f$ \emph{chung} và
$f(t_{2,a}-t_{2,b})=\langle b_2,\bar r_a-\bar r_b\rangle$, do đó
$\bar m_a=\bar m_b$; hai phần mở của đối thủ khi đó mâu thuẫn với
Định lý~\ref{thm:binding} cho $t_a$ hoặc cho $t_b$. Việc liên kết tới một chứng
chỉ đã đăng ký còn xác thực thêm bên nắm giữ của nó, như trong
Định lý~\ref{thm:auth}, bởi vì việc trích xuất cũng cho ra một phần mở nới lỏng
của $t_c$. Tính che giấu dùng bộ mô phỏng cho $\Pi_{\mathsf{eq}}$ và hai bước
lai M-LWE; sự bằng nhau của hai tải trọng được bộc lộ theo thiết kế.

\subsection{Tính toàn vẹn của nội dung video}
\label{sec:sec-integrity}

$\Canon$ chia luồng bằng bộ phân đoạn đã ghim (\texttt{ffmpeg -c:v copy
-f segment}, $F$ khung mỗi đoạn) và kết hợp các digest của từng đoạn với
$F$ và các số đếm khung. Một đoạn mang các chunk envelope được băm qua một phép
tuần tự hóa có gắn thẻ, phân định bằng độ dài: với mỗi đơn vị NAL là chỉ số,
kiểu, \texttt{forbidden\_zero\_bit}, \texttt{nal\_ref\_idc} và kích thước mã
khởi đầu (start-code); RBSP nguyên văn cho các đơn vị NAL không có vật mang; và,
bên trong một đơn vị NAL IDR có vật mang, các bit RBSP giữa các khối vật mang
cho tới bit dừng (stop bit), với mỗi khối vật mang được thay bằng vector hệ số
của nó trong đó mỗi vật mang được chuẩn hóa (độ lớn làm tròn xuống số chẵn, giữ
dấu). Các đoạn khác được băm từng byte.

\begin{lemma}[Digest và envelope xác định những gì]
\label{lem:canon}
Phép tuần tự hóa giải mã được một cách duy nhất, nên $V$ bằng nhau kéo theo
$\Canon$ bằng nhau ngoại trừ với xác suất $3\Adv^{\mathsf{cr}}_{\mathsf{SHA2}}$.
Dưới quy tắc vật mang ổn định ($|\ell|\ge4$, một bit chẵn-lẻ), một mức vật mang
là $\mathrm{sgn}(\ell)(2\lfloor|\ell|/2\rfloor+b)$, nên $(\Canon,P)$ và các bit
vật mang xác định toàn bộ vector hệ số của mọi khối vật mang. Nếu chuỗi bit vật
mang là mã hóa duy nhất của $E$ dưới hồ sơ đã ghim, thì $(\Canon,P,E)$ xác định:
sự phân đoạn và các số đếm khung; mọi byte của các đoạn không có vật mang; mọi
trường header NAL và bit RBSP bên ngoài các khối vật mang; và các vector hệ số
của các khối vật mang.
\end{lemma}

\begin{theorem}[Tính toàn vẹn]
\label{thm:integrity}
Giả sử một lần nhúng trung thực tạo ra $\vid_h$ cho phiên $N_s$ với cam kết $t$
(hoặc chứng chỉ $t_c$), và $\mathcal A$ xuất ra $\vid^*$ được chấp nhận cho
$N_s$ với cùng cam kết (hoặc chứng chỉ), không tiết lộ. Khi đó
$\Canon(\vid^*)=\Canon(\vid_h)$, $P^*=P_h$ và $E^*=E_h$, ngoại trừ với xác suất
tối đa
\begin{multline*}
4\Adv^{\mathsf{cr}}_{\mathsf{SHA2}}+\epszk(Q_P,Q_H)+\Adv^{\MSIS}+Q_H2^{-256}\\
+\sqrt{Q_H\bigl(\Adv^{\mathsf{row}}_{Q_P}+Q_H2^{-256}+2^{-257}\bigr)}.
\end{multline*}
\end{theorem}
\noindent
Nếu $x^*\neq x_h$, bộ ba là mới và việc trích xuất cho ra một phần mở nới lỏng
của một cam kết trung thực. Nếu $x^*=x_h$, giá trị băm vị trí và $V$ khớp nhau,
và Bổ đề~\ref{lem:canon} được áp dụng; khi đó chứng minh chính là chứng minh
trung thực, ngược lại trường hợp điểm-đã-lập-trình của Định lý~\ref{thm:se} cho
một nghiệm M-SIS.

\paragraph{Các giới hạn của khẳng định toàn vẹn}
\begin{itemize}
\item Chế độ in-band ràng buộc nội dung vào một \emph{cam kết}, chứ không vào
  một con người: bất kỳ ai nắm một $N_s$ còn sống đều có thể nhúng lại bất kỳ
  video đã chỉnh sửa nào với một cam kết mới của riêng nó, và nó sẽ được chấp
  nhận.
\item Các vật mang chỉ được bảo vệ như các bit envelope. Định lý~\ref{thm:integrity}
  cần các bit vật mang là mã hóa duy nhất của $E$. Bên kiểm chứng chia lại chunk
  của $E$ theo kích thước chunk đã ghim và từ chối trừ khi các mảnh trích xuất
  được đúng là mã hóa đó; nếu không có phép kiểm này, một envelope một-chunk sẽ
  để trống trường \texttt{chunk\_size} 16 bit của nó (mục~D1 trong
  Phụ lục~\ref{app:code-check}, đã sửa).
\item Không được ràng buộc: các byte không sống sót qua bộ phân đoạn đã ghim (dữ
  liệu bộ chứa, các luồng khác, dữ liệu ngoài các đơn vị NAL); các bit căn lề
  RBSP sau bit dừng trong các đơn vị NAL IDR có mang vật mang; các từ mã CAVLC
  thay thế cho cùng một vector hệ số của một khối vật mang; tầng chống giả lập
  (digest bao phủ RBSP); và phiên bản của bộ phân đoạn, vốn được chính sách
  codec nêu tên nhưng không ghim.
\item Không có khẳng định về tính đúng đắn codec: digest bao phủ cú pháp, không
  bao phủ các điểm ảnh đã giải mã. Không có khẳng định về nguồn gốc: không gì
  nói rằng một máy quay đã ghi $\vid_h$ hay rằng tải trọng là đúng sự thật.
\end{itemize}

\subsection{Các đối thủ lượng tử}
\label{sec:sec-qrom}

Mọi định lý ở trên đều là các phát biểu trong ROM cổ điển. Việc tua lại
(Định lý~\ref{thm:se}) và việc lập trình thích nghi (Bổ đề~\ref{lem:sim}) không
chuyển sang QROM như đã phát biểu. Barbosa và cộng sự~\cite{ref:barbosa} đưa ra
các chứng minh QROM của bước CMA-to-NMA cho Fiat--Shamir with aborts, và Don và
cộng sự~\cite{ref:dfms19} cùng Liu và Zhandry~\cite{ref:lz19} chỉ ra rằng phép
biến đổi Fiat--Shamir bảo toàn tính đúng đắn (tri thức) của các $\Sigma$-protocol
phù hợp trước các đối thủ lượng tử với tổn thất đa thức. Chúng tôi chưa kiểm
chứng rằng các giả thiết của họ đúng cho bộ trích xuất nới lỏng và bên chứng
minh có loại bỏ được dùng ở đây, nên chúng tôi không khẳng định định lý QROM
nào. Cách diễn đạt của chúng tôi là ``hậu lượng tử một cách hợp lý'': các phép
rút gọn nhắm tới M-SIS, M-LWE, ML-DSA và va chạm băm, không cái nào có một tấn
công lượng tử hiệu quả đã biết ở các tham số này.
, and after \appendices add
% % Full proofs for sections/security.tex and the specification/code
% cross-check. Input after \appendices in main.tex. Uses the macros that
% sections/security.tex provides.

\section{Các chứng minh an toàn}
\label{app:security-proofs}

\subsection{Quy ước}

Các oracle ngẫu nhiên được lấy mẫu lười (lazily). Một điểm của $\Hfs$ được
\emph{xác định} một khi nó đã được truy vấn hoặc lập trình; một truy vấn tại
một điểm chưa xác định là \emph{mới} (fresh). $Q_H$ chặn mọi truy vấn oracle
ngẫu nhiên của đối thủ, bao gồm các truy vấn của bên kiểm chứng trong
$\mathcal O_{\mathsf{Ver}}$ và một lần kiểm chứng cuối, còn $Q_P$ chặn số chứng
minh trung thực. Mọi phép rút gọn dưới đây tự chạy toàn bộ thí nghiệm, điều mà
nó làm được vì việc kiểm chứng là công khai và mọi cam kết trung thực đều do
bên thách thức sinh ra. Đối với oracle mở rộng khóa có một điểm duy nhất (thẻ
khóa theo sau bởi $\pid$); một phép rút gọn cần cấy một ma trận thách thức sẽ
lập trình điểm này trước khi đối thủ bắt đầu, điều này khiến $\mathsf{ck}$ được
phân phối đúng như một khóa đều.

\subsection{Chứng minh Bổ đề~\ref{lem:hvzk}}

Cố định $r\in S_1^k$ và $c\in\Cset$, và đặt $v=cr\in\mathbb Z^{kN}$ là tích
negacyclic nguyên. Vì $\norm{v}_\infty\le\norm{c}_1\norm{r}_\infty
\le\kappa<q/2$, việc tâm hóa tích theo modulo $q$ khôi phục $v$ một cách chính
xác, đúng như những gì phần hiện thực tính (\texttt{opening\_proof.py:79--80}).
Với một tọa độ và bất kỳ $a\in[-B,B]$, $\Pr[y_i+v_i=a]=\Pr[y_i=a-v_i]=1/2\gamma$,
bởi vì $a-v_i\in[-B-\kappa,B+\kappa]=[-\gamma+1,\gamma-1]$ nằm trong dải mặt nạ.
Các tọa độ độc lập, nên $\Pr[z=z_0]=(2\gamma)^{-kN}$ với mọi $z_0\in Z$. Lấy
tổng trên $Z$ cho xác suất chấp nhận $|Z|(2\gamma)^{-kN}=p_{\mathrm{acc}}$, và
với điều kiện được chấp nhận thì $z$ đều trên $Z$; không đại lượng nào phụ thuộc
vào $r$ hay $c$. \qed

\subsection{Chứng minh Bổ đề~\ref{lem:sim}}

\emph{Bên chứng minh thực} (\texttt{opening\_proof.py:86--109}). Trên $(x,t,r)$:
rút $e\leftarrow\{0,1\}^{256}$; $K=H_{\mathsf{hg}}(e\|x,t,r\text{ có tiền tố độ dài})$;
với $i=1,\dots,\Amax$: $y_i=H_{\mathsf y}(K\|i)\in Y$, $w_i=B_1y_i$,
$h_i=\Hfs(\pid,x,t,w_i)$, $c_i=G(h_i)$, $z_i=y_i+c_ir$; xuất ra $(h_i,z_i)$ tại
$i$ đầu tiên có $z_i\in Z$, hoặc $\bot$ sau $\Amax$ lần loại bỏ.

\emph{Bộ mô phỏng} $\mathcal S(x,t)$. Xuất ra $\bot$ với xác suất
$(1-p_{\mathrm{acc}})^{\Amax}$. Ngược lại rút $h\leftarrow\{0,1\}^{256}$ và
$z\leftarrow Z$, đặt $w=B_1z-G(h)t_1$; nếu $\Hfs$ đã xác định tại
$(\pid,x,t,w)$ thì loại bỏ; ngược lại lập trình $\Hfs(\pid,x,t,w):=h$ và xuất ra
$(h,z)$. Bộ mô phỏng không bao giờ dùng $r$.

Chúng tôi chặn độ thay đổi trên mỗi chứng minh; một cận hợp (union bound) trên
$Q_P$ chứng minh cho ta bổ đề. Viết $\eta_y=nN\log_2(2\gamma)$. Bởi vì
$B_1=[I_n|B_1']$, $n$ hàng đầu của $w=B_1y$ là $y_{[1..n]}+B_1'y_{[n+1..k]}$,
một ảnh song ánh của $y_{[1..n]}$ với $y_{[n+1..k]}$ cố định; do đó
$H_\infty(w)\ge\eta_y$ cho một mặt nạ đều,
$H_\infty(w)\ge\eta_y-\log_2\frac1{1-p_{\mathrm{acc}}}\ge\eta$ cho một mặt nạ
với điều kiện bị loại bỏ, và $H_\infty(w)\ge\eta$ cho một $w=B_1z-ct_1$ được
mô phỏng với $z$ đều trên $Z$.

\begin{itemize}
\item[$G_0$:] bên chứng minh thực.
\item[$G_1$:] $K$ được thay bằng một giá trị đều. Hai trò chơi chỉ khác nhau
  nếu một truy vấn nào đó tới $H_{\mathsf{hg}}$ khác với truy vấn của chính bên
  chứng minh bắt đầu bằng $e$: tối đa $(Q_H+Q_P)2^{-256}$.
\item[$G_2$:] mỗi $y_i$ được thay bằng một phần tử đều mới của $Y$. Hai trò
  chơi chỉ khác nhau nếu $H_{\mathsf y}$ bị truy vấn ở nơi khác tại $K\|i$: tối
  đa $(Q_H+Q_P)2^{-256}$.
\item[$G_3$:] trong một lần thử bị loại, bên chứng minh rút $h_i$ một cách đều
  và đặt $c_i=G(h_i)$ mà không truy vấn hay xác định $\Hfs$. Hai trò chơi chỉ
  khác nhau nếu điểm $(\pid,x,t,w_i)$, tại bất kỳ thời điểm nào, bị $\mathcal A$
  truy vấn hoặc được một chứng minh trung thực khác xác định: tối đa
  $\Amax(Q_H+\Amax Q_P)2^{-\eta}$. Đây là bước liên quan tới các truy vấn
  oracle ngẫu nhiên ẩn của các lần thử bị loại~\cite{ref:barbosa,ref:devevey}.
\item[$G_4$:] điểm của lần thử được chấp nhận được xác định bởi truy vấn của
  bên chứng minh; nó là mới ngoại trừ với xác suất $(Q_H+(\Amax+1)Q_P)2^{-\eta}$.
  Với điều kiện là mới, $h_i$ đều trong mọi lần thử và, theo
  Bổ đề~\ref{lem:hvzk}, lần thử $i$ được chấp nhận với xác suất đúng bằng
  $p_{\mathrm{acc}}$ bất kể $h_i$, $c_i$ và $r$. Vậy chỉ số của lần chấp nhận
  đầu tiên là hình học bị cắt cụt, và với điều kiện được chấp nhận, $h$ đều và
  $z$ đều trên $Z$, độc lập với nhau và với $r$. Các lần thử bị loại không còn
  liên kết với oracle nào sau $G_3$ và có thể được bỏ đi. Khi đó đầu ra và định
  nghĩa oracle $\Hfs(\pid,x,t,B_1z-G(h)t_1)=h$ được phân phối đúng như trong
  $\mathcal S$, mà sự kiện loại bỏ của nó chính là sự kiện không-mới vừa được
  chặn.
\end{itemize}
Tổng trên mỗi chứng minh là $2(Q_H+Q_P)2^{-256}+(\Amax+1)(Q_H+(\Amax+1)Q_P)2^{-\eta}$.

Đối với $\Pi_{\mathsf{eq}}$ (\texttt{equality\_proof.py:58--104}), bên chứng
minh rút hai mặt nạ mỗi lần thử, thông báo $(w_a,w_b,u)=(B_1y_a,B_1y_b,\langle b_2,y_a-y_b\rangle)$
và chấp nhận khi và chỉ khi cả hai đáp ứng đều nằm trong $Z$. Theo
Bổ đề~\ref{lem:hvzk} áp dụng cho mỗi đáp ứng, sự chấp nhận có xác suất đúng bằng
$p_{\mathrm{acc}}^2$ và $(z_a,z_b)$ được chấp nhận đều trên $Z^2$. Bộ mô phỏng
lấy mẫu $(h,z_a,z_b)$, đặt $w_a=B_1z_a-ct_{1,a}$, $w_b=B_1z_b-ct_{1,b}$,
$u=\langle b_2,z_a-z_b\rangle-c(t_{2,a}-t_{2,b})$ và lập trình $H_{\mathsf{eq}}$;
thông báo có min-entropy ít nhất $\eta$ thông qua $w_a$. Cùng các bước lai cho
cùng cận. \qed

\subsection{Một bổ đề trích xuất bằng rẽ nhánh}

\begin{lemma}[Trích xuất]
\label{lem:fork}
Gọi $\mathcal B$ là một thuật toán với các đồng xu $\rho$ thực hiện nhiều nhất
$Q$ truy vấn mới tới $\Hfs$, được trả lời bởi $h_1,\dots,h_Q\leftarrow\{0,1\}^{256}$
(oracle $G$ và mọi điểm được $\mathcal B$ lập trình đều do $\rho$ xác định).
$\mathcal B$ xuất ra $\bot$ hoặc $(J,x^*,t^*,\pi^*)$ sao cho
$\pi^*=(h_J,z^*)$ được chấp nhận cho $(x^*,t^*)$ và truy vấn mới thứ $J$ là
$(\pid,x^*,t^*,w^*)$ với $w^*=B_1z^*-G(h_J)t_1^*$. Đặt
$\mathrm{acc}=\Pr[\mathcal B\neq\bot]$. Tồn tại $\mathcal E$, chạy
$\mathcal B$ hai lần, xuất ra $x^*$ của lần chạy thứ nhất và một phần mở nới
lỏng của $t^*$ với xác suất $\varepsilon_{\mathsf{ext}}\ge\mathrm{acc}(\mathrm{acc}/Q-2^{-256})
-2^{-257}-Q2^{-256}$; tương đương
\[
\mathrm{acc}\le Q2^{-256}+\sqrt{Q\bigl(\varepsilon_{\mathsf{ext}}+2^{-257}+Q2^{-256}\bigr)}.
\]
\end{lemma}
\begin{proof}
$\mathcal E$ là thuật toán rẽ nhánh của Bellare và Neven~\cite{ref:bn06}
(Bổ đề~1, bổ đề rẽ nhánh tổng quát; xem thêm~\cite{ref:ps00}): chạy
$\mathcal B(\rho;h_1..h_Q)$; nếu nó trả về $J$, chạy lại nó trên
$(\rho;h_1..h_{J-1},h'_J..h'_Q)$ với $h'_J,\dots,h'_Q$ mới; thành công nếu lần
chạy thứ hai trả về cùng $J$ và $h'_J\neq h_J$. Bổ đề đó cho xác suất thành công
ít nhất $\mathrm{acc}(\mathrm{acc}/Q-2^{-256})$. Hai lần chạy giống hệt nhau cho
tới truy vấn mới thứ $J$, nên chúng chia sẻ $(x^*,t^*,w^*)$, và trả về $(h_J,z)$,
$(h'_J,z')$ với $B_1z-ct_1^*=w^*=B_1z'-c't_1^*$ cho $c=G(h_J)$, $c'=G(h'_J)$.
Giá trị $h'_J$ đều và độc lập với lần chạy thứ nhất; trừ khi nó là một trong
nhiều nhất $Q$ đầu vào mà tại đó $G$ đã được tính trong lần chạy thứ nhất
(xác suất $\le Q2^{-256}$), $G(h'_J)$ là một phần tử đều mới của $\Cset$ và
bằng $c$ với xác suất $1/|\Cset|\le2^{-257}$. Ngược lại
$f=c-c'\in\Cdiff$ và $\bar r=z-z'$ thỏa $\norm{\bar r}_\infty\le2B$ và
$B_1\bar r=ft_1^*$. Dạng thứ hai suy ra bằng cách giải bất đẳng thức bậc hai:
$\mathrm{acc}^2-Q2^{-256}\mathrm{acc}-Q\delta\le0$ với
$\delta=\varepsilon_{\mathsf{ext}}+2^{-257}+Q2^{-256}$ kéo theo
$\mathrm{acc}\le Q2^{-256}+\sqrt{Q\delta}$.
\end{proof}

\subsection{Chứng minh Định lý~\ref{thm:se}}

$G_0$ là $\mathsf{SE}$. Trong $G_1$, mọi chứng minh trung thực đều do bộ mô
phỏng tạo ra, nên $|\Pr[W_0]-\Pr[W_1]|\le\epszk(Q_P,Q_H)$, trong đó $W_i$ là sự
kiện thắng. Gọi $\mathsf{Prog}$ là sự kiện $\mathcal A$ thắng trong $G_1$ và
điểm $(\pid,x^*,t^*,w^*)$ của bên kiểm chứng đã được bộ mô phỏng lập trình, cho
một bộ ba $(x_i,t_i,\pi_i)\in L$. Khi đó $x^*=x_i$, $t^*=t_i$, $w^*=w_i$, và
sự chấp nhận buộc $h^*=h_i$, nên $c^*=c_i$. Vì $(x^*,t^*,\pi^*)\notin L$, ta
có $z^*\neq z_i$, và $B_1z^*-c_it_1^*=B_1z_i-c_it_1^*$ cho
$B_1(z^*-z_i)=0$ với $0<\norm{z^*-z_i}_\infty\le2B\le\beta$. Một đối thủ M-SIS
cấy ma trận của nó làm $B_1'$, tự rút $b_2'$, chạy $G_1$ và xuất ra $z^*-z_i$
thành công mỗi khi $\mathsf{Prog}$ xảy ra, nên
$\Pr[\mathsf{Prog}]\le\Adv^{\MSIS}$. Ngược lại, điểm của bên kiểm chứng được
xác định bởi một truy vấn mới, và $\mathcal B$ (chạy $G_1$ và xuất ra chỉ số
của truy vấn đó) thỏa giả thiết của Bổ đề~\ref{lem:fork} với
$Q=Q_H$ và $\mathrm{acc}\ge\varepsilon-\epszk-\Adv^{\MSIS}$. Thay vào cho ta
định lý.

Truy vấn bị rẽ nhánh chứa $x^*$, thứ mà bên kiểm chứng dựng từ các mục ghim của
nó và từ $\mathsf{Read}(\vid^*)$ (\texttt{video\_pipeline.py:278--284}); nó
không bao giờ lấy một trường phát biểu nào từ bên chứng minh ngoại trừ $N_s^*$,
thứ mà nó kiểm tra đối chiếu với kho của nó (hoặc với giá trị cố định) và bản
thân là một phần của $x^*$. Do đó phần mở trích xuất được được ràng buộc vào các
giá trị của bên kiểm chứng cho $\mathit{rid}$, $\pid$, gốc sổ đăng ký và epoch,
$N_s$, chính sách codec, hồ sơ vật mang, $|P|$, $\mathrm{SHA256}(P)$ và $V$. Hai
phát biểu khác nhau ở bất kỳ trường nào đều khác nhau khi là chuỗi, vì mọi
trường đều có độ rộng cố định (\texttt{video\_proof.py:68--77}). \qed

\subsection{Chứng minh Định lý~\ref{thm:binding}}

Cho $(\bar r,f,m)$ và $(\bar r',f',m')$ mở $t$ với $m\neq m'$, trong đó một
phần mở trung thực có $f=1$ và $\norm{\bar r}_\infty\le1$. Từ $ft_1=B_1\bar r$
và $f't_1=B_1\bar r'$ ta được $B_1x=0$ cho $x=f'\bar r-f\bar r'$. Từ
$ft_2=\langle b_2,\bar r\rangle+fm$ và $f't_2=\langle b_2,\bar r'\rangle+f'm'$,
\[
\langle b_2,x\rangle=f'(ft_2-fm)-f(f't_2-f'm')=ff'(m'-m).
\]
Nếu $x\equiv0\pmod q$ thì $ff'(m'-m)=0$ và, vì $f$ và $f'$ khả nghịch
(Bổ đề~\ref{lem:invertible}; $1$ hiển nhiên), $m=m'$, mâu thuẫn. Chuẩn
của $x$ tối đa là $\norm{f'}_1\norm{\bar r}_\infty+\norm{f}_1\norm{\bar r'}_\infty$,
tức $2\kappa\cdot2B+2\kappa\cdot2B=8\kappa B$ cho hai phần mở nới lỏng,
$2\kappa\cdot1+1\cdot2B$ cho một nới lỏng và một trung thực, và $2$ cho hai phần
mở trung thực. Tất cả đều tối đa $\beta<q/2$ (hàm khởi tạo tham số bắt buộc
$2\beta<q$, \texttt{params.py:53--54}), nên $x$ là một vector ngắn khác không
cả theo modulo $q$ lẫn trên $\mathbb Z$. Phép rút gọn cấy ma trận của nó làm
$B_1'$, lấy mẫu $b_2'$ và xuất ra $x$. \qed

\subsection{Proof of Lemma~\ref{lem:row}}

\begin{lemma}[Hiding instance]
\label{lem:mlwe-instance}
There is an efficient map that, for any $m$ chosen after seeing
$\mathsf{ck}$, sends an M-LWE instance $(A,u)$, $A\in\Rq^{(n+1)\times\lambda}$,
to $(\mathsf{ck},t)$ with $\mathsf{ck}$ uniform, such that $t=\Com(m;r)$ with
$r\leftarrow S_1^k$ if $u=As+e$, and $t$ is uniform if $u$ is uniform.
\end{lemma}
\begin{proof}
Write $A=\bigl[\begin{smallmatrix}A_1\\a_2\end{smallmatrix}\bigr]$ and
$u=(u_1,u_2)$. Draw $B_{1,2}'\leftarrow\Rq^{n\times1}$ and set $b_2'=a_2$,
$B_{1,3}'=A_1+B_{1,2}'a_2$, which makes $\mathsf{ck}$ uniform. Given $m$, set
$t_2=u_2+m$ and $t_1=u_1+B_{1,2}'u_2$. If $u=As+e$ with $e=(e_1,e_2)$ and
$r=(e_1,e_2,s)$, then $\langle b_2,r\rangle=e_2+a_2s=u_2$ and
$B_1r=e_1+B_{1,2}'e_2+(A_1+B_{1,2}'a_2)s=u_1+B_{1,2}'u_2$, so $t=\Com(m;r)$
with $r$ uniform ternary. If $u$ is uniform, so is $t$, because
$(u_1,u_2)\mapsto(u_1+B_{1,2}'u_2,u_2+m)$ is a bijection.
\end{proof}

Now let $\mathcal A$ output $(j,\bar r,f)$ with $f\in\Cdiff$,
$\norm{\bar r}_\infty\le2B$ and $ft^{(j)}_1=B_1\bar r$, and put $x=fr_j-\bar r$.
Then $B_1x=ft_1^{(j)}-ft_1^{(j)}=0$ and $\norm{x}_\infty\le2\kappa+2B<q/2$.

\emph{Case $x\neq0$.} An M-SIS adversary plants its matrix as $B_1'$, samples
$b_2'$, generates every $t^{(j)}$ itself (so it knows every $r_j$), and outputs
$x$. Hence $\Pr[x\neq0]\le\Adv^{\MSIS}$.

\emph{Case $x=0$.} Then $\bar r\equiv fr_j$, so $r_j$ is the centered
representative of $f^{-1}\bar r\bmod q$. An M-LWE distinguisher guesses
$j_0\leftarrow[Q]$, builds $\mathsf{ck}$ and $t^{(j_0)}$ from its instance by
Lemma~\ref{lem:mlwe-instance} (with the message $m_{j_0}$ that $\mathcal A$
chooses), generates the other commitments honestly, and outputs $1$ iff
$\mathcal A$ returns $j=j_0$ and $r'=\mathrm{center}(f^{-1}\bar r)$ is ternary
with $\Com(m_{j_0};r')=t^{(j_0)}$. On a real instance it outputs $1$ with
probability at least $\Pr[x=0]/Q$. On a uniform instance $t^{(j_0)}$ is uniform
and independent of $m_{j_0}$, so a ternary opening exists with probability at
most $|S_1^k|/q^{(n+1)N}=3^{kN}q^{-(n+1)N}$. Hence
$\Pr[x=0]\le Q(\Adv^{\MLWE}+3^{kN}q^{-(n+1)N})$. \qed

\subsection{Proof of Theorem~\ref{thm:hiding}}

$H_0$ is $\mathsf{IND}^0$. In $H_1$ all honest proofs, including the challenge
proof, are simulated: $|H_0-H_1|\le\epszk(Q_P+1,Q_H)$. In $H_2$ the challenge
commitment is replaced by a uniform pair; a distinguisher between $H_1$ and
$H_2$ is an M-LWE distinguisher by Lemma~\ref{lem:mlwe-instance} with
$m=\mathsf{enc}(p_0)$: it builds $\mathsf{ck}$ from its instance, generates the
other commitments honestly, answers $\mathcal O_{\mathsf{Disc}}$ and
$\mathcal O_{\mathsf{Link}}$ for them, and the challenge randomness is never
needed because the challenge proof is simulated and disclosure and linkage of
the challenge are excluded. So $|H_1-H_2|\le\Adv^{\MLWE}$. The stego video in
$H_2$ is $\mathsf{Write}(\vid,P,E)$ with $P$, $D$, $x$ independent of the
payload (fixed from a placeholder of the final length,
\texttt{video\_pipeline.py:216--227}), and $E$ contains only $\pid$, $N_s$, a
uniform $t$ and a simulated proof; $H_2$ is therefore independent of $b$. The
same three hybrids from $\mathsf{IND}^1$ end in the same $H_2$, which gives
the bound. The payload length is hidden because $\mathsf{enc}$ maps every
payload of at most 766 bytes to one polynomial and the envelope has a fixed
size. \qed

\subsection{Proof of Theorem~\ref{thm:replay}}

We model the store as a linearizable table keyed by
$d(N)=\mathrm{SHA256}(\text{tag}\|N)$ (\texttt{zkp\_sessions.py:113--117}).
Rows are inserted only by $\mathsf{Issue}$ with a fresh uniform $N$
(\texttt{:132--151}) and deleted only once expired (\texttt{:140--143}).
Acceptance requires the preflight \texttt{is\_active} and then
\texttt{consume}, an \texttt{UPDATE} with the conditions
\texttt{consumed\_at IS NULL}, matching binding and \texttt{expires\_at>now},
committed inside \texttt{BEGIN IMMEDIATE}, returning success iff exactly one
row changed (\texttt{:181--205}).

(i) After the first successful \texttt{consume} of a row, the condition
\texttt{consumed\_at IS NULL} fails for that row forever; a second success for
the same $N_s$ needs another row with key $d(N_s)$, i.e.\ a later issuance of
some $N'$ with $d(N')=d(N_s)$: either $N'=N_s$ (probability at most
$Q_I^22^{-257}$ over all pairs of issuances) or a SHA-256 collision.
(ii) Any row with key $d(N_s^*)$ comes from an issuance of some $N$ with
$d(N)=d(N_s^*)$, so $N_s^*$ was issued unless SHA-256 collides; the binding
and expiry are part of the \texttt{WHERE} clause, evaluated at consumption
time. (iii) The session field occupies bytes $106$--$137$ (0-indexed) of every statement,
and $x^*$ is a substring of the point $(\pid,x^*,t^*,w^*)$. A query made
before the issuance of $N_s^*$ fixes this field before $N_s^*$ is drawn, so it
matches with probability $2^{-256}$; a union bound over $Q_H$ queries and $Q_I$
sessions gives $Q_HQ_I2^{-256}$. \qed

\subsection{Proof of Theorem~\ref{thm:auth}}

Let $W_a$, $W_b$ be the two winning cases.

$W_a$: $\mathsf{Pins}$ verifies $\mathsf{Reg}^*$ under $\mathit{pk}_I$
(\texttt{video\_pipeline.py:146--147}). If its signing bytes were never signed
through $\mathcal O_{\mathsf{Reg}}$, it is an EUF-CMA forgery; the reduction
runs the whole experiment, using its signing oracle for
$\mathcal O_{\mathsf{Reg}}$. Otherwise $\mathsf{Reg}^*$ is a signed registry
and contains a descriptor equal to the registered descriptor of $t^*$
(\texttt{:150--153}). That descriptor has the registered constraint-module
hash, so it is not in any $\Delta$ (in-band descriptors have a different
module hash, unless SHA-256 collides on the two fixed module strings), and it
was computed from some $t'\in T_{\mathsf{reg}}$. Since $t^*\notin T_{\mathsf{reg}}$,
$t'\neq t^*$ and the canonical encodings differ, so
$\mathrm{SHA256}(B_1'\|b_2'\|t^*)=\mathrm{SHA256}(B_1'\|b_2'\|t')$
(\texttt{:119--128}) is a collision. Hence
$\Pr[W_a]\le\Adv^{\mathsf{euf\text{-}cma}}_{\mathsf{ML\text{-}DSA}}+\Adv^{\mathsf{cr}}_{\mathsf{SHA2}}$.

$W_b$: simulate all honest proofs ($\epszk$) and handle the programmed-point
case as in Theorem~\ref{thm:se} ($\Adv^{\MSIS}$). The remaining runner
$\mathcal B$ receives the honest credentials as a $Q_C$-instance challenge of
Lemma~\ref{lem:row}, never uses their randomness (all proofs are simulated,
and a disclosure query for an instance aborts the run, which does not affect
runs in which $\mathcal A$ wins with an undisclosed $j$), and outputs the
index of the verifier's fresh query. Lemma~\ref{lem:fork} extracts a relaxed
opening of $t^*=t_c^{(j)}$, which is a Lemma~\ref{lem:row} adversary, so
$\varepsilon_{\mathsf{ext}}\le\Adv^{\mathsf{row}}_{Q_C}$. Adding the cases
gives the theorem.

A tight alternative for $W_b$ follows the ML-DSA analysis: after replacing the
credential by a uniform pair (M-LWE), an accepted proof is a solution of the
SelfTargetMSIS problem for $[B_1\,|\,t_1]$ and $\Hfs$~\cite{ref:kls18}, without
rewinding. We state the rewinding bound because it uses only the estimated
M-SIS and M-LWE instances. \qed

\subsection{Proofs for linkage (Theorem~\ref{thm:link})}

\begin{lemma}[Special soundness of the equality proof]
\label{lem:eq-ss}
Two accepting transcripts $((w_a,w_b,u),c,z_a,z_b)$ and
$((w_a,w_b,u),c',z_a',z_b')$ with $c\neq c'$ yield relaxed openings
$(\bar r_a,f)$ of $t_a$ and $(\bar r_b,f)$ of $t_b$, $f=c-c'$, with equal
messages.
\end{lemma}
\begin{proof}
The verifier checks $B_1z_a=w_a+ct_{1,a}$, $B_1z_b=w_b+ct_{1,b}$ and
$\langle b_2,z_a-z_b\rangle=u+c(t_{2,a}-t_{2,b})$
(\texttt{equality\_proof.py:135--145}). Subtracting the two transcripts gives
$B_1\bar r_a=ft_{1,a}$, $B_1\bar r_b=ft_{1,b}$ and
$\langle b_2,\bar r_a-\bar r_b\rangle=f(t_{2,a}-t_{2,b})$ with
$\bar r_\ast=z_\ast-z_\ast'$, $\norm{\bar r_\ast}_\infty\le2B$. Then
$\bar m_a-\bar m_b=(t_{2,a}-t_{2,b})-f^{-1}\langle b_2,\bar r_a-\bar r_b\rangle=0$.
\end{proof}

\emph{Soundness.} Simulate $\mathcal O_{\mathsf{Link}}$ ($\epszk$). If the
verifier's $H_{\mathsf{eq}}$ point was programmed for the same context and
commitments, then $c$ is the programmed one and $(z_a,z_b)\neq(z_a^{(i)},z_b^{(i)})$
with $B_1(z_a-z_a^{(i)})=B_1(z_b-z_b^{(i)})=0$, an M-SIS solution of norm $2B$.
Otherwise Lemma~\ref{lem:fork}, with the forked oracle $H_{\mathsf{eq}}$ and
Lemma~\ref{lem:eq-ss} in place of the opening relation, extracts openings with
a common message $\bar m$. The adversary's first-run openings have messages
$m_a\neq m_b$, so $m_a\neq\bar m$ or $m_b\neq\bar m$, and
Theorem~\ref{thm:binding} applied to $t_a$ or $t_b$ gives an M-SIS solution.
Hence $\varepsilon_{\mathsf{ext}}\le\Adv^{\MSIS}$ in the bound of
Theorem~\ref{thm:se}.

\emph{Hiding.} Simulate the proof ($\epszk(Q_P+1,Q_H)$), then replace $t_a$ and
$t_b$ by uniform pairs, one M-LWE hybrid each (Lemma~\ref{lem:mlwe-instance};
the randomness is independent). The final game is independent of $b$; doing
the same from $b=1$ gives $2\epszk+4\Adv^{\MLWE}$.

\emph{Disclosure by proof.} The equality proof with $t_b=(0,\mathsf{enc}(p))$
and $r_b=0$ (a ternary opening of $t_b$) proves that $t_a$ hides $p$ without
revealing $r_a$: extraction gives a relaxed opening $(\bar r_b,f)$ of $t_b$
with $B_1\bar r_b=0$, so $\bar r_b=0$ unless M-SIS is broken, and then
$\bar m_b=\mathsf{enc}(p)=\bar m_a$. This is the fix proposed for item~D3.
\qed

\subsection{Proof of Lemma~\ref{lem:canon} and Theorem~\ref{thm:integrity}}

\emph{Unique decodability.} For a carrier segment the hashed string is the
domain tag followed, for each NAL unit, by \texttt{"NAL\textbackslash0"} and a
fixed 12-byte header, then either \texttt{"RAW\textbackslash0"}, a 64-bit
length and the RBSP, or a sequence of \texttt{"BITS\textbackslash0"}, a 64-bit
bit count and the packed bits, alternating with \texttt{"BLOCK\textbackslash0"},
the 16-byte key and length, and the signed 64-bit levels
(\texttt{video\_canonicalization.py:234--307}). The four tags differ in their
first two bytes, every field is fixed-width or length-prefixed, and the
padding of packed bits is determined by the bit count, so a left-to-right
parser recovers the whole sequence. The segment combiner hashes a tag, the
32-bit frames-per-segment and segment count, and fixed-width
$(i,f_i,d_i)$ triples (\texttt{blind\_sync\_streaming.py:225--233}); $V$ hashes a
tag and $D$. A difference in $\Canon$ is therefore a difference in the input
of one of the three SHA-256 layers whose outputs are equal, which is a
collision. Carrier segments are exactly those containing positions of $P$, so
for equal $P$ both sides hash each segment in the same mode.

\emph{What is determined.} Fix $\Canon$ and $P$. Segments without carriers
are fixed byte for byte. In a carrier segment, every NAL header field in the
serialization, every RBSP of a NAL unit without carriers, and every RBSP bit
outside carrier blocks up to the stop bit are fixed. Each carrier block is
fixed as a level vector except at carriers, where only the normalized value
$\nu=\mathrm{sgn}(\ell)2\lfloor|\ell|/2\rfloor$ is fixed. The pipeline enforces
the stable rule (\texttt{blind\_sync\_streaming.py:268--271}; carriers have
$|\ell|\ge4$ and the embedder flips the magnitude parity), so
$\ell=\nu+\mathrm{sgn}(\nu)b$ with $b$ the carrier bit; for the optional
trailing-one sign carrier, $\ell=\pm1$ with the sign as the bit. The carrier
bits are written in the order of $P$ as the chunk-framed encoding of $E$.
If that encoding is unique, $E$ fixes the bits and hence every carrier level.

\emph{Theorem~\ref{thm:integrity}.} Let $(x_h,t,\pi_h)\in L$ be the honest
triple. If $(x^*,t,\pi^*)\notin L$, the event is bounded as in
Theorem~\ref{thm:auth}, case $W_b$, with the $Q_P$ honest commitments as the
Lemma~\ref{lem:row} instances. Otherwise $x^*=x_h$ and $\pi^*=\pi_h$ (each
honest in-band commitment occurs in one triple except with negligible
probability; a credential occurs in several, all honest). Equal statements
give equal $|P|$ and $\mathrm{SHA256}(P)$, hence $P^*=P_h$ unless SHA-256
collides (the positions are hashed as canonical JSON of integer triples), and
equal $V$, hence equal $\Canon$ by the first part. The envelope is
$E^*=\texttt{magic}\|1\|\pid\|N_s\|t\|\pi^*$, and its parser accepts only this
exact layout (\texttt{video\_proof.py:142--153, 218--228}), so $E^*=E_h$. \qed

\subsection{Completeness}

An accepted attempt satisfies $B_1z-ct_1=B_1y+cB_1r-cB_1r=w$, and for the
equality proof $\langle b_2,z_a-z_b\rangle-c(t_{2,a}-t_{2,b})=\langle b_2,y_a-y_b\rangle$
because $t_{2,a}-t_{2,b}=\langle b_2,r_a-r_b\rangle$ for equal messages. By
Lemma~\ref{lem:hvzk} the number of attempts is geometric with parameter
$p_{\mathrm{acc}}$ (mean $5.23$) for the opening proof and $p_{\mathrm{acc}}^2$
(mean $27.4$) for the equality proof. With the caps $\Amax=512$
(\texttt{opening\_proof.py:33}) and $\Amax=4\,096$
(\texttt{equality\_proof.py:50}) the prover fails with probability
$2^{-156.7}$ and $2^{-220.0}$ respectively. Under the shared cap of 512 that
the reviewed code used, the equality proof failed with probability
$2^{-35.0}$ for \texttt{bdlop-open-128-v1} (means $4.65$ and $21.6$); this is
item~D4.

\section{Specification and implementation cross-check}
\label{app:code-check}

Each item lists the property relied on above and where the implementation
realizes it, at commit \texttt{ea96d60}. Paths are relative to
\texttt{VideoLevel/src}; \texttt{opening\_proof.py}, \texttt{equality\_proof.py},
\texttt{bdlop.py}, \texttt{params.py}, \texttt{ring.py}, \texttt{codec.py},
\texttt{sampling.py}, \texttt{video\_proof.py}, \texttt{video\_pipeline.py}
and \texttt{cli.py} are in \texttt{lattice\_zk/}.

\begin{enumerate}
\item \emph{Statement.} Field order and widths: \texttt{u16} version,
  $\mathit{rid}$, $\pid$, root, \texttt{u64} epoch, $N_s$, codec policy,
  carrier profile, \texttt{u64} carrier count, positions hash, $V$; every
  digest checked to be 32 bytes (\texttt{video\_proof.py:58--77}).
  Matches the specification, \S3.6. \textbf{Consistent.}
\item \emph{Fiat--Shamir input.} SHAKE-256 over the tag, $\pid$,
  \texttt{u64}$(|x|)\|x$, the full commitment (6\,720 bytes) and $w$ as
  big-endian 64-bit words; 32-byte output
  (\texttt{opening\_proof.py:42--50}). Only $x$ is length-prefixed; the other
  fields have lengths fixed by $\pid$. Equality proof: same with both
  commitments and $(w_a,w_b,u)$ (\texttt{equality\_proof.py:40--51}).
  \textbf{Consistent}; the specification's ``with length prefixes'' (\S4.7)
  was imprecise and has been corrected.
\item \emph{Hedging.} $K=\mathrm{SHA3}(e\|\text{length-prefixed }x,t,r)$
  (\texttt{opening\_proof.py:53--63,99--100}). \textbf{Consistent} (see D7).
\item \emph{Samplers.} Masks: exact uniform on $[-\gamma+1,\gamma]$ (18-bit
  draws below $2\gamma=2^{18}$, \texttt{sampling.py:60--66}); challenges:
  sign bits then Fisher--Yates with exact rejection, as in ML-DSA
  \textsf{SampleInBall} (\texttt{:69--81}); ternary randomness by rejection on
  2-bit draws (\texttt{:54--57}). \textbf{Consistent} with the uniformity used in
  Lemmas~\ref{lem:hvzk} and~\ref{lem:fork}.
\item \emph{Verification.} Exact proof length, packed responses below
  $2B+1$, zero padding, recomputed $w'$, constant-time seed comparison
  (\texttt{opening\_proof.py:118--138}). The packing has no padding bits at
  these parameters, so the proof encoding is canonical. \textbf{Consistent.}
\item \emph{Parameter checks.} Prime $q\equiv5\pmod8$, $2\sqrt2<\sqrt q$,
  $2\cdot8\kappa B<q$ (\texttt{params.py:37--54}, \texttt{ring.py:52--58});
  production pins require $|\Cset|\ge2^{256}$
  (\texttt{params.py:127--130}, \texttt{video\_pipeline.py:144--145}).
  \textbf{Consistent.}
\item \emph{Envelope.} In-band: exact length, magic \texttt{ZKLB}, version,
  pinned $\pid$, canonical commitment ($<q$ per coefficient), proof
  (\texttt{video\_proof.py:142--153}); the header session must equal the
  statement's (\texttt{:183}). Registered: \texttt{ZKLR}, header and proof only,
  pinned commitment (\texttt{:218--244}). \textbf{Consistent} with \S3.7.
\item \emph{Sessions.} One-use semantics as in the proof of
  Theorem~\ref{thm:replay}; binding $\mathsf{bind}(\mathsf{pins})$ hashes
  $\pid$, $\mathit{rid}$, root, epoch, codec policy and profile
  (\texttt{video\_pipeline.py:166--172}), so in registered mode it also
  depends on the credential through $\mathit{rid}$. Stateless mode requires
  an explicit \texttt{expected\_session} (\texttt{:262--263}).
  \textbf{Consistent.}
\item \emph{Registry.} ML-DSA-65 over a tagged canonical JSON encoding; the
  registered descriptor hashes the key and the credential
  (\texttt{video\_pipeline.py:109--131}, \texttt{zkp\_registry.py:128--147}).
  \textbf{Consistent}, but see D2.
\item \emph{Canonical digest.} As described in
  Section~\ref{sec:sec-integrity}; the specification (\S3.6) described $D$ as
  a single-stream digest and has been corrected to the segmented form.
\item \emph{Domain tags.} SHAKE-256 tags \texttt{.../bdlop-open/v1/\{key,
  randomness, fiat-shamir, mask\}\textbackslash0},
  \texttt{.../bdlop-equal/v1/\{fiat-shamir, mask\}\textbackslash0} and
  \texttt{.../challenge-poly/v1\textbackslash0} each end in the only NUL byte
  and are pairwise distinct, hence prefix-free. \textbf{Consistent.}
\end{enumerate}

\subsection{Discrepancies and gaps}

\begin{description}
\item[D1 (medium): carrier bits are malleable with a single chunk.]
  \texttt{blind\_payload\_chunks.py:12,84--137} accept a one-chunk set for any
  \texttt{chunk\_size} in $[|E|,65\,535]$, and the verifier never compares the
  extracted chunking with the pinned \texttt{profile.chunk\_bytes}
  (\texttt{video\_pipeline.py:} \texttt{267--303}). When the profile's chunk
  size is at least the envelope length, an adversary can rewrite the 16-bit
  field and the outer CRC-32: up to 48 carrier coefficients change by $\pm1$
  while $E$, $P$ and $D$ are unchanged, and the video is accepted. Rewriting
  8\,000 to 20\,000 changed 21 carrier bits in a direct test of the framing
  functions. This contradicts goal (iii) of the threat model (``any change to
  the carriers is rejected''). Fix: the verifier recomputes
  $\mathsf{Enc}_{\mathsf{prof}}(E)$, the chunks of $E$ under the pinned chunk
  size, and requires it to equal the extracted pieces (or requires the
  header field to equal the pinned size); add a negative test.
\item[D2 (medium): no holder identity in the signed registry.] The registered
  descriptor (\texttt{video\_pipeline.py:109--131}; field set
  \texttt{zkp\_registry.py:29--35}) binds a credential commitment, not a
  holder. ``Only the credential holder can produce the proof'' holds per
  commitment; which person that is is not covered by the ML-DSA signature.
  Fix: add a holder identifier (or its digest) to the registered descriptor.
\item[D3 (medium): disclosure hands over the commitment.]
  \texttt{verify\_disclosure} takes the randomness
  (\texttt{video\_proof.py:84--89,188--197}). The recipient can then produce
  fresh opening proofs for that commitment on new statements, i.e.\ accepted
  in-band videos with the creator's commitment, or impersonate a holder if a
  credential is disclosed. Fix: disclose $p$ with the equality proof against
  $(0,\mathsf{enc}(p))$ and zero randomness, bound to a disclosure context
  (proof above).
\item[D4 (low): equality-proof completeness $2^{-35}$.]
  \texttt{equality\_proof.py:29,83} reuses $\Amax=512$ with acceptance
  $p_{\mathrm{acc}}^2$ (figures for \texttt{bdlop-open-128-v1}, the set
  reviewed; $2^{-27.5}$ for \texttt{bdlop-open-128-v2}). Fix: a separate cap,
  e.g.\ 4\,096 attempts ($\approx2^{-280}$ for v1, $2^{-220}$ for v2); a larger cap does not affect zero knowledge because
  acceptance is witness-independent.
\item[D5 (low): scope of the canonical digest.] The digest is over
  ffmpeg-segmented files (\texttt{blind\_sync\_streaming.py:117--215}); segments
  without carriers are hashed as raw files without a domain tag
  (\texttt{:236--241,358,549}); RBSP bits after the stop bit of a carrier-bearing
  IDR NAL unit are excluded (\texttt{video\_canonicalization.py:}
  \texttt{279, 303, 310--317});
  carrier blocks are hashed as level vectors; the codec policy string
  (\texttt{video\_pipeline.py:25,98--99}) does not pin the segmenter version.
  None breaks Theorem~\ref{thm:integrity}, but each widens the set of
  unbound changes listed in Section~\ref{sec:sec-integrity}. Fix: tag the
  raw-segment digests, hash the full RBSP of carrier NAL units, and pin or
  replace the segmenter.
\item[D6 (low): API does not tie pins to the credential.]
  \texttt{verify\_video\_proof} accepts \texttt{pins} and \texttt{credential}
  separately (\texttt{video\_pipeline.py:250--265}) and does not check that
  \texttt{pins.relation\_id} is the relation of that credential. The CLI derives
  both together (\texttt{cli.py:98--107}), but a caller that pins one
  credential and passes another verifies an unregistered credential. Fix:
  recompute the descriptor of \texttt{credential} and compare relation
  identifiers inside \texttt{verify\_video\_proof}.
\item[D7 (info): hedging hash has no domain tag.]
  \texttt{opening\_proof.py:59} hashes $e\|\dots$ directly. Separation from the
  tagged SHA3-256 uses relies on $e$ being uniform; add a tag.
\item[D8 (info): unreduced commitments.] \texttt{Commitment} objects built
  without \texttt{from\_bytes} may hold coefficients in $[q,2^{30})$, which
  \texttt{pack\_uints} accepts (\texttt{codec.py:10--16}). Envelopes are
  unaffected (their decoder checks $<q$); pinned credentials loaded by other
  code paths are. Fix: range check in \texttt{Commitment}.
\item[D9 (info): one global key.] $\mathsf{ck}$ is expanded from a fixed public
  seed (\texttt{params.py:133--149}); all users share one M-SIS/M-LWE
  instance, so precomputation against it amortizes.
\item[D10 (info):] \texttt{verify\_statement\_binding} hard-codes the
  \texttt{lazer-v1} suite (\texttt{zkp\_registry.py:206--210}) and cannot serve
  \texttt{bdlop-open-v1}; the BDLOP path uses \texttt{pins\_from\_registry}.
\item[D11 (info): specification wording.] \S4.5 wrote the bijection as
  $(t_1,t_2)\mapsto(t_1-B_{1,2}'t_2,t_2)$; the correct map uses $t_2-m$
  (Lemma~\ref{lem:mlwe-instance}). Corrected. \S4.2 cites an ``ML-DSA
  specification, Appendix~B'' for the zero-knowledge argument; this reference
  should be checked.
\end{description}

\paragraph{Status after the review.} D1--D4 and D6 are fixed in the code
evaluated in Section~\ref{sec:evaluation}, each with a negative test: the
verifier compares the extracted chunks with the canonical chunking of $E$
under the pinned chunk size (D1); the registered descriptor can bind a
holder-identifier digest (D2); disclosure is the zero-randomness equality
proof of Section~\ref{sec:sec-soundness} and the randomness-revealing API was
removed (D3); the equality proof has its own 4\,096-attempt cap (D4); and
\texttt{verify\_video\_proof} recomputes the relation identifier of the
supplied parameters, profile, credential and holder identifier and refuses
mismatching pins (D6). D5, D7 and D8 and the reviewer question on
production-strength enforcement are closed in a second round, again each
with negative tests (\texttt{src/runtest/test\_hardening\_d5\_d7\_d8.py}).
D7: the hedged mask key hashes a domain tag, the protocol label
(opening, equality or binary proof) and the parameter identifier before the
entropy. Only \texttt{bdlop-open-128-v1} keeps the untagged derivation, so
its known-answer vectors are unchanged; the \texttt{bdlop-open-128-v2} and
binary-proof vectors were regenerated (commitments unchanged). The verifier
never recomputes a mask, so proofs from either derivation verify, which the
tests check. D8: a \texttt{Commitment} carries its parameter set and its
constructor rejects coefficients outside $[0,q)$ and wrong shapes, storing
read-only copies; verifiers still re-check objects built around the
constructor. Production strength: every application-level verifier
(\texttt{verify\_video\_proof}, the envelope, disclosure and linkage
verifiers, and \texttt{pins\_from\_registry}) refuses a parameter set with a
challenge space below $2^{256}$ unless the caller sets
\texttt{allow\_toy\_parameters}; hand-built pins with the toy set are refused.
D5 is closed by a new, opt-in codec policy (v2) with its own identifier, so
it changes profile, policy, relation and statement identifiers, while codec
policy v1 stays byte-identical for existing videos (recorded identifiers and
digests are tested). Under v2 (a) a segment without carriers is hashed with a
domain tag and its length and every segment digest is typed in the combined
digest; (b) a carrier-bearing IDR NAL unit is hashed over its complete RBSP
with its length, so zero bits after the stop bit are bound (under v1 a zero
byte appended there is accepted; under v2 it is rejected, through the real
H.264 transport); (c) the segmentation is pinned: the segments must be
exactly the access-unit partition of the input, compared byte for byte, so
the digest is a function of the input alone and a segmenter that cuts
elsewhere fails closed. A registered \texttt{bdlop-open-128-v2} run on a
200-frame CIF cut under v2 verified from the video alone, and the replayed
session, a one-bit change and the zero-extended carrier NAL unit were
rejected (\texttt{benchmark/results/q2/hardening\_e2e\_container\_cif\_200f\_v2.json}).
The ffmpeg version is not hashed, because that
would bind the verifier's installation rather than the video. Carrier blocks
are still hashed as level vectors, which is what lets the proof bind the
video without depending on its own bits. D9--D11 remain as described.
. Requires amsmath, amssymb, amsthm and the
% theorem environments definition, lemma, theorem defined in main.tex. Extra
% bibliography entries are in sections/security_refs.bib
% (\bibliography{references,sections/security_refs}).
%
% Every code reference is to commit ea96d60 of feature/lattice-pq-stack.

\providecommand{\Rq}{\mathcal{R}_q}
\providecommand{\norm}[1]{\lVert #1 \rVert}
\providecommand{\Adv}{\mathsf{Adv}}
\providecommand{\pid}{\mathit{pid}}
\providecommand{\Cset}{\mathcal{C}}
\providecommand{\Cdiff}{\bar{\mathcal{C}}}
\providecommand{\Com}{\mathsf{Com}}
\providecommand{\vid}{\mathsf{vid}}
\providecommand{\Hfs}{H_{\mathsf{fs}}}
\providecommand{\Stmt}{\mathsf{Stmt}}
\providecommand{\Canon}{\mathsf{Canon}}
\providecommand{\MSIS}{\textsf{M-SIS}}
\providecommand{\MLWE}{\textsf{M-LWE}}
\providecommand{\epszk}{\varepsilon_{\mathsf{zk}}}
\providecommand{\Amax}{A_{\max}}

\section{Mô hình và phân tích an toàn}
\label{sec:security-model}

Mục này phát biểu hệ thống chứng minh được điều gì và dưới những giả thiết nào.
Mọi kết quả đều nằm trong mô hình oracle ngẫu nhiên cổ điển (ROM). Không có định
lý nào trong mô hình oracle ngẫu nhiên lượng tử (QROM) được khẳng định
(Mục~\ref{sec:sec-qrom}). Các chứng minh đầy đủ nằm trong
Phụ lục~\ref{app:security-proofs}; việc kiểm tra chéo từng định nghĩa đối chiếu
với phần hiện thực nằm trong Phụ lục~\ref{app:code-check}.

\subsection{Ký hiệu, tham số và giả thiết}
\label{sec:sec-notation}

Đặt $\mathcal R=\mathbb Z[X]/(X^N+1)$ và $\Rq=\mathcal R/q\mathcal R$, với
$N=256$ và số nguyên tố $q=2^{30}-35\equiv5\pmod 8$. Khóa cam kết là
$\mathsf{ck}=(B_1,b_2)$ với $B_1=[I_n\,|\,B_1']\in\Rq^{n\times k}$,
$b_2=[0^n\,|\,1\,|\,b_2']\in\Rq^{1\times k}$, $k=n+\lambda+1$, $n=6$, $\lambda=7$ (nên $k=14$).
$S_1$ là tập các đa thức có hệ số trong $\{-1,0,1\}$. Với một thông điệp
$m\in\Rq$ và $r\in S_1^k$,
$\Com(m;r)=(t_1,t_2)=(B_1r,\ \langle b_2,r\rangle+m)$. Tập thách thức là
$\Cset=\{c\in\mathcal R:\norm{c}_\infty=1,\ \norm{c}_1=\kappa\}$ với
$\kappa=60$, nên $|\Cset|=\binom{256}{60}2^{60}=2^{257.01}$; tập hiệu của nó
là $\Cdiff=\{c-c':c\neq c'\in\Cset\}$. Các mặt nạ đều (uniform) trên
$Y=[-\gamma+1,\gamma]^{kN}$ với $\gamma=2^{17}$; các đáp ứng được chấp nhận khi
và chỉ khi chúng nằm trong $Z=[-B,B]^{kN}$, $B=\gamma-\kappa-1$. Chúng tôi dùng
$\norm{fv}_\infty\le\norm{f}_1\norm{v}_\infty$ xuyên suốt. Các đại lượng cụ thể
dùng trong các cận được tập hợp trong Bảng~\ref{tab:sec-numbers}.

\begin{table}[t]
\centering
\caption{Các đại lượng cụ thể cho \texttt{bdlop-open-128-v2}}
\label{tab:sec-numbers}
\footnotesize
\begin{tabular}{@{}lr@{}}
\toprule
Quantity & Value \\
\midrule
$|\Cset|$ & $2^{257.01}$\\
Fiat--Shamir seed length & 256 bits\\
$p_{\mathrm{acc}}=((2B+1)/2\gamma)^{kN}$ (opening) & $0.1912$\\
$p_{\mathrm{acc}}^2$ (equality proof) & $0.0365$\\
Attempt cap $\Amax$ (opening; equality) & $512$; $4\,096$\\
$(1-p_{\mathrm{acc}})^{512}$; $(1-p_{\mathrm{acc}}^2)^{4096}$ & $2^{-156.7}$; $2^{-220.0}$\\
$\eta=nN\log_2(2B+1)$ (min-entropy of $w$) & $27\,646.98$ bits\\
$\beta=8\kappa B$ (M-SIS bound) & $62\,885\,280\approx2^{25.91}$\\
$2\kappa+2B$ (one honest, one relaxed opening) & $262\,142\approx2^{18.0}$\\
$\sqrt{q/2}$ (invertibility threshold) & $23\,170.5$\\
$3^{kN}q^{-(n+1)N}$ & $\approx2^{-48\,079}$\\
Statement $x$ & 274 bytes\\
Envelope, in-band / registered & 14\,885 / 8\,165 bytes\\
\bottomrule
\end{tabular}
\end{table}

\begin{lemma}[Hiệu của các thách thức là khả nghịch]
\label{lem:invertible}
Mọi $f\in\Cdiff$ đều khả nghịch trong $\Rq$.
\end{lemma}
\begin{proof}
Vì $q\equiv5\pmod8$, $X^N+1$ tách theo modulo $q$ thành đúng hai nhân tử bất
khả quy, và mọi $y\in\Rq$ khác không với $\norm{y}_\infty<q^{1/2}/\sqrt2$ đều
khả nghịch (Lyubashevsky--Seiler~\cite{ref:ls18}, Hệ quả~1.2 với $k=2$). Mọi
$f\in\Cdiff$ đều khác không với $\norm{f}_\infty\le2<23\,170$. Hàm khởi tạo của
bộ tham số bắt buộc cả hai điều kiện (\texttt{params.py:49--51},
\texttt{ring.py:52--58}).
\end{proof}

\begin{definition}[Các phần mở]
\label{def:relaxed}
Một \emph{phần mở nới lỏng} của $t=(t_1,t_2)$ là một cặp $(\bar r,f)$ với
$f\in\Cdiff$, $\norm{\bar r}_\infty\le2B$ và $ft_1=B_1\bar r$. Nó xác định thông
điệp $\bar m=t_2-f^{-1}\langle b_2,\bar r\rangle$, là $\bar m$ duy nhất thỏa
$ft_2=\langle b_2,\bar r\rangle+f\bar m$ (Bổ đề~\ref{lem:invertible}). Một
\emph{phần mở trung thực} là $(r,1)$ với $r\in S_1^k$ và $t_1=B_1r$, có thông
điệp $t_2-\langle b_2,r\rangle$; đó là thứ mà một lần tiết lộ bộc lộ ra.
\end{definition}

\paragraph{Các giả thiết}
$\Adv^{\MSIS}$ là lợi thế của việc tìm $x\in\mathcal R^k$ với
$B_1x\equiv0\pmod q$, $0<\norm{x}_\infty\le\beta=8\kappa B$, cho $B_1'$ đều.
$\Adv^{\MLWE}$ là lợi thế của việc phân biệt $(A,As+e)$ với đều, cho
$A\in\Rq^{(n+1)\times\lambda}$ đều và $s\in S_1^\lambda$, $e\in S_1^{n+1}$.
Đây là hai thể hiện được ước lượng ở 140/127 và 164/150 bit (core-SVP cổ
điển/lượng tử, Bảng~\ref{tab:params}); ước lượng M-SIS lượng tử chỉ vừa dưới
128 bit. Mọi nghiệm M-SIS trong mục này đều có chuẩn tối đa $\beta$, nên không
cần thêm giả thiết lưới nào nữa.
$\Adv^{\mathsf{cr}}_{\mathsf{SHA2}}$ và $\Adv^{\mathsf{cr}}_{\mathsf{SHA3}}$ là
các lợi thế tìm va chạm đối với SHA-256 và SHA3-256, và
$\Adv^{\mathsf{euf\text{-}cma}}_{\mathsf{ML\text{-}DSA}}$ là lợi thế EUF-CMA đối
với ML-DSA-65~\cite{ref:fips204}.

\paragraph{Các oracle ngẫu nhiên}
Chúng tôi mô hình hóa như các oracle ngẫu nhiên độc lập các thể hiện SHAKE-256
được dùng cho hạt giống thách thức Fiat--Shamir $\Hfs$ (và $H_{\mathsf{eq}}$ cho
chứng minh đẳng thức), phép mở rộng thách thức $G:\{0,1\}^{256}\to\Cset$
(\textsf{SampleInBall}), luồng mặt nạ, phép mở rộng khóa, và hàm băm phòng hộ
(hedging) SHA3-256. Các thẻ miền của chúng là không-tiền-tố đôi một
(Phụ lục~\ref{app:code-check}), nên một oracle SHAKE-256 tách thành các oracle
độc lập. $G$ đều chính xác trên $\Cset$ với một đầu vào mới: các dấu và các chỉ
số Fisher--Yates được rút bằng lấy mẫu loại bỏ chính xác. Phép mở rộng khóa
dưới dạng một oracle ngẫu nhiên làm cho $\mathsf{ck}$ đều; đó là một khóa toàn
cục duy nhất được dẫn xuất từ một hạt giống công khai (không có khóa riêng cho
từng người dùng). $Q_H$ ký hiệu tổng số truy vấn oracle ngẫu nhiên của đối thủ,
bao gồm cả những truy vấn được thực hiện thay mặt nó bởi các truy vấn kiểm
chứng, cộng thêm một.

\subsection{Cú pháp của lược đồ hợp thành}
\label{sec:sec-syntax}

Tầng vận chuyển là một tập các hàm tất định của một hồ sơ vật mang đã ghim
$\mathsf{prof}$ và một khóa đồng bộ $k_s$: $\mathsf{Place}(\vid,\ell)$ trả về
các vị trí vật mang có thứ tự $P$ cho một envelope $\ell$ byte;
$\Canon(\vid,P)$ trả về digest đã chuẩn hóa theo vật mang $D$;
$\mathsf{Write}(\vid,P,E)$ ghi envelope $E$ vào các vật mang;
$\mathsf{Read}(\vid)$ trả về $(E,P,D)$ hoặc $\bot$. \emph{Tính đúng đắn vận
chuyển} nghĩa là $\mathsf{Read}(\mathsf{Write}(\vid,P,E))=(E,P,\Canon(\vid,P))$
với $P=\mathsf{Place}(\vid,|E|)$; đó là một tính chất hàm, và việc nó thất bại
chỉ gây ra từ chối.

\begin{itemize}
\item $\mathsf{Setup}$: cố định bộ tham số, $\pid$ (một digest SHA3-256 của tất
  cả các tham số và hạt giống thiết lập) và $\mathsf{ck}$.
\item $\mathsf{RegKeyGen}$, $\mathsf{RegSign}$: bên phát hành nắm một cặp khóa
  ML-DSA-65 $(\mathit{pk}_I,\mathit{sk}_I)$ và ký một sổ đăng ký
  $\mathsf{Reg}=(e,\Delta)$: một epoch và một tập các mô tả quan hệ. Một mô tả
  băm $\pid$, khóa, module quan hệ và chính sách (chính sách codec và hồ sơ vật
  mang); trong chế độ đã đăng ký nó còn băm một cam kết chứng chỉ $t_c$.
\item $\mathsf{Pins}$: cho $\mathsf{Reg}$, $\mathit{pk}_I$, $\mathsf{prof}$,
  một ngưỡng chống quay lui $e_{\min}$ và (chế độ đã đăng ký) $t_c$, kiểm tra
  chữ ký, tính lại mô tả kỳ vọng, đòi hỏi đúng một kết quả khớp, và xuất ra
  $\mathsf{pins}=(\mathit{rid},\mathit{root},e,\mathit{cpid},
  \mathit{prof\_id})$, hoặc thất bại.
\item $\mathsf{Issue}$: rút một thách thức đều 32 byte $N_s$ và lưu nó cùng ràng
  buộc $\mathsf{bind}(\mathsf{pins})$ và một thời điểm hết hạn.
\item $\mathsf{Embed}$: trên vật phủ $\vid$, tải trọng $p$ (in-band) hoặc chứng
  chỉ $(t_c,r_c)$ (đã đăng ký), $\mathsf{pins}$ và $N_s$: tính
  $P=\mathsf{Place}(\vid,\ell)$ từ một placeholder bằng không có độ dài cuối
  cùng $\ell$, $D=\Canon(\vid,P)$ và $x=\Stmt(\mathsf{pins},N_s,P,D)$; với
  in-band, lấy mẫu $r\leftarrow S_1^k$, đặt $t=\Com(\mathsf{enc}(p);r)$; chứng
  minh $\pi\leftarrow\mathsf{P}(x,t,r)$; lập $E=\texttt{"ZKLB"}\|1\|\pid\|N_s\|t\|\pi$
  (hoặc $\texttt{"ZKLR"}\|1\|\pid\|N_s\|\pi$) và xuất ra
  $\mathsf{Write}(\vid,P,E)$ và, với in-band, phần mở $(p,r)$.
\item $\mathsf{Verify}$: trên $\vid^*$: $(E^*,P^*,D^*)=\mathsf{Read}(\vid^*)$;
  phân tích $E^*$ một cách nghiêm ngặt; lấy $N_s^*$ từ $E^*$ (chế độ có lưu
  trữ) hoặc dùng một $N_s$ cố định (chế độ không trạng thái);
  $x^*=\Stmt(\mathsf{pins},N_s^*,P^*,D^*)$; chấp nhận khi và chỉ khi
  $\mathsf{V}(x^*,t^*,\pi^*)=1$ với $t^*$ từ $E^*$ (in-band) hoặc $t_c$ đã ghim
  (đã đăng ký) và, trong chế độ có lưu trữ, $N_s^*$ còn sống đối với
  $\mathsf{bind}(\mathsf{pins})$ và được tiêu thụ một cách nguyên tử.
\item $\mathsf{VerDisc}$: chấp nhận $(E,p,r)$ khi và chỉ khi $(r,1)$ là một phần
  mở trung thực của cam kết trong $E$ tới $\mathsf{enc}(p)$.
\item $\mathsf{Link}$, $\mathsf{VerLink}$: chứng minh đẳng thức
  $\Pi_{\mathsf{eq}}$ trên hai cam kết, được ràng buộc vào một digest ngữ cảnh
  $\mathit{lx}=\mathrm{SHA3}(\text{tag}\|\pid\|\text{ngữ cảnh có tiền tố độ dài})$.
\end{itemize}

Phát biểu là chuỗi 274 byte
\begin{multline*}
\Stmt=\mathtt{u16}(1)\|\mathit{rid}\|\pid\|\mathit{root}\|\mathtt{u64}(e)\|N_s\|
\mathit{cpid}\|\mathit{prof\_id}\\
\|\mathtt{u64}(|P|)\|\mathrm{SHA256}(P)\|V,\quad
V=\mathrm{SHA256}(\text{tag}\|D),
\end{multline*}
với $\mathrm{SHA256}(P)$ được lấy trên một mã hóa JSON chuẩn tắc của các vị trí
có thứ tự. Chứng minh mở $\Pi_{\mathsf{open}}$ xuất ra $\pi=(h,z)$ với
$h=\Hfs(\text{tag}\|\pid\|\mathtt{u64}(|x|)\|x\|\mathsf{enc}(t)\|\mathsf{enc}(w))$,
$w=B_1y$, $c=G(h)$, $z=y+cr\in Z$; bên kiểm chứng tính lại $w'=B_1z-ct_1$ và
kiểm tra $h$. Mọi trường sau $x$ đều có độ dài được cố định bởi $\pid$, nên đầu
vào của hàm băm giải mã được một cách duy nhất.

\subsection{Mô hình đối thủ}
\label{sec:sec-adversary}

Đối thủ $\mathcal A$ là một thuật toán xác suất cổ điển. Nó kiểm soát mọi kênh,
chọn các vật phủ và các tải trọng, biết $k_s$ (không tính chất nào dưới đây phụ
thuộc vào tính bí mật của nó), và có các oracle sau. $\mathcal O_{\mathsf{Iss}}$
trả về một $N_s$ mới. $\mathcal O_{\mathsf{Emb}}(\vid,p,N_s)$ chạy một
$\mathsf{Embed}$ in-band trung thực và trả về video stego;
$\mathcal O_{\mathsf{EmbR}}(j,\vid,N_s)$ làm điều tương tự cho chứng chỉ đã đăng
ký của bên nắm giữ trung thực $j$. Chúng tôi cho phép các oracle này chứng minh
các phát biểu $x$ tùy ý do $\mathcal A$ chọn, điều này chỉ làm mạnh thêm các kết
quả. Các đầu ra của chúng tạo thành danh sách $L$ các bộ ba $(x,t,\pi)$.
$\mathcal O_{\mathsf{Ver}}$ chạy $\mathsf{Verify}$ với kho dùng chung.
$\mathcal O_{\mathsf{Link}}$ trả về các chứng minh liên kết trung thực giữa các
cam kết đã ghi. $\mathcal O_{\mathsf{Disc}}(i)$ trả về phần mở của cam kết trung
thực thứ $i$ và đánh dấu nó là \emph{đã tiết lộ}.
$\mathcal O_{\mathsf{Reg}}(e,\Delta,T)$ khiến bên phát hành ký một sổ đăng ký
với epoch $e$ mà các mô tả của nó là các mô tả in-band $\Delta$ và các mô tả đã
đăng ký mà bên phát hành tính từ các cam kết chứng chỉ trong $T$, tất cả do
$\mathcal A$ chọn; $T_{\mathsf{reg}}$ là hợp của mọi $T$ đã truy vấn. (Một bên
phát hành chỉ ký một mô tả đã đăng ký cho một chứng chỉ mà nó đã thẩm định, nên
các mô tả đã đăng ký chỉ đi vào qua $T$.) $Q_P$ chặn số chứng minh được tạo bởi
các bên trung thực.

Bên phát hành và bên kiểm chứng là trung thực; kho phiên là tuyến tính hóa được
(linearizable) và bền vững. Việc làm hỏng bên nắm giữ là tĩnh (đối thủ đăng ký
các chứng chỉ của riêng nó qua $\mathcal O_{\mathsf{Reg}}$); tiết lộ là hình
thức làm hỏng thích nghi duy nhất, và các cam kết đã tiết lộ bị loại khỏi các
điều kiện thắng. Các kênh phụ, lỗi codec, và tính đúng đắn của tầng vận chuyển
nằm ngoài phạm vi.

\subsection{Tính không tiết lộ tri thức của chứng minh Fiat--Shamir-with-aborts}
\label{sec:sec-zk}

\begin{lemma}[Chấp nhận độc lập với nhân chứng]
\label{lem:hvzk}
Với mọi $r\in S_1^k$ và $c\in\Cset$, một lần thử với $y\leftarrow Y$ được chấp
nhận với xác suất đúng bằng $p_{\mathrm{acc}}$, và với điều kiện được chấp nhận
thì $z=y+cr$ đều trên $Z$.
\end{lemma}
\noindent
Mỗi hệ số $v$ của $cr$ thỏa $|v|\le\kappa$, nên $2B+1$ giá trị
$z_i\in[-B,B]$ tương ứng với các mặt nạ $z_i-v\in[-\gamma+1,\gamma-1]\subset
[-\gamma+1,\gamma]$, mỗi cái có xác suất $1/2\gamma$.

\begin{lemma}[Mô phỏng]
\label{lem:sim}
Tồn tại một bộ mô phỏng mà, chỉ với $(x,t)$ và quyền lập trình $\Hfs$, xuất ra
các chứng minh có phân phối hợp với các oracle ngẫu nhiên nằm trong khoảng
\begin{multline*}
\epszk(Q_P,Q_H)=2Q_P(Q_H+Q_P)2^{-256}\\
+(\Amax+1)Q_P\bigl(Q_H+(\Amax+1)Q_P\bigr)2^{-\eta}
\end{multline*}
so với $Q_P$ chứng minh thực, với mọi đối thủ thực hiện $Q_H$ truy vấn. Điều
tương tự đúng cho $\Pi_{\mathsf{eq}}$ với $p_{\mathrm{acc}}^2$ thay cho $p_{\mathrm{acc}}$.
\end{lemma}
\noindent
Bộ mô phỏng lấy mẫu số lần thử từ phân phối hình học bị cắt cụt với tham số
$p_{\mathrm{acc}}$, xuất ra $\bot$ với xác suất $(1-p_{\mathrm{acc}})^{\Amax}$,
và ngược lại lấy mẫu $h\leftarrow\{0,1\}^{256}$, $z\leftarrow Z$, đặt
$w=B_1z-G(h)t_1$ và lập trình $\Hfs$ tại $(\pid,x,t,w)$. Bước tế nhị là các truy
vấn oracle ngẫu nhiên mà bên chứng minh thực thực hiện trong các lần thử bị loại
và không bao giờ tiết lộ: các phân tích trước đây về Fiat--Shamir with aborts đã
xử lý sai việc tái lập trình tương ứng trong bước CMA-to-NMA, như được Barbosa
và cộng sự~\cite{ref:barbosa} chỉ ra và sửa chữa, và được Devevey và cộng
sự~\cite{ref:devevey} phân tích chi tiết. Ở đây bước này được giải quyết trong
ROM cổ điển nhờ hai sự kiện đặc thù của hiện thực này: sự chấp nhận là
\emph{hoàn toàn} độc lập với nhân chứng và với $c$ (Bổ đề~\ref{lem:hvzk}), và
mọi điểm truy vấn ẩn đều chứa $w=B_1y$ có min-entropy ít nhất $\eta$ (bởi vì
$B_1=[I_n|B_1']$). Tính độc lập chính xác cũng có nghĩa là số lần thử không rò
rỉ gì, dù rằng ngoài ra phần hiện thực không chạy trong thời gian hằng.

\subsection{Tính đúng đắn tri thức và ràng buộc vào video}
\label{sec:sec-soundness}

\begin{definition}[Ràng buộc video trích xuất được qua mô phỏng]
\label{def:se}
Trong thí nghiệm $\mathsf{SE}$, $\mathcal A$ tương tác với mọi oracle và xuất
ra $\vid^*$. Gọi $x^*$ là phát biểu mà bên kiểm chứng tính lại từ $\vid^*$ và
các mục ghim của chính nó, và $(t^*,\pi^*)$ là cam kết và chứng minh mà nó kiểm
chứng. $\mathcal A$ \emph{thắng} nếu $\mathsf{Verify}$ chấp nhận và
$(x^*,t^*,\pi^*)\notin L$. Một bộ trích xuất $\mathcal E^{\mathcal A}$
\emph{thành công} nếu nó xuất ra $x^*$ và một phần mở nới lỏng của $t^*$, trong
đó $x^*$ là phát biểu của bên kiểm chứng trong một lần chạy thắng của $\mathcal A$.
\end{definition}

\begin{theorem}[Tính đúng đắn tri thức ràng buộc vào phát biểu]
\label{thm:se}
Giả sử $\mathcal A$ thắng $\mathsf{SE}$ với xác suất $\varepsilon$. Tồn tại một
bộ trích xuất $\mathcal E$ chạy $\mathcal A$ hai lần với các oracle được mô
phỏng và thành công với xác suất $\varepsilon_{\mathsf{ext}}$ sao cho
\begin{multline*}
\varepsilon\le\epszk(Q_P,Q_H)+\Adv^{\MSIS}+Q_H2^{-256}\\
+\sqrt{Q_H\bigl(\varepsilon_{\mathsf{ext}}+Q_H2^{-256}+2^{-257}\bigr)}.
\end{multline*}
\end{theorem}
\noindent
\emph{Ý tưởng chứng minh.} Thay các chứng minh trung thực bằng các chứng minh
được mô phỏng (Bổ đề~\ref{lem:sim}). Nếu $\mathcal A$ tái dùng một điểm đã lập
trình, tức $(x^*,t^*)=(x_i,t_i)$ và $w^*=w_i$, thì $h^*=h_i$, do đó $c^*=c_i$ và
$B_1(z^*-z_i)=0$ với $0<\norm{z^*-z_i}_\infty\le2B$: một nghiệm M-SIS. Ngược
lại, truy vấn $(\pid,x^*,t^*,w^*)$ của bên kiểm chứng được trả lời bằng một giá
trị ngẫu nhiên mới, và bổ đề rẽ nhánh tổng quát của Bellare và
Neven~\cite{ref:bn06} (Bổ đề~1) với không gian đầu ra $2^{256}$ cho hai chứng
minh được chấp nhận có cùng $(x^*,t^*,w^*)$ và các hạt giống $h\neq h'$. Khi đó
$c=G(h)\neq G(h')=c'$ ngoại trừ với xác suất $2^{-257}+Q_H2^{-256}$, và
$B_1(z-z')=(c-c')t_1^*$ cho phần mở nới lỏng $(z-z',c-c')$. Hiệu thách thức là
khả nghịch theo Bổ đề~\ref{lem:invertible}, chính điều này khiến phần mở xác
định một thông điệp duy nhất. Bởi vì truy vấn bị rẽ nhánh chứa $x^*$, phần mở
được ràng buộc vào đúng phát biểu do bên kiểm chứng tính lại: cam kết video, các
vị trí vật mang, thách thức phiên, gốc sổ đăng ký và epoch, quan hệ và chính
sách của nó. \qed

Tổn thất căn bậc hai là cái giá thường thấy của việc tua lại (rewinding). Với
$Q_H=2^{64}$, cận biến một xác suất trích xuất $2^{-128}$ thành $\varepsilon
\lesssim2^{-32}$. Đối với tính xác thực (Định lý~\ref{thm:auth}), có một cận
chặt nhờ hướng SelfTargetMSIS của Kiltz, Lyubashevsky và
Schaffner~\cite{ref:kls18}, với cái giá là giả thiết phi chuẩn đó.

\begin{theorem}[Tính ràng buộc đối với các phần mở nới lỏng]
\label{thm:binding}
Không đối thủ nào có thể xuất ra $t$ và hai phần mở của $t$, mỗi phần mở trung
thực hoặc nới lỏng, với các thông điệp khác nhau, ngoại trừ với xác suất
$\Adv^{\MSIS}$.
\end{theorem}
\noindent
Cho $(\bar r,f,m)$ và $(\bar r',f',m')$, vector
$x=f'\bar r-f\bar r'$ thỏa $B_1x=0$; nếu $x=0$ thì
$\langle b_2,x\rangle=ff'(m-m')=0$, mâu thuẫn với tính khả nghịch. Chuẩn của nó
tối đa là $2\cdot2\kappa\cdot2B=\beta$ cho hai phần mở nới lỏng, $2\kappa+2B$
cho một trung thực và một nới lỏng, và $2$ cho hai trung thực, tất cả đều dưới
$q/2$ (\texttt{params.py:53}), nên $x\not\equiv0$ là một nghiệm M-SIS.

\noindent
Gộp lại, các Định lý~\ref{thm:se} và~\ref{thm:binding} nói rằng: ai tạo ra một
video in-band được chấp nhận đều bị ràng buộc vào thông điệp $\bar m$ mà bộ
trích xuất tính ra, và bất kỳ lần tiết lộ về sau nào mà $\mathsf{VerDisc}$ chấp
nhận đều bộc lộ một tải trọng $p$ với $\mathsf{enc}(p)=\bar m$, trừ khi M-SIS
bị phá.

\emph{Việc tiết lộ giữ cho tính ngẫu nhiên riêng tư.} Một lần tiết lộ là tải
trọng $p$ cùng với một chứng minh đẳng thức giữa $t$ và cam kết tất định
$\mathsf{Com}(\mathsf{enc}(p);0)$, được Fiat--Shamir ràng buộc vào một ngữ cảnh
do bên nhận chọn (định danh phát biểu và nonce). Bên nhận biết $p$ nhưng không
biết phần mở nào của $t$, nên nó không thể tạo các chứng minh mở mới cho $t$;
theo tính đúng đắn của chứng minh đẳng thức, một lần tiết lộ được chấp nhận vẫn
thỏa $\mathsf{enc}(p)=\bar m$ trừ khi M-SIS bị phá. Một phiên bản trước đây bộc
lộ chính $r$; đây là mục~D3 trong Phụ lục~\ref{app:code-check} và đã được sửa.
Các định lý dưới đây vẫn coi các cam kết mà \emph{phần mở} của chúng đã bị rò rỉ
(ví dụ do lộ khóa) là đã tiết lộ.

\begin{lemma}[Tính một chiều của các cam kết trung thực]
\label{lem:row}
Cho $\mathsf{ck}$ và $t^{(j)}=\Com(m_j;r_j)$, $j\le Q$, với
$r_j\leftarrow S_1^k$ và các thông điệp $m_j$ được chọn thích nghi trước khi mỗi
$t^{(j)}$ được rút, không đối thủ nào xuất ra một phần mở nới lỏng của bất kỳ
$t^{(j)}$ nào ngoại trừ với xác suất
$\Adv^{\mathsf{row}}_Q=\Adv^{\MSIS}+Q\bigl(\Adv^{\MLWE}+3^{kN}q^{-(n+1)N}\bigr)$.
\end{lemma}
\noindent
Đặt $x=fr_j-\bar r$, nên $B_1x=0$ và $\norm{x}_\infty\le2\kappa+2B$. Nếu
$x\neq0$ thì nó là một nghiệm M-SIS (phép rút gọn đã sinh mọi $t^{(j)}$ và biết
$r_j$). Nếu $x=0$ thì đối thủ đã tìm được $r_j=f^{-1}\bar r$, điều mà một bộ
phân biệt cấy mẫu M-LWE của nó tại một $j$ được đoán sẽ kiểm tra đối chiếu với
$t^{(j)}$; với $t^{(j)}$ đều, một phần mở bậc ba tồn tại với xác suất tối đa
$3^{kN}q^{-(n+1)N}$. Chỉ riêng tính ràng buộc là không đủ ở đây: một đối thủ
khôi phục được $r_j$ sẽ mở $t^{(j)}$ tới đúng thông điệp.

\subsection{Che giấu tải trọng}
\label{sec:sec-hiding}

\begin{definition}[Tính không phân biệt được tải trọng]
\label{def:hiding}
Trong $\mathsf{IND}^b$, $\mathcal A$ có mọi oracle và nộp một lần
$(\vid,p_0,p_1,N_s)$ với $|p_0|,|p_1|\le766$ byte; bên thách thức trả về
$\mathcal O_{\mathsf{Emb}}(\vid,p_b,N_s)$. $\mathcal A$ không được truy vấn
$\mathcal O_{\mathsf{Disc}}$ hoặc $\mathcal O_{\mathsf{Link}}$ trên cam kết
thách thức. $\Adv^{\mathsf{ind}}=|\Pr[\mathsf{IND}^1=1]-\Pr[\mathsf{IND}^0=1]|$.
\end{definition}

\begin{theorem}[Che giấu]
\label{thm:hiding}
$\Adv^{\mathsf{ind}}\le2\,\epszk(Q_P+1,Q_H)+2\Adv^{\MLWE}$.
\end{theorem}
\noindent
Video stego chỉ phụ thuộc vào $b$ thông qua $E$: $P$ và $D$ được cố định từ một
placeholder trước khi chứng minh, và envelope có kích thước cố định, nên độ dài
tải trọng cũng được che giấu. Mô phỏng chứng minh thách thức
(Bổ đề~\ref{lem:sim}), rồi thay $t$ bằng một cặp đều: với
$B_1'=[B_{1,2}'\,|\,B_{1,3}']$ và $r=(r_1,r_2,r_3)$,
$(t_1-B_{1,2}'(t_2-m),\,t_2-m)=A r_3+(r_1,r_2)$ với
$A=\bigl[\begin{smallmatrix}B_{1,3}'-B_{1,2}'b_2'\\ b_2'\end{smallmatrix}\bigr]$
đều, một mẫu M-LWE hạng $\lambda$ với $n+1$ mẫu vành; ánh xạ là một song ánh với
mỗi $m$ cố định. Khi đó thí nghiệm độc lập với $b$. Tính che giấu mang tính tính
toán và là tính chất duy nhất dựa trên M-LWE; nó đúng cho tải trọng, chứ không
cho sự kiện rằng một tải trọng tồn tại, vốn là một tính chất giấu tin được đánh
giá bằng thực nghiệm.

\subsection{Khả năng chống phát lại}
\label{sec:sec-replay}

\begin{definition}[Phát lại]
\label{def:replay}
Trong $\mathsf{RPL}$, $\mathcal A$ có mọi oracle đối với một bên kiểm chứng có
kho phiên và thắng nếu (i) hai lời gọi $\mathcal O_{\mathsf{Ver}}$ chấp nhận với
cùng $N_s$; hoặc (ii) một lời gọi chấp nhận với một $N_s$ không được phát hành
cho $\mathsf{bind}(\mathsf{pins})$, hoặc đã hết hạn; hoặc (iii) một lời gọi chấp
nhận $\vid^*$ có bộ ba $(x^*,t^*,\pi^*)\notin L$ và $\mathcal A$ đã truy vấn
$\Hfs$ tại điểm $(\pid,x^*,t^*,w^*)$ của bên kiểm chứng trước khi $N_s^*$ được
phát hành.
\end{definition}

\begin{theorem}[Khả năng chống phát lại]
\label{thm:replay}
$\Pr[\text{(i) hoặc (ii)}]\le\Adv^{\mathsf{cr}}_{\mathsf{SHA2}}+Q_I^22^{-257}$ và
$\Pr[\text{(iii)}]\le Q_HQ_I2^{-256}$, với $Q_I$ phiên được phát hành.
\end{theorem}
\noindent
(i) và (ii) suy ra từ ngữ nghĩa của kho: các hàng chỉ tồn tại cho các thách
thức đã phát hành (khóa theo một digest SHA-256 của $N_s$), và việc tiêu thụ là
một \texttt{UPDATE \dots\ WHERE consumed\_at IS NULL AND expires\_at>now} có
điều kiện đơn nhất bên trong một giao dịch ghi, nên nhiều nhất một bên gọi quan
sát được một hàng đã cập nhật (\texttt{zkp\_sessions.py:181--241}). Với (iii),
$N_s^*$ là một chuỗi con của $x^*$, nên một truy vấn thực hiện trước đó dự đoán
một giá trị 256 bit đều. Cùng với Định lý~\ref{thm:se}, một phiên được chấp nhận
tương ứng với một bên đã biết một phần mở và đã tính chứng minh sau khi phiên
được phát hành.

\emph{Điều mà một bên kiểm chứng không trạng thái không thể đem lại.} Với một
$N_s$ cố định (\texttt{expected\_session}), (i) thất bại: bất kỳ bản sao nào của
một video được chấp nhận đều lại được chấp nhận, bởi bất kỳ ai và bao nhiêu lần
tùy ý, vì một chứng minh không tương tác có thể chuyển nhượng. Nếu $N_s$ bị tái
dùng qua nhiều lần kiểm chứng, (iii) cũng thất bại và chỉ còn lại ràng buộc
video (Định lý~\ref{thm:se}). Không biến thể nào ngăn được việc tiếp chuyển
(relay): chế độ in-band không ràng buộc $N_s$ vào một bên chứng minh, nên ai
trình ra một video hợp lệ trước thì tiêu thụ phiên.

\subsection{Tính xác thực chế độ đã đăng ký}
\label{sec:sec-auth}

\begin{definition}[Tính xác thực]
\label{def:auth}
Trong $\mathsf{AUTH}$, bên thách thức tạo các chứng chỉ trung thực
$t_c^{(j)}=\Com(m_j;r_j)$, $j\le Q_C$, với $m_j$ do $\mathcal A$ chọn, và
$\mathcal A$ có mọi oracle. $\mathcal A$ xuất ra $(\mathsf{Reg}^*,\vid^*,t^*)$
và thắng nếu một bên kiểm chứng ghim $\mathit{pk}_I$ và $t^*$ chấp nhận $\vid^*$
với các mục ghim từ $\mathsf{Reg}^*$, bộ ba $(x^*,t^*,\pi^*)\notin L$, và hoặc
(a) $t^*\notin T_{\mathsf{reg}}$ (chưa bao giờ được bên phát hành đăng ký), hoặc
(b) $t^*=t_c^{(j)}$ cho một $j$ trung thực chưa tiết lộ.
\end{definition}

\begin{theorem}[Tính xác thực]
\label{thm:auth}
\begin{multline*}
\Adv^{\mathsf{auth}}\le\Adv^{\mathsf{euf\text{-}cma}}_{\mathsf{ML\text{-}DSA}}
+\Adv^{\mathsf{cr}}_{\mathsf{SHA2}}+\epszk(Q_P,Q_H)\\
+\Adv^{\MSIS}+Q_H2^{-256}\\
+\sqrt{Q_H\bigl(\Adv^{\mathsf{row}}_{Q_C}+Q_H2^{-256}+2^{-257}\bigr)}.
\end{multline*}
\end{theorem}
\noindent
Trường hợp (a) cho một chữ ký giả mạo ML-DSA trên các byte ký của sổ đăng ký
hoặc, nếu $\mathsf{Reg}^*$ đã được ký, một va chạm SHA-256 giữa các đầu vào có
độ dài cố định $B_1'\|b_2'\|t^*$ và $B_1'\|b_2'\|t'$ với một
$t'\in T_{\mathsf{reg}}$ nào đó, vì mô tả khớp chứa digest đó. Trong trường hợp
(b), Định lý~\ref{thm:se} áp dụng cho trò chơi được mô phỏng trích xuất một phần
mở nới lỏng của một chứng chỉ trung thực, điều mà Bổ đề~\ref{lem:row} loại trừ.
Định lý này nói về các chứng chỉ; sự liên kết tới một con người chỉ mạnh ngang
với việc bên phát hành kiểm tra bên nắm giữ trước khi ký. Khi bên phát hành ràng
buộc một định danh của bên nắm giữ, mô tả đã ký chứa một digest có tách miền của
nó bên cạnh cam kết, nên quan hệ đã ghim của bên kiểm chứng nêu tên cả hai
(Phụ lục~\ref{app:code-check}, mục~D2, đã sửa).

\subsection{Liên kết}
\label{sec:sec-linkage}

\begin{definition}[Tính đúng đắn và tính che giấu của liên kết]
\label{def:link}
Trong $\mathsf{LNK\text{-}S}$, $\mathcal A$ xuất ra một ngữ cảnh, $t_a,t_b$, một
$\pi_{\mathsf{eq}}$ được chấp nhận không do $\mathcal O_{\mathsf{Link}}$ xuất ra
cho bộ ba đó, và các phần mở (trung thực hoặc nới lỏng) của $t_a$ và $t_b$ với
các thông điệp khác nhau. Trong $\mathsf{LNK\text{-}H}^b$, $\mathcal A$ nộp
$(p_0,p_1,\mathit{ctx})$ và nhận $t_a=\Com(p_b;r_a)$,
$t_b=\Com(p_b;r_b)$ và $\pi_{\mathsf{eq}}$ cho $\mathit{ctx}$.
\end{definition}

\begin{theorem}[Liên kết]
\label{thm:link}
$\Adv^{\mathsf{lnk\text{-}s}}$ tuân theo cận của Định lý~\ref{thm:se} với
$\varepsilon_{\mathsf{ext}}\le\Adv^{\MSIS}$, và
$\Adv^{\mathsf{lnk\text{-}h}}\le2\,\epszk(Q_P+1,Q_H)+4\Adv^{\MLWE}$.
\end{theorem}
\noindent
Hai bản ghi bị rẽ nhánh với một thông báo $(w_a,w_b,u)$ cho các phần mở nới lỏng
$(\bar r_a,f)$, $(\bar r_b,f)$ với một $f$ \emph{chung} và
$f(t_{2,a}-t_{2,b})=\langle b_2,\bar r_a-\bar r_b\rangle$, do đó
$\bar m_a=\bar m_b$; hai phần mở của đối thủ khi đó mâu thuẫn với
Định lý~\ref{thm:binding} cho $t_a$ hoặc cho $t_b$. Việc liên kết tới một chứng
chỉ đã đăng ký còn xác thực thêm bên nắm giữ của nó, như trong
Định lý~\ref{thm:auth}, bởi vì việc trích xuất cũng cho ra một phần mở nới lỏng
của $t_c$. Tính che giấu dùng bộ mô phỏng cho $\Pi_{\mathsf{eq}}$ và hai bước
lai M-LWE; sự bằng nhau của hai tải trọng được bộc lộ theo thiết kế.

\subsection{Tính toàn vẹn của nội dung video}
\label{sec:sec-integrity}

$\Canon$ chia luồng bằng bộ phân đoạn đã ghim (\texttt{ffmpeg -c:v copy
-f segment}, $F$ khung mỗi đoạn) và kết hợp các digest của từng đoạn với
$F$ và các số đếm khung. Một đoạn mang các chunk envelope được băm qua một phép
tuần tự hóa có gắn thẻ, phân định bằng độ dài: với mỗi đơn vị NAL là chỉ số,
kiểu, \texttt{forbidden\_zero\_bit}, \texttt{nal\_ref\_idc} và kích thước mã
khởi đầu (start-code); RBSP nguyên văn cho các đơn vị NAL không có vật mang; và,
bên trong một đơn vị NAL IDR có vật mang, các bit RBSP giữa các khối vật mang
cho tới bit dừng (stop bit), với mỗi khối vật mang được thay bằng vector hệ số
của nó trong đó mỗi vật mang được chuẩn hóa (độ lớn làm tròn xuống số chẵn, giữ
dấu). Các đoạn khác được băm từng byte.

\begin{lemma}[Digest và envelope xác định những gì]
\label{lem:canon}
Phép tuần tự hóa giải mã được một cách duy nhất, nên $V$ bằng nhau kéo theo
$\Canon$ bằng nhau ngoại trừ với xác suất $3\Adv^{\mathsf{cr}}_{\mathsf{SHA2}}$.
Dưới quy tắc vật mang ổn định ($|\ell|\ge4$, một bit chẵn-lẻ), một mức vật mang
là $\mathrm{sgn}(\ell)(2\lfloor|\ell|/2\rfloor+b)$, nên $(\Canon,P)$ và các bit
vật mang xác định toàn bộ vector hệ số của mọi khối vật mang. Nếu chuỗi bit vật
mang là mã hóa duy nhất của $E$ dưới hồ sơ đã ghim, thì $(\Canon,P,E)$ xác định:
sự phân đoạn và các số đếm khung; mọi byte của các đoạn không có vật mang; mọi
trường header NAL và bit RBSP bên ngoài các khối vật mang; và các vector hệ số
của các khối vật mang.
\end{lemma}

\begin{theorem}[Tính toàn vẹn]
\label{thm:integrity}
Giả sử một lần nhúng trung thực tạo ra $\vid_h$ cho phiên $N_s$ với cam kết $t$
(hoặc chứng chỉ $t_c$), và $\mathcal A$ xuất ra $\vid^*$ được chấp nhận cho
$N_s$ với cùng cam kết (hoặc chứng chỉ), không tiết lộ. Khi đó
$\Canon(\vid^*)=\Canon(\vid_h)$, $P^*=P_h$ và $E^*=E_h$, ngoại trừ với xác suất
tối đa
\begin{multline*}
4\Adv^{\mathsf{cr}}_{\mathsf{SHA2}}+\epszk(Q_P,Q_H)+\Adv^{\MSIS}+Q_H2^{-256}\\
+\sqrt{Q_H\bigl(\Adv^{\mathsf{row}}_{Q_P}+Q_H2^{-256}+2^{-257}\bigr)}.
\end{multline*}
\end{theorem}
\noindent
Nếu $x^*\neq x_h$, bộ ba là mới và việc trích xuất cho ra một phần mở nới lỏng
của một cam kết trung thực. Nếu $x^*=x_h$, giá trị băm vị trí và $V$ khớp nhau,
và Bổ đề~\ref{lem:canon} được áp dụng; khi đó chứng minh chính là chứng minh
trung thực, ngược lại trường hợp điểm-đã-lập-trình của Định lý~\ref{thm:se} cho
một nghiệm M-SIS.

\paragraph{Các giới hạn của khẳng định toàn vẹn}
\begin{itemize}
\item Chế độ in-band ràng buộc nội dung vào một \emph{cam kết}, chứ không vào
  một con người: bất kỳ ai nắm một $N_s$ còn sống đều có thể nhúng lại bất kỳ
  video đã chỉnh sửa nào với một cam kết mới của riêng nó, và nó sẽ được chấp
  nhận.
\item Các vật mang chỉ được bảo vệ như các bit envelope. Định lý~\ref{thm:integrity}
  cần các bit vật mang là mã hóa duy nhất của $E$. Bên kiểm chứng chia lại chunk
  của $E$ theo kích thước chunk đã ghim và từ chối trừ khi các mảnh trích xuất
  được đúng là mã hóa đó; nếu không có phép kiểm này, một envelope một-chunk sẽ
  để trống trường \texttt{chunk\_size} 16 bit của nó (mục~D1 trong
  Phụ lục~\ref{app:code-check}, đã sửa).
\item Không được ràng buộc: các byte không sống sót qua bộ phân đoạn đã ghim (dữ
  liệu bộ chứa, các luồng khác, dữ liệu ngoài các đơn vị NAL); các bit căn lề
  RBSP sau bit dừng trong các đơn vị NAL IDR có mang vật mang; các từ mã CAVLC
  thay thế cho cùng một vector hệ số của một khối vật mang; tầng chống giả lập
  (digest bao phủ RBSP); và phiên bản của bộ phân đoạn, vốn được chính sách
  codec nêu tên nhưng không ghim.
\item Không có khẳng định về tính đúng đắn codec: digest bao phủ cú pháp, không
  bao phủ các điểm ảnh đã giải mã. Không có khẳng định về nguồn gốc: không gì
  nói rằng một máy quay đã ghi $\vid_h$ hay rằng tải trọng là đúng sự thật.
\end{itemize}

\subsection{Các đối thủ lượng tử}
\label{sec:sec-qrom}

Mọi định lý ở trên đều là các phát biểu trong ROM cổ điển. Việc tua lại
(Định lý~\ref{thm:se}) và việc lập trình thích nghi (Bổ đề~\ref{lem:sim}) không
chuyển sang QROM như đã phát biểu. Barbosa và cộng sự~\cite{ref:barbosa} đưa ra
các chứng minh QROM của bước CMA-to-NMA cho Fiat--Shamir with aborts, và Don và
cộng sự~\cite{ref:dfms19} cùng Liu và Zhandry~\cite{ref:lz19} chỉ ra rằng phép
biến đổi Fiat--Shamir bảo toàn tính đúng đắn (tri thức) của các $\Sigma$-protocol
phù hợp trước các đối thủ lượng tử với tổn thất đa thức. Chúng tôi chưa kiểm
chứng rằng các giả thiết của họ đúng cho bộ trích xuất nới lỏng và bên chứng
minh có loại bỏ được dùng ở đây, nên chúng tôi không khẳng định định lý QROM
nào. Cách diễn đạt của chúng tôi là ``hậu lượng tử một cách hợp lý'': các phép
rút gọn nhắm tới M-SIS, M-LWE, ML-DSA và va chạm băm, không cái nào có một tấn
công lượng tử hiệu quả đã biết ở các tham số này.
